\documentclass[12pt, a4paper]{report}

% Pass hyperref options safely
\PassOptionsToPackage{breaklinks}{hyperref}

\usepackage{graphicx} % For including images
\usepackage{titlesec} % For customizing section titles
\usepackage{tocloft} % For customizing table of contents
\usepackage{acro} % For acronyms

%% ____Bibliography____%%
\usepackage[numbers,sort&compress]{natbib}
\usepackage{amsmath}
\usepackage{amssymb} % for \mathbb (used by Block 8 DKG formulation)
% Page margins
\usepackage[left=1.5in, right=1in, top=1in, bottom=1in]{geometry}
% Hyperlinks (LOAD ONCE, LOAD LAST)
\usepackage{hyperref}
\usepackage{float}
\usepackage{placeins}
\usepackage{xcolor}     % for \rowcolor
 \usepackage{colortbl}   % for row shading
 \usepackage{makecell}   % for \makecell and \thead
\definecolor{gapgray}{gray}{0.92}

\usepackage{comment}
\usepackage{subcaption}
\usepackage{algorithm}
% \usepackage{algorithmic}
\usepackage{algpseudocode}
\usepackage{soul}    % in preamble
\usepackage{longtable}


%\hypersetup{colorlinks=true,citecolor=blue,linkcolor=blue,urlcolor=blue}


% Remove page number from the first page
\thispagestyle{empty}

% Customize table of contents, list of figures, and list of tables
\renewcommand{\cfttoctitlefont}{\hfill\Large\bfseries}
\renewcommand{\cftaftertoctitle}{\hfill}
\renewcommand{\cftloftitlefont}{\hfill\Large\bfseries}
\renewcommand{\cftafterloftitle}{\hfill}
\renewcommand{\cftlottitlefont}{\hfill\Large\bfseries}
\renewcommand{\cftafterlottitle}{\hfill}

% This is my customize variable
\newcommand{\mypar}{\vspace{0.8em}\noindent} 

% Define acronyms
% Define acronyms
\DeclareAcronym{ITS}{
	short = ITS,
	long = Intelligent Transportation Systems
}
\DeclareAcronym{VANET}{
	short = VANET,
	long = Vehicular Ad Hoc Network
}
\DeclareAcronym{V2V}{
	short = V2V,
	long = Vehicle-to-Vehicle
}
\DeclareAcronym{V2I}{
	short = V2I,
	long = Vehicle-to-Infrastructure
}
\DeclareAcronym{SDVN}{
	short = SDVN,
	long = Software-Defined Vehicular Network
}
\DeclareAcronym{SDN}{
	short = SDN,
	long = Software-Defined Networking
}
\DeclareAcronym{AI}{
	short = AI,
	long = Artificial Intelligence
}
\DeclareAcronym{ML}{
	short = ML,
	long = Machine Learning
}
\DeclareAcronym{LLM}{
	short = LLM,
	long = Large Language Model
}
\DeclareAcronym{FL}{
	short = FL,
	long = Federated Learning
}
\DeclareAcronym{RSU}{
	short = RSU,
	long = Roadside Unit
}
\DeclareAcronym{IDS}{
	short = IDS,
	long = Intrusion Detection System
}
\DeclareAcronym{OSR}{
	short = OSR,
	long = Open-Set Recognition
}
\DeclareAcronym{DAE}{
	short = DAE,
	long = Deep Auto-Encoder
}
\DeclareAcronym{LSTM}{
	short = LSTM,
	long = Long Short-Term Memory
}
\DeclareAcronym{KNN}{
	short = KNN,
	long = K-Nearest Neighbor
}
\DeclareAcronym{GA}{
	short = GA,
	long = Genetic Algorithm
}
\DeclareAcronym{PDR}{
	short = PDR,
	long = Packet Delivery Ratio
}
\DeclareAcronym{E2E}{
	short = E2E,
	long = End-to-End
}

\DeclareAcronym{ANN}{
	short = ANN,
	long = Artificial Neural Network
}
\DeclareAcronym{RF}{
	short = RF,
	long = Random Forest
}
\DeclareAcronym{SVM}{
	short = SVM,
	long = Support Vector Machine
}
\DeclareAcronym{DSRC}{
	short = DSRC,
	long = Dedicated Short-Range Communications
}
\DeclareAcronym{VLC}{
	short = VLC,
	long = Visible Light Communication
}

\begin{document}

\include{cover}

% This is my customize variable (this removed - commented)
% \setlength{\parindent}{0pt}   % No indentation
% \setlength{\parskip}{0.8em}   % Gap between paragraphs


\renewcommand{\thepage}{\roman{page}} % Start page numbering in roman

\chapter*{Abstract}
\begin{comment}
Software defined vehicular networks (SDVN) allow vehicles to communicate effectively although they are susceptible to gray hole attacks, which are malicious nodes that use selective packet dropping to disrupt on-traffic and safety. The Traditional approaches that are in use are not completely effective or work well since they are not trustful, flexible, and offer real time protection in dynamic networks (dynamic network).

\mypar
This project proposes a secure mechanism to locate and prevent gray hole attacks in SDVNs. The blockchain maintains an immutable, secure register of the actions of the nodes, thus everything is transparent and just. Large Language Models (LLMs) observe the patterns in the network in real-time to identify threats and generate intelligent solutions, such as new routing policies. Federated Learning allows vehicles to learn models collaboratively without exchanging personal information, enabling detection to be improved and privacy maintained.

\mypar
There are three attack scenarios in the system. They are fixed rate selective dropping, time based intermittent dropping and target specific intelligent gray hole attacks. NS3 and SUMO will be used to perform simulations and Hyperledger Fabric will be used to implement the blockchain.

\mypar
In this project, We will verify such key results as the ratio of packet delivery, latency and the reliability of the network in this project. This project aims at a helpful, secure and expandable system that will ensure SDVN communication is more dependable and secure to smart cities and self driving vehicles without malicious communication.
\end{comment}
% ============================================================
% REPLACE ABSTRACT
% ============================================================

Grey hole attacks in Software-Defined Vehicular Networks (SDVNs) present a
uniquely challenging security threat: malicious nodes selectively suppress
packet forwarding while mimicking normal mobility-induced packet loss, evading
detection across both the data plane and the SDN controller plane. Existing
solutions address only subsets of the attack space, omit controller-plane
vectors, and lack the post-quantum cryptographic protection and lightweight
deployability required for real-world vehicular deployment. This paper presents
SHIELD-GH, a dual-mode detection and mitigation framework that formalizes six
distinct grey hole attack signatures across two network planes and provides
both a rule-based lightweight mode and a full AI-cryptographic mode. The
lightweight mode employs mobility-corrected sliding-window
signature detection covering fixed-rate, intermittent, and
target-specific variants across both the data plane and
controller plane for low-overhead deployment. The full mode integrates a fine-tuned Large Language Model (LLM) with an RSU-maintained, blockchain-verified Federated Learning (FL) aggregator for semantic anomaly detection, combined with post-quantum cryptographic mitigation employing
a controller-independent PQC Logical Key Hierarchy (PQC-LKH)
with Distributed Key Generation, ML-DSA-87 signatures,
VRF-based dynamic endorser selection, zk-SNARK-based packet
forwarding proofs, and $(k,n)$-threshold signatures for
collectively authorized node isolation without controller
key monopoly. A novel Mobility-Aware Trust Decay (MATD) mechanism,
Dual-Evidence Blockchain Smart Contract (DEBSC), and
multi-controller flat architecture with blockchain-backed
trust scoring and ordered failover ensure that RSU handoff
ambiguity and controller compromise do not undermine
detection integrity.

% Comment out this line as you don't have results yet.

% Simulation results in NS-3 with SUMO mobility traces and Hyperledger Fabric blockchain demonstrate that SHIELD-GH achieves a detection accuracy exceeding X\% across all six attack variants, with false positive rate below Y\%, while maintaining packet delivery ratio above Z\% and end-to-end delay within acceptable bounds for safety-critical vehicular applications.

\newpage

% Table of Contents
\tableofcontents
\newpage

% List of Figures
\listoffigures
\newpage

% List of Tables
\listoftables
\newpage

% Acronyms
\addcontentsline{toc}{chapter}{Acronyms} % Add to table of contents
\acuseall % Use all acronyms to ensure they appear in the list
\printacronyms
\newpage

% List of Notations
\addcontentsline{toc}{chapter}{List of Notations}
\chapter*{List of Notations}

\noindent The following table defines all mathematical symbols used in the equations throughout this document. Symbols are grouped by the subsystem in which they are primarily defined. Units are given where dimensionally applicable; dimensionless or structural quantities are marked with~``---''.

\vspace{0.8em}

\begin{longtable}{p{2.6cm} p{9.2cm} p{1.9cm}}
\hline
\textbf{Symbol} & \textbf{Description} & \textbf{Units} \\
\hline
\endfirsthead
\multicolumn{3}{l}{\textit{(Continued from previous page)}} \\[2pt]
\hline
\textbf{Symbol} & \textbf{Description} & \textbf{Units} \\
\hline
\endhead
\hline
\multicolumn{3}{r}{\textit{(Continued on next page)}} \\
\endfoot
\hline
\endlastfoot

%--- Vehicular Network ---
\multicolumn{3}{l}{\textit{\textbf{Vehicular Network}}} \\[2pt]
$v_i$ & Vehicle node identifier (index $i$) & --- \\
$s_i(t)$ & Speed of vehicle $v_i$ at time $t$ & m\,s$^{-1}$ \\
$R_\mathrm{RSU}$ & RSU coverage radius & m \\
$\Delta t_\mathrm{ho}$ & Average duration of an RSU handoff transition & s \\
$\Delta t$ & Observation slot duration & s \\
$N$ & Total number of vehicles in the network & --- \\
$M$ & Number of registered SDN controllers & --- \\
$\mathcal{H}_i$ & Set of all historical RSU interactions of $v_i$ recorded on the blockchain ledger & --- \\[6pt]

%--- Packet Forwarding and PDR ---
\multicolumn{3}{l}{\textit{\textbf{Packet Forwarding and PDR}}} \\[2pt]
$n_i^{\text{fwd}}(\tau)$ & Number of packets forwarded by $v_i$ in time slot $\tau$ & packets \\
$n_i^{\text{rx}}(\tau)$ & Number of packets received by $v_i$ in time slot $\tau$ & packets \\
$n_i^{\text{drop}}(t)$ & Number of packets dropped by $v_i$ at time $t$; equals $n_i^{\text{rx}}(t) - n_i^{\text{fwd}}(t)$ & packets \\
$W$ & Sliding observation window length & slots \\
$\mathrm{PDR}_i(t,W)$ & Observed Packet Delivery Ratio of $v_i$ over window $W$ ending at $t$ & [0,\,1] \\
$\widehat{\mathrm{PDR}}_i(t,W)$ & Mobility-corrected PDR after adding back expected handoff loss & [0,\,1] \\
$\overline{\mathrm{PDR}}_i(W)$ & Mean PDR of $v_i$ over window $W$ & [0,\,1] \\
$\delta_i(t)$ & Instantaneous drop rate of $v_i$ at slot $t$; equals $1 - \mathrm{PDR}_i(t,1)$ & [0,\,1] \\
$\sigma^2_i(W)$ & PDR variance of $v_i$ over window $W$ & --- \\[6pt]

%--- Mobility-Aware Correction ---
\multicolumn{3}{l}{\textit{\textbf{Mobility-Aware Correction (MATD)}}} \\[2pt]
$\rho_\mathrm{ho}(v_i,t)$ & Expected handoff-induced packet loss rate for $v_i$ at time $t$ & [0,\,1] \\
$\rho_\mathrm{max}$ & Worst-case calibrated handoff loss rate & [0,\,1] \\
$\lambda_s$ & Mobility decay coefficient in the MATD exponential penalty & --- \\[6pt]

%--- Vehicle Trust and Blockchain Reputation ---
\multicolumn{3}{l}{\textit{\textbf{Vehicle Trust and Blockchain Reputation}}} \\[2pt]
$\mathcal{T}_i(t)$ & Instantaneous Bayesian trust score of vehicle $v_i$ & [0,\,1] \\
$\alpha,\;\beta$ & Beta distribution prior parameters for the trust model & --- \\
$\mathcal{T}_i^{\text{mob}}(t)$ & Mobility-adjusted trust score of $v_i$ (output of MATD) & [0,\,1] \\
$\mathcal{R}_i(t)$ & Blockchain reputation score of $v_i$, aggregated over $\mathcal{H}_i$ & [0,\,1] \\
$\theta_R$ & Reputation deficit isolation threshold for the DEBSC statistical gate & [0,\,1] \\
$\Lambda_i(t)$ & Suspicion level: number of consecutive rounds in which the statistical gate has fired for $v_i$ & rounds \\
$W_s$ & Sliding window used for suspicion level computation & rounds \\
$\lambda_1,\;\lambda_2$ & Suspicion level thresholds triggering graduated responses 1 and 2 & rounds \\
$N_\mathrm{min}$ & Minimum blockchain interactions required to exit probationary mode & --- \\[6pt]

%--- Attack Detection Signatures ---
\multicolumn{3}{l}{\textit{\textbf{Attack Detection Signatures (S1--S6)}}} \\[2pt]
$\tau$ & Generic time slot index variable & slot \\
$\tau_f$ & Minimum acceptable forwarding ratio (Signature S1: DP-FR) & [0,\,1] \\
$\varepsilon_f$ & Maximum allowable PDR variance (Signature S1: DP-FR) & --- \\
$m_i(\tau)$ & Binary drop indicator for $v_i$ at slot $\tau$; equals $\mathbf{1}[\mathrm{PDR}_i(\tau,1)<\tau_{it}]$ & \{0,\,1\} \\
$\tau_{it}$ & Per-slot PDR threshold for the binary drop indicator (Signature S2: DP-IT) & [0,\,1] \\
$T^*$ & Candidate autocorrelation period & slots \\
$T_\mathrm{min},\;T_\mathrm{max}$ & Expected attack cycle range for the autocorrelation test & slots \\
$\gamma_{it}$ & Autocorrelation significance threshold (Signature S2: DP-IT) & [0,\,1] \\
$\tau_{ts}$ & KL-divergence threshold for target-specific detection (Signature S3: DP-TS) & nat \\
$\tau_c$ & Flow rule drop probability threshold (Signature S4: CP-FR) & [0,\,1] \\
$\gamma_c$ & Autocorrelation significance threshold (Signature S5: CP-IT) & [0,\,1] \\
$\tau_\mathrm{sub}$ & Lower bound of the sub-threshold suspicious drop probability range & [0,\,1] \\
$\mathcal{F}_c(t)$ & Set of active flow rules installed by controller $c$ at time $t$ & --- \\
$\mathcal{W}_\mathrm{BC}$ & Blockchain-stored Flow Rule Whitelist (smart contract) & --- \\
$\mathcal{C}_\mathrm{safety}$ & Set of safety-critical traffic classes (emergency, collision, ITS) & --- \\[6pt]

%--- Controller Trust and Anomaly Scoring ---
\multicolumn{3}{l}{\textit{\textbf{Controller Trust and Anomaly Scoring}}} \\[2pt]
$\Psi_c(t)$ & Aggregate controller behavioural anomaly score at time $t$ & count \\
$W_c$ & Controller observation window & slots \\
$\psi_\mathrm{thresh}$ & Threshold for the aggregate anomaly score to trigger a trust decrement & count \\
$\mathcal{T}_c(t)$ & Per-controller trust score (initialised at 1 at bootstrap) & [0,\,1] \\
$\delta_\mathrm{sig}$ & Trust decrement weight applied per confirmed S4--S6 trigger & --- \\
$\delta_\mathrm{agg}$ & Trust decrement weight applied per aggregate anomaly breach & --- \\
$\theta_c$ & Controller trust eligibility threshold for segment assignment and failover & [0,\,1] \\
$T_q$ & Timeout for RSU reachability check in degraded mode & s \\[6pt]

%--- ZKP Forwarding Proofs ---
\multicolumn{3}{l}{\textit{\textbf{ZKP Forwarding Proofs}}} \\[2pt]
$C_i$ & Pedersen commitment to the forwarded count $n_i^{\text{fwd}}$ by vehicle $v_i$ & --- \\
$r_i$ & Random Pedersen commitment blinding factor chosen by $v_i$ & --- \\
$g,\;h$ & Public group generators for the Pedersen commitment scheme & --- \\
$p$ & Large prime group modulus & --- \\
$\pi_i$ & Zero-knowledge forwarding proof generated by vehicle $v_i$ & --- \\
$\Pi_\mathrm{ZKP}(v_i,t)$ & Three-state ZKP evidence result: \texttt{PASS}, \texttt{FAIL}, or \texttt{ABSENT} & --- \\
$T_\mathrm{zkp}$ & ZKP proof submission deadline per observation window & s \\
$\widehat{n}_i^{\text{fwd}}(t)$ & RSU-independently-observed forwarded packet count for $v_i$ at time $t$ & packets \\
$\varepsilon_\mathrm{obs}$ & RSU cross-reference consistency tolerance & packets \\[6pt]

%--- Post-Quantum Cryptographic Primitives ---
\multicolumn{3}{l}{\textit{\textbf{Post-Quantum Cryptographic Primitives}}} \\[2pt]
$\mathbf{pk}$ & CRYSTALS-Kyber public key & --- \\
$\mathbf{sk}$ & CRYSTALS-Kyber private key & --- \\
$K$ & Kyber-encapsulated session key & --- \\
$\mathbf{pk}_c,\;\mathbf{sk}_c$ & Controller CRYSTALS-Dilithium signing key pair & --- \\
$\sigma$ & Dilithium digital signature & --- \\
$\mathcal{M}$ & SDN flow modification command (FlowMod message) & --- \\
$\sigma_j$ & Partial Dilithium threshold signature from RSU $j$ & --- \\
$\sigma^*$ & Combined $(k,n)$-threshold aggregate signature & --- \\
$\mathbf{pk}_\mathrm{group}$ & Collective public key of the $(k,n)$ RSU threshold scheme & --- \\
$k$ & Minimum RSU co-signature threshold for blacklisting and re-keying & RSUs \\
$n$ & Total number of RSUs in the threshold scheme & RSUs \\[6pt]

%--- VRF-Based Endorser Selection ---
\multicolumn{3}{l}{\textit{\textbf{VRF-Based Dynamic Endorser Selection}}} \\[2pt]
$\mathbf{pk}_j,\;\mathbf{sk}_j$ & Dilithium public/private key pair of RSU $j$, used for VRF evaluation and partial signing & --- \\
$\mathcal{T}_{r_j}(t)$ & Trust score of RSU $j$ at time $t$, maintained by BFT consensus on the blockchain & [0,\,1] \\
$\theta_{\text{RSU}}$ & Minimum trust score required for RSU endorser candidacy & [0,\,1] \\
$\mathcal{H}_j$ & Set of blockchain-recorded interactions of RSU $j$ (analogous to $\mathcal{H}_i$ for vehicles) & --- \\
$\mathcal{E}(t)$ & Eligible endorser pool: RSUs satisfying both the trust threshold and minimum observation count & --- \\
$s_{tx}$ & Per-transaction VRF seed, bound to tx\_id, block height, timestamp, and previous block hash & --- \\
$\beta_j$ & 256-bit pseudorandom VRF output produced by RSU $j$ & \{0,1\}$^{256}$ \\
$\pi_j^{\text{VRF}}$ & Verifiable proof that $\beta_j$ was computed correctly from $\mathbf{sk}_j$ and $s_{tx}$ & --- \\
$k_{\text{end}}$ & Minimum number of endorsers selected per transaction via VRF ranking & RSUs \\
$\alpha_{\text{end}}$ & Endorser pool fraction parameter; $\alpha_{\text{end}} > 1/3$ guarantees at least one honest endorser & --- \\
$f_{\max}$ & Maximum Byzantine RSUs in the eligible pool: $\lfloor(|\mathcal{E}(t)|-1)/3\rfloor$ & RSUs \\
$\Omega(t)$ & VRF-selected endorser set for the current transaction: top-$k_{\text{end}}$ RSUs ranked by $\beta_j$ ascending & --- \\[6pt]

%--- Federated Learning ---
\multicolumn{3}{l}{\textit{\textbf{Federated Learning}}} \\[2pt]
$\mathbf{w}$ & FL model weight vector & --- \\
$\mathbf{w}^{(r)}$ & Global FL model weights at aggregation round $r$ & --- \\
$\Delta\mathbf{w}_i$ & Local gradient update submitted by vehicle $v_i$ & --- \\
$\mathcal{D}_i$ & Local labelled forwarding-log dataset of vehicle $v_i$ & samples \\
$\mathcal{A}$ & Set of vehicles whose gradient updates passed the blockchain integrity check & --- \\
$\mathcal{L}_i(\mathbf{w})$ & Local cross-entropy loss function of vehicle $v_i$ & --- \\
$\mathcal{C}_i^\mathrm{BC}$ & On-chain hash commitment for gradient update $\Delta\mathbf{w}_i$ & --- \\[6pt]

%--- LLM Threat Scoring ---
\multicolumn{3}{l}{\textit{\textbf{LLM Semantic Threat Scoring}}} \\[2pt]
$\mathbf{x}_i^{(t)}$ & Tokenised blockchain forwarding log sequence of $v_i$ at time $t$ & --- \\
$\boldsymbol{\theta}$ & Fine-tuned LLM model parameters & --- \\
$\mathcal{Q}_i(t)$ & LLM semantic threat score (malicious class softmax probability) for $v_i$ & [0,\,1] \\
$\varepsilon_u$ & Edge LLM confidence threshold for escalation to Tier 2 cloud LLM & [0,\,1] \\[6pt]

%--- Detection Fusion ---
\multicolumn{3}{l}{\textit{\textbf{Detection Fusion}}} \\[2pt]
$\mathcal{S}_\mathrm{total}(v_i)$ & Maximum binary attack signature score across all six signatures S1--S6 & \{0,\,1\} \\
$\hat{y}_i(t)$ & Final binary detection verdict for $v_i$ ($1 =$ malicious, $0 =$ benign) & \{0,\,1\} \\
$\mu_1,\;\mu_2,\;\mu_3$ & Fusion weights for rule-based score, LLM score, and reputation deficit; $\mu_1+\mu_2+\mu_3=1$ & --- \\
$\theta_\mathrm{det}$ & Weighted fusion threshold for the final detection decision & [0,\,1] \\[6pt]

%--- PQC-LKH and DKG ---
\multicolumn{3}{l}{\textit{\textbf{PQC-LKH Group Re-Keying and DKG}}} \\[2pt]
$K_\mathrm{grp}$ & Group session key distributed to all non-isolated vehicles & --- \\
$\mathcal{K}_j$ & Set of Kyber path keys held by vehicle $v_j$ in the PQC-LKH tree & --- \\
$k_\mathrm{key}$ & Re-keying co-signature threshold (controller plus RSUs) & --- \\
$\mathbf{pk}_\mathrm{root}$ & DKG aggregate root public key for the PQC-LKH tree & --- \\
$\sigma_\mathrm{rekey}^*$ & Threshold re-keying authorisation aggregate signature & --- \\
$\mathcal{P}$ & Set of DKG participants (active controller plus registered RSUs) & --- \\

\end{longtable}

\newpage

\renewcommand{\thepage}{\arabic{page}} % Start page numbering in arabic 
\setcounter{page}{1} % start page numbering from 1

% Main Content
\chapter{Introduction}
\section{Background}

\subsection{Evolution of Vehicular Networks}
Intelligent Transportation Systems (ITS) have brought a new generation of connected vehicles, which is able to communicate with each other and roadside infrastructure. Vehicular Ad Hoc Networks (VANETs) were the foundation of early vehicular communication systems and provided simple vehicle-to-vehicle (V2V) and vehicle-to-infrastructure (V2I) communication. Nevertheless, VANETs have a number of challenges because of high vehicle mobility, high topology changes, and lack of centralized network control.

\mypar
Software-Defined Vehicular Networks (SDVNs) have been proposed in order to address these shortcomings by extending the concept of software defined networking to motor vehicle contexts. SDVNs decouple the control plane and data plane to enable centralized controllers to make routing decisions, traffic flow, and resource allocation more efficiently. This is a programmable and flexible architecture that is needed to support safety-critical ITS applications like collision avoidance systems, traffic management, and emergency response services. \cite{ieee10402120, islam2021sdvn}.


\subsection{Security Challenges and Grey Hole Attacks in SDVN}

Since SDVNs are becoming more and more significant in the context of road safety and autonomous driving, their safety and reliability have become one of the most significant concerns. The mobile nature, high velocity of vehicles and intermittent wireless connections are features of a dynamic topology that expose the vulnerability to security that cannot be completely addressed by conventional security measures.

\mypar
The Grey Hole attack is one of the security threats that are hard to detect. In this attack, malicious nodes are used to selectively drop data packets as they pass through other nodes to ensure that the network looks normal. Such selective packet dropping may interfere with vital communication, including emergency alerts and collision warnings, without necessarily arousing suspicion. Consequently, the attacks of Grey Hole may cause a considerable decrease in network performance and reliability that can be a serious threat to the safety of traffic and efficient transportation. \cite{rathod2017survey}.

% This is the Introduction Section of your project proposal report. Provide background information and context for your project. For example, \ac{AI} and \ac{ML} are transforming industries \cite{smith2020ai}.


\section{Introduction to the Problem}

One of the significant threats of SDVNs is grey hole attacks. In this attack, a set of data packets are selectively dropped by malicious vehicles as they forward them, and they look normal. This may interfere with crucial messages, including emergency messages, collision messages, and traffic messages. SDVNs are extremely mobile and have dynamic topology and it is quite impossible to detect such attacks. This issue is critical to the safe, reliable, and efficient communication on vehicles.


\section{Limitations of Existing Work}

The previous methods of ensuring vehicular networks are usually based on hard and fast rules, signature based detection, or simple machine learning. The traditional approaches find it difficult to identify selectively dropping attacks by the Grey Hole since the sporadic nature is similar to the normal network variations, and they are not able to respond well to the evolving traffic conditions. Moreover, most of the solutions need centralized data gathering or intensive computing, which pose privacy threats and scalability issues in dynamic vehicular settings \cite{ijaer2018, arizaga2025}.

\section{Problem Statement}
% Clearly define the problem the project aims to solve.
\begin{comment}
Software-Defined Vehicular Networks can be attacked by selective dropping of Grey Hole attacks, in which malicious vehicles send packets at random to avoid being noticed. This behavior is very similar to the normal packet loss that occurs due to mobility and changing network conditions in high dynamic vehicular environments. Consequently, the current security strategies do not have effective and privacy-sensitive schemes to reliably differentiate between the Grey Hole attacks and the normal network operation, which restricts their accuracy and scalability in the real-world implementation of SDVN.
\end{comment}

Grey hole attacks in Software-Defined Vehicular Networks (SDVNs) exploit the
indistinguishability between mobility-induced packet loss and selective malicious dropping — across both data-plane forwarding nodes and controller-plane flow rule injection by an untrusted SDN controller — to silently suppress safety-critical vehicular communication, while no existing detection or mitigation framework simultaneously addresses all attack variants with the controller-independent enforcement, mobility-aware trust management, and post-quantum cryptographic security required for real-world SDVN deployment.

\section{Objectives and Scope}
% State the objectives of your project clearly.

\subsection{Objectives}

The objectives of this project are to establish a comprehensive security system that identifies and prevents both data-plane and control-plane grey hole attacks within Software-Defined Vehicular Networks using Blockchain, Federated Learning, and secure logging; to create privacy-preserving detection schemes that integrate Blockchain, Large Language Models, and Federated Learning so that vehicles can cooperate without transferring sensitive information; and to evaluate the efficiency of the proposed framework under realistic network conditions and demonstrate its effectiveness over current methods.

% --- SUPERVISOR REVISION BLOCK 18: Updated Research Questions ---
\section{Research Questions}
\label{sec:rq}

This work is guided by the following six research questions, each
corresponding to a specific experimental evaluation or ablation study:

\begin{enumerate}
  \item[\textbf{RQ1.}] Under default operating conditions, does SHIELD-GH
  detect all six grey hole attack variants with higher MCC and complete
  Attack Variant Coverage Rate compared to existing state-of-the-art
  methods evaluated at the per-variant level?

  \item[\textbf{RQ2.}] As attacker penetration and attack intensity
  jointly increase across the full two-dimensional attack space, does
  SHIELD-GH maintain superior detection quality, mitigation
  effectiveness, and false isolation safety compared to baselines?

  \item[\textbf{RQ3.}] Does SHIELD-GH maintain detection accuracy and low
  false isolation rate across the vehicle speed range from urban to
  motorway conditions, and what is the measurable contribution of the
  MATD mechanism relative to non-mobility-corrected baselines?

  \item[\textbf{RQ4.}] Does SHIELD-GH sustain detection performance and
  acceptable protocol overhead as the vehicle population scales from
  100 to 400 nodes, and does PQC-LKH empirically validate its
  $O(\log N)$ re-keying complexity advantage over unicast alternatives?

  \item[\textbf{RQ5.}] As grey hole attackers become increasingly
  evasive by simultaneously optimising intermittent period, malicious
  epoch duty cycle, and targeting selectivity, does SHIELD-GH's
  dual-mode architecture maintain detection coverage beyond the point
  where lightweight rule-based signatures fail?

  \item[\textbf{RQ6.}] What is the independent quantitative contribution
  of each individually formulated SHIELD-GH design component to
  detection and mitigation performance, and which components produce
  the largest gain in their target metric when ablated?
\end{enumerate}
% --- END SUPERVISOR REVISION BLOCK 18 ---

\subsection{Scope}

The project will utilize the simulation of a Software-Defined Vehicular Network (SDVN) and study of the attacks of the data plane and control plane, namely the Grey Hole attacks. It entails the use of Blockchain in documentation of vehicle interactions in a secure way and Federated Learning in training detecting models collaboratively and keeping vehicle information confidential. The performance of the framework in the aspects of detecting the attacks, network reliability, and efficiency is also analyzed as a part of the project. The evaluation of performance through simulation will be done on the basis of conventional network and security indicators, but actual implementation of the project in the real world and hardware installation is beyond the scope of the project.



\mypar
\section{Novelty}
\begin{comment}
The originality of the research is that it has developed a single, privacy sensitive detection method against the attacks of the Grey Hole in Software-Defined Vehicular Networks. In contrast to the current solutions, the proposed framework models packet-forwarding trust based on blockchain verified behavior, learns pattern of attack collaboratively under federated detection in the context of vehicular mobility, and approaches the problem of detecting attack as an adaptive learning problem based on Large Language Models, instead of rule-based mechanisms. Besides that, the work presents formal modeling of various behaviors of the Grey Hole attacks in realistic SDVN conditions, which allow proper detection in highly dynamic environments.
\end{comment}

% ============================================================
% REPLACED SECTION 1.6 CONTENT
% ============================================================

The novelty of this work is twofold: in the problem formulation, this is the
first work to formally characterize grey hole attacks across all six variants
arising from the SDVN-specific combination of controller-plane programmability
— including by an untrusted SDN controller capable of malicious flow rule
injection — and vehicular mobility, and to rigorously model how RSU handoff
ambiguity, trust score volatility, and non-IID traffic distributions specifically
enable grey hole attackers to evade detection in ways that are not possible in
static SDN or conventional VANET. In the solution, SHIELD-GH introduces a
dual-mode framework that uniquely combines formal attack signature-driven
rule-based detection (deployable with minimal compute) and LLM-FL full-mode
detection (providing semantic reasoning over temporal behavioral patterns),
unified under a mobility-aware trust mechanism and
multi-controller flat architecture with blockchain trust
scoring, ordered failover, and VRF-based dynamic endorser
selection, and enforced by a post-quantum cryptographic
mitigation layer that employs zero-knowledge forwarding
proofs, threshold signatures, and PQC-LKH group re-keying
with DKG-based controller-independent key generation and
CRYSTALS-family PQC — the first such integrated framework
providing controller-independent, trust-governed mitigation
enforcement for grey hole attack detection in SDVN.



\mypar
\section{Contribution}

\begin{comment}
This project makes a contribution to the field of vehicular network security through supporting scalable, privacy-sensitive, and methodologically novel infrastructure of detecting and addressing Grey Hole attacks in SDVNs. The main contributions are,

\begin{itemize}
	\item \textbf{Blockchain based verification modeling:} Constructs a model in which blockchain is not only used to log data, but the mathematical verification of the packet forwarding behavior of the vehicles is performed for detecting suspicious nodes even in the case that attacks resemble normal traffic.
	
	\item \textbf{Federated Detection:} Federated learning is used to enable vehicles to jointly learn detection models, taking mobility, resource constraints, and local traffic dynamics into account, without sacrificing privacy and enhancing detectors in evolving settings.
	
	\item \textbf{Adaptive threat analysis using LLMs:} Invents the application of Large Language Models to scanned blockchain verified logs and identifies emerging trends in the pattern of the Grey Hole attacks, and does not perform attack detection based on pre determined rules.
	
	\item \textbf{Thorough attack modeling:} Formal models of various types of attacks to the Grey Hole have been studied: fixed-rate, time-based, and target specific under realistic vehicular network settings, which enables the targeted detection and mitigation of attacks.
	
	\item \textbf{Performance evaluation framework:} Conducts experiments to check the proposed framework in realistic SDVN settings and demonstrates good detection, minimal overheads, and resilience to various attacks of the SDVN under different behaviors of Grey Hole attackers.
\end{itemize}

\end{comment}

% ============================================================
% REPLACE SECTION 1.7 CONTENT
% ============================================================

The principal contributions of this work are:

\begin{enumerate}
  \item \textbf{Six-Variant Grey Hole Attack Taxonomy and Formal Signatures:}
  We provide the first formal characterization of six grey hole attack variants
  (fixed-rate, intermittent, and target-specific, in both the data plane and
  the controller plane of SDVN) with mathematically defined detection signatures
  (Eqs.~\eqref{eq:sig_dpfr}--\eqref{eq:sig_cpts}), enabling precise, per-variant
  rule-based detection.

  \item \textbf{Dual-Mode Detection Framework:}
  A lightweight rule-based detection mode and a full LLM-FL detection mode are
  jointly proposed and evaluated, ensuring practical deployability across the
  full spectrum of vehicular hardware capability. Both modes are unified under
  the SHIELD-GH architecture.

  \item \textbf{Mobility-Aware Trust Decay (MATD):}
  A novel trust scoring mechanism that corrects for RSU handoff-induced packet
  loss as a function of vehicle speed (Eqs.~\eqref{eq:matd},~\eqref{eq:handoff_loss}),
  preventing false isolation of legitimate high-speed vehicles — an SDVN-specific
  contribution absent in all prior works.

  \item \textbf{PQC-Enforced Mitigation with ZKP Forwarding Proofs:}
  Post-quantum cryptographic mitigation combining CRYSTALS-Kyber (KEM),
  CRYSTALS-Dilithium (signatures), zk-SNARK-based packet forwarding proofs,
  and $(k,n)$-threshold signatures for collectively authorized node blacklisting,
  providing quantum-resistant security guarantees in the mitigation layer.

  \item \textbf{Dual-Evidence Blockchain Smart Contract (DEBSC):}
  A novel smart contract design that requires joint satisfaction of both a
  statistical reputation threshold and a ZKP forwarding proof failure before
  triggering node isolation (Eq.~\eqref{eq:debsc}), minimizing false positives
  in adversarial mobility conditions.

  \item \textbf{Blockchain-Verified FL Gradient Integrity:}
  FL model updates are hash-committed to the blockchain before aggregation,
  preventing poisoning attacks by malicious vehicles submitting corrupted
  gradient updates in the full detection mode (Eq.~\eqref{eq:gradient_verify}).

  \item \textbf{Multi-Controller Flat Architecture with Trust-Scored Ordered
  Failover:}
  A flat multi-controller architecture in which $M$ registered SDN controllers
  each hold a per-controller trust score $\mathcal{T}_c(t)$ maintained on the
  blockchain by RSU consensus
  (Eqs.~\eqref{eq:controller_agg}--\eqref{eq:controller_trust}), with
  zone-based ordered failover that selects the lowest-ranked eligible controller
  from each segment's blockchain-stored list (Eq.~\eqref{eq:failover_select}),
  eliminating the dominant single-controller failure mode without concentrating
  fallback authority on any one controller.

  \item \textbf{DKG-Based Threshold Key Generation Eliminating Controller Key
  Monopoly:}
  The root Kyber key pair for the PQC-LKH tree is generated jointly by the
  active controller and RSUs via Pedersen Distributed Key Generation
  (Eqs.~\eqref{eq:dkg_share}--\eqref{eq:rekey_sig}), ensuring no single entity
  holds complete key material and all re-keying broadcasts require a threshold
  co-signature from the controller and $k_{\text{key}}-1$ RSUs, eliminating
  controller key monopoly while maintaining controller participation.
\end{enumerate}


\mypar
\section{Implication}

This Project Suggests Implications to Vehicular Network Security, Studies, and Real Implementation.

From a technological standpoint, combining Blockchain with Federated Learning offers a novel method of securing vehicle networks, and the strategy is transferable to smart cities, traffic management, and other interconnected networks.
The project also has educational and research value, providing a practical demonstration of contemporary security methods that enables students and researchers to understand how distributed networks can be protected.
\textbf{Privacy Preserving Collaborative Security} is central to the design: Federated Learning allows vehicles to collaborate in detecting attacks without sharing raw personal data, keeping the system secure and scalable.
The framework improves road safety and reliability by ensuring consistent vehicle-to-vehicle communication and reducing the likelihood of attack-induced accidents or traffic disruptions.
Finally, the project establishes a foundation for future transportation systems, providing a roadmap for secure, flexible, and privacy-conscious vehicular networks.



 

\chapter{Literature Review}
\begin{comment}
Being centered around the convergence of blockchain technology, large language models (LLMs), and federated learning (FL), this literature review focuses on existing literature related to countermeasures of grey hole attacks in Software-Defined Vehicular Networks (SDVN). Owing to its dynamic network structure, a mobile network structure, and a distributed character of vehicular networks, routing integrity in SDVNs is vulnerable to grey hole attacks, which represent a sort of black hole attack as nodes might reject packets of interest selectively. SDVNs remain vulnerable to attacks on packet routing and availability since they adopt software-defined networks (SDN) and aim at building a centralized control structure alongside a distributed routing structure. The literature review is divided into sub-headings in relation to key tech: FL for securing and supporting a collaborative approach of handling large datasets, LLMs for superior threat analysis capabilities, as well as blockchain technology for a secure distributed network structure.


\section{Blockchain in Vehicular Networks}


Due to its unique decentralized and tamper-resistant properties, blockchain technology nowadays has been one of the great spotlights for enhancing the security issues in both VANET and SDVN. By storing immutable ledgers, it makes blockchain enable data integrity, transparency, and trust between involved vehicular nodes. In, a blockchain-based smart contract framework was presented in order to avoid black hole and grey hole routing attacks in VANETs. For the detection and isolation of malicious vehicles, this scheme utilized self-classification blockchain-based contracts (SCBC) and voting-classification blockchain-based contracts (VCBC). The network performance is assessed with packet delivery ratio (PDR) and throughput (TP) and the results showed good resiliency even in active attack conditions. It could be verified on node behavior, thanks to blockchain consensus mechanisms, without any trust in centralized authority \cite{alabdulatif2024blockchain}.

\mypar
In a different study, blockchain technology was integrated with federated Q-learning to provide a secure method of sharing data in VANETs. This system used extended elliptic cryptography (EX-ECC) for encryption and the interplanetary file system (IPFS) to provide decentralized file storage, with delegated practical Byzantine fault tolerance (DPBFT) used as the consensus algorithm. In order to provide privacy in VANETs, decentralization of updating was used in a manner that limits privacy in centralized architectures. The system showed high throughput and low latency in simulations conducted on up to 100 vehicles. In addition, the blockchain’s immutable ledger of transactions also provides support to ensure the tracing of evil acts, such as selective packet dropping in grey hole attacks, performed in VANETs. The major shortcoming of most studies, including VANETs, is that blockchain technology limits vulnerabilities associated with single point failures in software-defined VANET architectures, including those associated with software-defined network (SDN) controllers.


\section{Decentralized Security through Blockchain and Federated Learning}

A recent trend in SV research includes decentralizing the architecture of intrusion detection systems (IDS), which is prone to the problem of single points of failure and also a privacy concern as it involves the collection of data in a centralized manner. To counter these drawbacks, federated learning (FL) has gained popularity as a decentralized learning technique in VANETs. FL helps vehicles learn together to develop an IDS model from their respective generated data without revealing their vehicle details to others. To further protect against inference attacks, various studies have utilized differential privacy tools that involve the injection of controlled Gaussian noise \cite{ahmed2024federated}.

\mypar
At the same time, there has been an integration in FL frameworks utilizing blockchain technology, and more specifically, the use of consortium blockchains, to ensure trust and accountability among FL nodes. Specifically, consortium blockchains are employed to record and aggregate model updates into an immutable, transparent ledger, thus ensuring data integrity and unaltered data during their storage. Still, FL frameworks that support decentralized learning have a major deficiency, namely, they are susceptible to poisoning attacks carried out by malicious nodes that may provide compromised or deceptive updates. Recently, new frameworks, referred to as Zero-X, have been introduced, which include enhanced consensus protocols, such as Proof-of-Accuracy (PoA), to assess FL model updates based on their efficacy in enhancing global accuracy, thus making FL IDSs more resilient regarding decentralized learning configurations within highly dynamic vehicular environments.

\subsection{Large Language Models in Security}

Recently, research has emerged to apply the concept of Large Language Models (LLMs) to the security domain using models such as BERT (Bidirectional Encoder Representations from Transformers). In particular, attention has grown towards using models such as BERT in intrusion detection systems (IDS) with a focus on sequence-based anomaly detection in SDNVs. A research contribution specifically used a privacy-preserving framework for an IDS system utilizing Federated learning with BERT (FL-BERT) in software-defined VANETs. In this proposal for improvement, attack behaviors in terms of different attack patterns such as grey hole attacks with patterns of intermittent packet drops have been incorporated by modeling network traffic logs as a sequence to capture dependencies in such patterns. Based on context provided by VeReMi, FL-BERT provided better accuracy for IDS over traditional ML methods such as Random Forest and Support Vector Machines by maintaining data locality via federated learning \cite{zhang2025anomalous}.

\mypar
LLMs have been very effective in understanding semantic relations in network flow data, which helps in recognizing subtle routing anomalies in network flows to detect potentially selectively malicious activities as opposed to other kinds of threats. In correlated research studies, black hole attack prevention in VANETs has been done using threshold-based packet validation; nevertheless, this modeling has been found to have potentially vast improvement using LLMs with smart understanding capabilities for forged route request messages as well as anomalous control massages. In addition to this, edge models for unstructured or complex data representations like encrypted flow data patterns and obfuscated attack signs have been found to have immense applicability for LLMs as opposed to other models for determining corresponding data patterns in SDNs. Despite such benefits, research applications of LLM models in vehicular network security are still in their infancy and mostly involve offline process analysis based on their high computation as well as memory complexities. While potentially allowing real-time validation of flow rule data in SDNs by their collaboration with SDN controller models for grey hole attack validations, most correlated research models have failed to address edge device complexities in SDVN.


\subsection{Federated Learning in VANETs}

Federated learning (FL) has become a promising technique for jointly learning models in Vehicular Ad Hoc Networks (VANETs) and Software-Defined Vehicular Networks (SDVNs) to perform learning on models in a distributed manner without sharing data directly. A learning technique that is distributed in nature helps overcome the problem of intrusion detection systems related to user privacy in a vehicular network where accurate tracking is required for vehicles with proper location details without violating user privacy rules. A federated learning framework based on blockchain for VANET was designed using a blockchain-based federated Q-learning system for learning models based on Q-learning approach using deep Q-network learning models based on local observations and using InterPlanetary File System (IPFS) for model update off-chain storage with negligible time delay for threat detection with a high level of network throughput with an optimized blockchain architecture. This technique becomes significant for preventing Grey Hole attacks on a vehicular network where learning is performed using models based on federated learning with varying patterns during attacks based on different models for different vehicles based on different routes with accurate location tracking for preventing attacks on user privacy in a vehicular network \cite{ahsan2024flbert, malik2022blackhole}.

\mypar
Another research effort introduced an efficient dynamic method for detecting black hole attacks in VANETs by analyzing forged route request packets and applying threshold-based calculations within the AODV routing protocol. This technique aligns well with the principles of federated learning, as FL supports continuous and distributed model updates across network nodes in response to changing attack behaviors. In the context of software-defined VANETs, FL-BERT has been adopted to federate BERT-based intrusion detection models, enabling privacy-preserving training on misbehavior datasets such as VeReMi. This federated approach achieved higher attack classification accuracy compared to non-federated learning methods, demonstrating FL’s effectiveness in SDN-based vehicular environments. While FL significantly reduces the risks associated with data centralization in SDN architectures by keeping datasets local to each vehicle, it also introduces challenges, particularly communication overhead in high-mobility scenarios. Although existing studies primarily evaluate FL-based solutions in simulated VANET settings and report improvements in packet delivery ratio (PDR) and end-to-end delay, the combined integration of blockchain for tamper-proof model updates and large language models for advanced feature extraction remains largely unexplored, especially for SDVN-specific grey hole attack mitigation.


\section{Gaps in Literature}

\begin{itemize}
	\item \textbf{Limited Technology Integration:} While most studies concentrate on individual technologies (e.g., blockchain for consensus or FL for privacy), very few integrate blockchain, LLMs, and FL into a cohesive framework for mitigating grey hole attacks in SDVN. This results in fragmented solutions that ignore synergistic advantages like secure, intelligent, and privacy-preserving detection.
	
	\item \textbf{Inadequate Attention to the Grey Hole Specificity:} While black hole attacks are frequently addressed, grey hole variants—which have selective and erratic packet dropping—are not as well studied. This is especially true in SDVN, where SDN controllers may be used to cause partial disruptions, leaving gaps in adaptive detection mechanisms.
	
	\item \textbf{Resource Constraints in Edge Devices:} Real-time applicability in dynamic SDVN environments is hampered by existing works that seldom take into account the computational overhead of LLMs like BERT in resource-limited vehicular nodes, frequently assuming unlimited processing power.
	
	\item \textbf{Absence of Real-World Datasets and Simulations:} The majority of evaluations rely on simulated datasets (such as VeReMi), with little testing conducted on actual SDVN deployments. As a result, factors that impact the effectiveness of attack detection, such as heterogeneous vehicle hardware, signal interference, and varying mobility patterns, are ignored.
	
	\item \textbf{Privacy vs. Performance Trade-offs:} FL-based methods protect privacy but frequently sacrifice model accuracy because of non-IID data distributions in VANETs; there are gaps in optimizing these trade-offs, particularly when combined with the overhead of blockchain technology and the complexity of LLMs.
	
	\item \textbf{Scalability in High-Mobility Scenarios:} Research demonstrates strong performance in small-scale simulations (e.g., 100 vehicles), but scalability to large SDVN with frequent topology changes is still unresolved, which could result in higher latency and consensus failures in blockchain-FL hybrids.
	
	\item \textbf{Lack of Holistic Threat Models:} The robustness of suggested solutions against combined or evolving threats is limited by the literature's lack of comprehensive threat models that incorporate SDN-specific vulnerabilities (such as controller hijacking) alongside grey hole attacks.
\end{itemize}
\end{comment}

% ============================================================
% REPLACE ALL OF CHAPTER 2 WITH THE FOLLOWING
% ============================================================

%\chapter{Literature Review}
%\label{chap:litreview}

This chapter reviews work directly relevant to grey hole attack detection and
mitigation in vehicular and software-defined networks, followed by a review of
LLM- and federated-learning-based approaches in vehicular security. The review
is deliberately focused on works with strong relevance to the proposed framework
in order to derive verifiable novelty claims.

\section{Grey Hole Attack Detection and Mitigation in Vehicular Networks}
\label{sec:lit_greyhole}

Grey hole attacks — in which a malicious node selectively drops a fraction of
packets while appearing to participate normally — are considerably less studied
than black hole attacks, despite posing comparable threat to safety-critical
vehicular communication.

Kaur \textit{et al.}~\cite{kaur2022greyhole} proposed a detection and prevention
mechanism for grey hole attacks in VANET using the AODV routing protocol
simulated in NS-2. The method monitors forwarding behaviour and isolates
malicious nodes when packet drop patterns are identified. While the results
show improvement in PDR and throughput after detection, the approach relies on
AODV-specific control messages that are unavailable in SDN-based architectures,
and no machine learning or cryptographic component is employed. The work also
does not distinguish between fixed-rate and intermittent dropping variants, nor
does it address mobility-induced packet loss ambiguity.

Malik \textit{et al.}~\cite{malik2023greyhole} introduced a more comprehensive
approach targeting both Smart GHA and Sequence Number-based GHA variants in
VANETs, proposing a dynamic threshold-based detection mechanism evaluated on
NS-3. The study identifies two distinct attack behaviours and applies
variant-specific thresholds to improve detection accuracy, achieving a PDR
improvement and a reduction in end-to-end delay. However, the solution is
evaluated in a pure VANET context without SDN control-plane separation, does not
consider intermittent or target-specific variants, and proposes no cryptographic
mitigation. Mobility-aware correction of legitimate packet loss is also absent.

Remya-krishnan and Kumar~\cite{remyakrishnan2022dtw} proposed a time-series
analysis approach using Dynamic Time Warping (DTW) to detect both smart black
hole and grey hole nodes in VANET. By modelling the sequence of dropped packets
per node as a time series and applying DTW-based similarity matching against
malicious behaviour templates, the method achieves competitive detection
performance including in the intermittent attack case. This is one of the few
works to explicitly address temporally varying grey hole behaviour. Nonetheless,
the approach requires pre-built attack templates, making it unsuitable for
novel or evolving attack patterns, and it does not address the controller-plane
attack surface introduced by SDN.

Younas \textit{et al.}~\cite{younas2022collaborative} presented a neural
network-based technique for the collaborative detection of both black hole and
grey hole attacks in VANET, using cooperative awareness messages (CAMs) as
input features to an anomaly detection model. The cooperative nature of the
detection scheme anticipates the federated learning direction, but raw feature
data is shared between vehicles rather than model updates, introducing privacy
concerns. Moreover, the solution neither employs blockchain for trust management
nor incorporates any cryptographic mitigation, and the SDN architecture is not
considered.

Alabdulatif \textit{et al.}~\cite{alabdulatif2024mitigating} employed
blockchain-based smart contracts to classify vehicles as benign or malicious
based on their packet forwarding ratios, using Self-Classification (SCBC) and
Voting-Classification (VCBC) blockchain contracts in an Ethereum environment.
Under combined black hole and grey hole attack conditions, VCBC achieves a PDR
of 78\%. The blockchain provides a decentralised trust mechanism resilient to
single points of failure. However, no AI-based detection is employed, the work
is evaluated in stable highway scenarios without mobility-induced ambiguity, and
the SDN controller plane attack vector is entirely absent.

Arízaga-Silva \textit{et al.}~\cite{arizagasilva2025ml} proposed a machine
learning-powered IDS for grey hole attack detection in VANET, evaluating
multiple classifiers — including random forests and gradient boosting — on
features extracted from NS-3 simulations. The study is among the most recent
specifically targeting grey hole detection and demonstrates that ML-based
detection significantly outperforms threshold-based methods. Nevertheless, the
approach is confined to the data plane of a conventional VANET architecture,
does not incorporate blockchain or cryptographic mechanisms, and does not address
the controller-plane attack scenario unique to SDVN.

\section{Grey Hole Attack Detection in Software-Defined Networks}
\label{sec:lit_sdn_greyhole}

The SDN architecture introduces an additional attack surface not present in
conventional VANET: a compromised or manipulated controller can install
malicious flow rules that cause selective packet dropping across entire network
segments without any individual vehicle behaving abnormally.

Hsieh and Ku~\cite{hsieh2018sdn} specifically studied time-based and
random-based grey hole attacks in SDN, where compromised switches selectively
drop packets based on control-plane instructions. They proposed a detection
method based on weighted K-Nearest Neighbour (KNN) combined with a Genetic
Algorithm (GA) applied to flow counter statistics collected from OpenFlow
switches. This is the only work identified in the literature that specifically
addresses grey hole attacks in the SDN infrastructure rather than in the
vehicular or wireless network layer. However, the method does not extend to the
vehicular mobility context, does not address intermittent or target-specific
variants, proposes no mitigation mechanism beyond detection, and does not
integrate any trust or cryptographic component.

\section{LLM and Federated Learning for Vehicular Network Security}
\label{sec:lit_llm_fl}

Hamhoum and Cherkaoui~\cite{hamhoum2024mistralbsm} proposed MistralBSM, an
LLM-based misbehavior detection system that fine-tunes the Mistral-7B model on
sequences of Basic Safety Messages (BSMs) from the extended VeReMi dataset.
Operating within an edge-cloud framework, the edge LLM processes short BSM
windows for real-time detection while a cloud LLM performs deeper analysis.
MistralBSM achieves 98\% binary classification accuracy and 96\% multiclass
accuracy, outperforming BERT- and LLAMA2-7B-based baselines. The work
demonstrates that LLMs are effective for vehicular misbehavior detection when
applied to sequential communication logs. However, MistralBSM targets position
falsification and similar beacon-level misbehavior; its adaptation to
routing-layer grey hole attacks — particularly in SDVN where forwarding logs
rather than BSMs are the input — remains unaddressed. Furthermore, no
blockchain, federated learning, or cryptographic component is employed.

Ahsan \textit{et al.}~\cite{ahsan2024flbert} introduced a privacy-preserving
IDS for software-defined VANET using Federated Learning with BERT (FL-BERT),
training local BERT-based sequence classifiers on the VeReMi dataset at each
vehicle without sharing raw data. The federated model achieves higher attack
classification accuracy than non-federated baselines, including Random Forest
and SVM, and represents the closest existing work to the proposed full-mode
detection pipeline. Nonetheless, the work does not address controller-plane grey
hole attacks, proposes no cryptographic mitigation, and does not account for
the mobility-induced packet loss ambiguity that is specific to SDVN.

Ahmed \textit{et al.}~\cite{ahmed2024blockchain} integrated blockchain with
federated Q-learning for secure data sharing in VANET, using Extended Elliptic
Curve Cryptography (EX-ECC) for encryption and the InterPlanetary File System
(IPFS) for decentralised storage, achieving high throughput and low consensus
latency. While the work establishes that blockchain-FL integration is feasible
in vehicular environments, it does not address grey hole detection specifically,
does not incorporate LLMs, and the cryptographic scheme is not post-quantum
secure.

% ============================================================
% INSERT AS NEW SECTION 2.5 — BEFORE GAP ANALYSIS
% This section justifies the proposed methodology by reviewing
% how similar techniques are used for OTHER attacks/problems.
% ============================================================

\section{State-of-the-Art in Similar Detection and Mitigation Techniques}
\label{sec:lit_techniques}

Beyond the grey hole attack literature, understanding how blockchain,
federated learning, and LLMs have been applied to analogous security
problems in SDN and vehicular environments is essential for justifying
the proposed methodology and positioning its design contributions.

\subsection{Blockchain Integration with Software-Defined Networking}
\label{sec:lit_blockchain_sdn}

Nguyen \textit{et al.}~\cite{nguyen2021blockchain} provided a comprehensive
survey of blockchain technologies applied to SDN, covering both security
and non-security use cases. The survey categorises blockchain-SDN integration
into two principal paradigms: using blockchain to decentralise and secure
the SDN control plane (mitigating its single-point-of-failure vulnerability),
and using blockchain-based smart contracts to authenticate and audit the
flow rules installed by SDN controllers. Of particular relevance to the
proposed framework, the survey identifies flow rule integrity verification
and node trust management as open problems in SDN security where blockchain
offers the most direct architectural benefit. The survey also identifies
key challenges: the high latency of blockchain consensus is incompatible
with the sub-millisecond flow update requirements of SDN in vehicular
environments, and naive application of public blockchains introduces
transaction costs and throughput limitations. These findings directly
motivate the design of the Dual-Evidence Blockchain Smart Contract (DEBSC)
proposed in Section~\ref{sec:signatures}: rather than logging all forwarding
events on-chain (which would generate prohibitive transaction volumes in
a high-mobility SDVN), the DEBSC triggers a blockchain transaction only
when a joint statistical and cryptographic anomaly threshold is crossed,
dramatically reducing on-chain activity while preserving the integrity
guarantees.

% ....
\begin{comment}
@article{nguyen2021blockchain,
  author    = {Nguyen, Hai Nam and Tran, Hai Anh and Fowler, Scott
               and Souihi, Sami},
  title     = {A Survey of Blockchain Technologies Applied to
               Software-Defined Networking: Research Challenges
               and Solutions},
  journal   = {IET Wireless Sensor Systems},
  volume    = {11},
  number    = {6},
  pages     = {233--247},
  year      = {2021},
  doi       = {10.1049/wss2.12031}
}

@article{li2025ddos,
  author    = {Li, [Authors to verify on IEEE Xplore]},
  title     = {{DDoS} Attack Detection in {SDN}-Assisted Federated
               Learning Environment Based on Contrastive Learning},
  journal   = {IEEE Access},
  year      = {2025},
  doi       = {10.1109/ACCESS.2025.3582445},
  note      = {Early Access. Verify author names on IEEE Xplore
               before submission.}
}
\end{comment}

\subsection{Federated Learning for Attack Detection in SDN Environments}
\label{sec:lit_fl_sdn}

Li \textit{et al.}~\cite{li2025ddos} proposed a contrastive learning-based
DDoS detection framework specifically designed for SDN-assisted FL environments,
where the SDN control plane is exploited to optimise FL model parameter
routing among participants. The framework uses a two-layer classification
architecture: an LSTM-SVM model in the first layer and a CNN-BiGRU model
in the second layer, with the second layer refined via contrastive
classification loss from the first. On a Mininet-generated SDN dataset,
the LSTM-SVM model achieves 99.75\% detection accuracy and 99.80\% recall,
while CNN-BiGRU achieves 99.36\% accuracy. Critically, the work demonstrates
that SDN's centralised flow management can itself be used to attack the FL
network — an insight that directly informs the controller-plane attack
variants (CP-FR, CP-IT, CP-TS) modelled in this work. However, the
framework does not address vehicular mobility, considers only DDoS rather
than selective dropping, and proposes no cryptographic mitigation. The
contrastive learning approach also requires large labelled datasets, which
are unavailable for grey hole attacks in SDVNs — motivating the use of
the LLM's in-context reasoning and the lightweight rule-based mode in
the proposed dual-mode framework. The high detection accuracy of 99.75\%
in a static SDN environment also serves as a methodological performance
benchmark for the proposed framework's full-mode detector.

\subsection{Post-Quantum Cryptography in Connected and Autonomous Vehicles}
\label{sec:lit_pqc_cav}

Current vehicular security standards — including IEEE~1609.2 for V2X
communications and ETSI~ITS-G5 — are built on classical public-key
cryptography (ECDSA over P-256 and ECIES), which is known to be broken
by Shor's algorithm on a sufficiently powerful quantum computer. As
connected and autonomous vehicles are expected to remain operational for
15--20~years, cryptographic mechanisms deployed today face a
harvest-now, decrypt-later threat: an adversary can record encrypted
V2X traffic today and retrospectively decrypt it once a
cryptographically relevant quantum computer becomes
available~\cite{nist2024fips203}.

In August 2024, NIST finalised the first post-quantum cryptographic
(PQC) standards: FIPS~203 (ML-KEM, based on CRYSTALS-Kyber) for key
encapsulation, and FIPS~204 (ML-DSA, based on CRYSTALS-Dilithium) for
digital signatures~\cite{nist2024fips203,nist2024fips204}. Both
algorithms are lattice-based, offering conjectured security against both
classical and quantum adversaries. CRYSTALS-Dilithium at security
level~5 (ML-DSA-87) produces signatures of 4,595~bytes and public
keys of 2,592~bytes, compared to 64~bytes for an ECDSA signature
over P-256 — a 72$\times$ increase in signature size with significant
implications for bandwidth-constrained V2X broadcast channels where
Basic Safety Messages (BSMs) and Cooperative Awareness Messages (CAMs)
are transmitted at 10~Hz per vehicle.

Benchmarking on ARM Cortex-A class processors — representative of
automotive-grade on-board units (OBUs) — shows that ML-DSA signing and
verification are computationally feasible for periodic safety-message
authentication, though the increased message size strains dedicated
short-range communication (DSRC) and C-V2X channel capacity under dense
traffic conditions~\cite{nist2024fips204}. For resource-constrained RSU
deployments, the overhead of batch signature verification across
hundreds of simultaneous vehicle transmissions remains an active
engineering challenge. Hybrid classical-PQC schemes, maintaining
backward compatibility while introducing quantum resistance, are
recommended by NIST for transitional deployments.

No work reviewed in the grey hole detection or SDVN security literature
integrates PQC mechanisms for authentication of forwarding proofs or
SDN flow rule integrity verification. This gap directly motivates the
post-quantum mitigation layer in SHIELD-GH, which employs
CRYSTALS-Kyber for session key establishment between SDN controllers
and vehicular nodes, and CRYSTALS-Dilithium for signing forwarding-proof
blockchain transactions, targeting the harvest-now threat model for
long-lived vehicular deployments.

\section{Research Gap Analysis}
\label{sec:gaps}

The literature review exposes significant gaps in detection coverage.
No existing framework addresses all three grey hole variants — fixed-rate,
intermittent, and target-specific — simultaneously within a single
SDVN-aware system.

\cite{malik2023greyhole} cover only two variants restricted to the data
plane, and \cite{remyakrishnan2022dtw} target intermittent dropping
exclusively, leaving target-specific attacks on emergency traffic
undetected.
The controller-plane grey hole, where a compromised SDN controller injects
malicious flow rules to redirect or suppress traffic, is entirely absent
from VANET-focused literature.
\cite{hsieh2018sdn} examine SDN grey holes but without vehicular context,
and no work provides controller-plane signatures or blockchain-backed
flow-rule integrity verification.
These two gaps collectively motivate the six formal attack signatures
(S1--S6) and the DEBSC smart contract in SHIELD-GH\@.

A further gap concerns the context-awareness of detection and the
deployability of proposed solutions.
No reviewed detection scheme distinguishes mobility-induced packet loss from
deliberate dropping, which causes high false-positive rates when high-speed
vehicles undergo RSU handoffs.
SHIELD-GH addresses this through the Mobility-Aware Trust Decay (MATD)
mechanism with handoff-aligned detection windows
(Eqs.~\eqref{eq:matd} and~\eqref{eq:handoff_loss}).
Equally, no prior work proposes a tiered architecture that supports both a
lightweight rule-based mode and a full AI-based mode: existing proposals
require either full inference or fixed thresholds, making them unsuitable
across all hardware tiers.
SHIELD-GH provides both modes as independently evaluable components,
enabling deployment on resource-constrained on-board units as well as
infrastructure-side edge servers.

The reviewed literature also leaves critical cryptographic and
intelligence-layer gaps unaddressed.
Post-quantum cryptography has not been applied to grey hole mitigation;
\cite{ahmed2024blockchain} rely on EX-ECC, which is vulnerable to
harvest-now-decrypt-later attacks on long-lived vehicular deployments,
motivating the adoption of CRYSTALS-Kyber (ML-KEM) and
CRYSTALS-Dilithium (ML-DSA) in SHIELD-GH\@.
Zero-knowledge proof-based forwarding verification has likewise not been
integrated into grey hole detection: without cryptographic evidence of
forwarding behaviour, a dishonest node can falsify blockchain records and
undermine the trust layer.
SHIELD-GH closes this gap through Pedersen commitment-based ZKP forwarding
proofs and a DEBSC smart contract that requires both a ZKP failure and a
statistical threshold breach before vehicle isolation.
Finally, LLM-assisted detection in vehicular security has been confined to
beacon-level misbehavior~\cite{hamhoum2024mistralbsm} or federated-learning
aggregation without semantic reasoning~\cite{ahsan2024flbert}.
No work exploits the temporal pattern-reasoning capability of LLMs at the
routing layer; SHIELD-GH fills this gap by applying a fine-tuned LLM to
tokenised blockchain forwarding logs for semantic grey hole pattern
recognition in its full detection mode.


\vspace{0.5em}
% \textbf{Related Works Comparison}
\noindent\textbf{Related Works Comparison}

\begin{table*}[!t]
\renewcommand{\arraystretch}{1.3}
\caption{Comparison of Related Works}
\label{tab:relatedwork}
\centering
\scriptsize
\begin{tabular}{|p{2.1cm}|p{1.1cm}|p{1.5cm}|p{1.4cm}|p{1.2cm}|p{1.2cm}|p{0.9cm}|p{0.9cm}|p{0.85cm}|p{0.9cm}|}
\hline
\textbf{Work} & \textbf{Network} & \textbf{Attack Variants} &
\textbf{Detection Method} & \textbf{AI/ML} & \textbf{Crypto.} &
\textbf{Mob.\ Aware} & \textbf{LW Mode} &
\textbf{Ctrl.\ Plane} & \textbf{PQC} \\
\hline
Kaur \textit{et al.}~\cite{kaur2022greyhole} (2022) &
VANET & GH (general) & AODV-based threshold & None & None & No & No & No & No \\
\hline
Malik \textit{et al.}~\cite{malik2023greyhole} (2023) &
VANET & Smart GH, Seq.-No.\ GH & Dynamic threshold & None & None & No & No & No & No \\
\hline
Remya-krishnan \& Kumar~\cite{remyakrishnan2022dtw} (2022) &
VANET & BH + GH (intermittent) & DTW time-series & None & None & No & No & No & No \\
\hline
Younas \textit{et al.}~\cite{younas2022collaborative} (2022) &
VANET & BH + GH & Neural network (collab.) & ANN & None & No & No & No & No \\
\hline
Hsieh \& Ku~\cite{hsieh2018sdn} (2018) &
SDN & GH (time- \& random-based) & KNN + GA (flow counters) & KNN & None & No & No & Partial & No \\
\hline
Alabdulatif \textit{et al.}~\cite{alabdulatif2024mitigating} (2024) &
VANET & BH + GH & Blockchain smart contracts & None & Blockchain consensus & No & No & No & No \\
\hline
Arízaga-Silva \textit{et al.}~\cite{arizagasilva2025ml} (2025) &
VANET & GH (intermittent) & ML classifiers (RF, GB) & RF, GB & None & No & No & No & No \\
\hline
Ahsan \textit{et al.}~\cite{ahsan2024flbert} (2024) &
SD-VANET & GH (general) & FL-BERT (IDS) & BERT, FL & None & Partial & No & No & No \\
\hline
Ahmed \textit{et al.}~\cite{ahmed2024blockchain} (2024) &
VANET & General security & FL Q-learning + blockchain & FL, RL & EX-ECC & No & No & No & No \\
\hline
Hamhoum \& Cherkaoui~\cite{hamhoum2024mistralbsm} (2024) &
VANET & BSM misbehavior & LLM fine-tuning (Mistral-7B) & LLM & None & No & No & No & No \\
\hline
\textbf{Proposed} & \textbf{SDVN} &
\textbf{All 6 GH variants} &
\textbf{Rule-sig (LW) + LLM-FL (full)} &
\textbf{LLM, FL} &
\textbf{Kyber, Dilithium, ZKP, TS} &
\textbf{Yes} & \textbf{Yes} & \textbf{Yes} & \textbf{Yes} \\
\hline
\end{tabular}
\vspace{1mm}
\footnotesize{BH: Black Hole; GH: Grey Hole; LW: Lightweight; Ctrl.: Controller;
Mob.: Mobility; RF: Random Forest; GB: Gradient Boosting; DTW: Dynamic Time
Warping; TS: Threshold Signatures; PQC: Post-Quantum Cryptography.}
\end{table*}


% SUPERVISOR COMMENT:
% The methodology conflates detection and mitigation throughout. The LLM role is
% described as both "detecting" and "generating routing policies" without a clear
% logical boundary. Blockchain is described only as a logging mechanism
% (mitigation-adjacent) but its detection role (baseline comparison) is
% underspecified. No attack signatures are formally defined, which is the
% foundational step for any rule-based detector. The cryptographic component
% (PyCryptodome is mentioned in tools) is completely absent from the methodology.
% A dual-mode framework (lightweight + full) is required. The following
% formulations address all of these.

\chapter{Methodology}
\label{sec:methodology}


% SUPERVISOR FIGURE INSTRUCTIONS:
% Figure A — Dual-Mode System Architecture (Block Diagram):
%   Two-column block diagram. Top labels: "LIGHTWEIGHT MODE" (left, blue)
%   and "FULL VERSION MODE" (right, orange). Left column: (1) OpenFlow Counter
%   Collector → (2) Signature Engine [S1–S6] → (3) HMAC + SDN FlowMod
%   Mitigation. Right column: (1) Blockchain Log Tokenizer → (2) parallel:
%   LLM Anomaly Scorer | FL Global Model → (3) Fusion Engine [Eq. 20] →
%   (4) PQC Mitigation Engine [Kyber + Dilithium + ZKP + Threshold Sig].
%   Shared green "Blockchain Plane" box between columns with bidirectional
%   arrows. Bottom: "Mobility-Aware Trust Decay (MATD)" feeds both columns.
%
% Figure B — Lightweight System Sequence Diagram (UML):
%   Lifelines: Vehicle (Legitimate) | Vehicle (Attacker) | RSU/OpenFlow |
%   Blockchain Ledger | SDN Controller.
%   Sequence: (1) Attacker drops packets → (2) RSU detects PDR anomaly
%   via S1–S3 → (3) RSU queries Blockchain → (4) Blockchain returns
%   historical PDR → (5) RSU computes MATD-corrected score → (6) If
%   score > threshold: RSU sends FlowMod → (7) Controller isolates
%   attacker → (8) Blockchain reputation updated.
%   Dashed vertical dividers label detection and mitigation phases.
%
% Figure C — Full-Version System Sequence Diagram (UML):
%   Same lifelines as B plus: LLM Agent | FL Aggregator.
%   Sequence: (1) Attacker intermittently drops packets → (2) Vehicles
%   generate ZKP proofs → (3) Proofs logged to Blockchain → (4) LLM
%   tokenizes log → (5) LLM produces threat score → (6) FL vehicles
%   compute local updates → (7) FL Aggregator (blockchain-verified)
%   produces global model → (8) Fusion Engine combines LLM + FL +
%   Reputation → (9) Threshold sig collected from k RSUs → (10)
%   Dilithium-signed FlowMod issued → (11) Kyber re-keying.
%
% Figure D — Cryptographic Mitigation Flowchart:
%   Start: "Detection Decision: Malicious".
%   Step 1: Retrieve ZKP Proof → "ZKP Valid?" → No: "Flag Confirmed Attacker".
%   Step 2: Collect k Partial Threshold Sigs → Combine σ* [Eq. 26] →
%           "σ* Valid?" → No: "Abort (Insufficient Evidence)".
%   Step 3 (Yes): Dilithium-sign FlowMod [Eq. 19] → Install Block Rule →
%                 Kyber Re-keying [Eq. 17] → Update BC Reputation → End.

% SUPERVISOR COMMENT:
% The current paper treats grey hole attacks in SDVN almost identically to how
% they would be treated in a generic VANET or SDN. There is no explicit analysis
% of how vehicle mobility specifically enables, amplifies, and complicates the
% detection of grey hole attacks. This is a critical omission for an
% SDVN-specific paper. Mobility causes rapid topology changes, intermittent RSU
% handoffs, non-IID traffic patterns across vehicles, and trust score volatility
% — all of which a grey hole attacker can exploit. The paper must explicitly
% address this to distinguish itself from existing VANET/generic SDN work.
% See Part B for the LaTeX insertion.

\section{Proposed System Architecture}

The architecture separates the \emph{detection pipeline} from the \emph{mitigation pipeline} across four functional layers \cite{karof2014}. Detection determines whether a node is malicious; mitigation acts on that verdict. No layer performs both functions simultaneously, and no component is described as doing both without an explicit handoff boundary.

% \begin{figure}[htbp]
% 	\centering
% 	\includegraphics[width=\textwidth,height=0.45\textheight,keepaspectratio]{Figures/figure1_architecture_diagram}
% 	\caption{SHIELD-GH Four-Layer System Architecture}
% 	\label{fig:shield-gh-architecture}
% \end{figure}

% below is without space after figure

% \begin{figure}[H]
% 	\centering
% 	\includegraphics[
% 		width=\textwidth,
% 		height=0.9\textheight,
% 		keepaspectratio
% 	]{Figures/figure1_architecture_diagram}
% 	\caption{SHIELD-GH Four-Layer System Architecture}
% 	\label{fig:shield-gh-architecture}
% \end{figure}

%below is without space before and after for figure
\begin{figure}[!t]
	\centering
	\includegraphics[
		width=\textwidth,
		height=0.7\textheight,
		keepaspectratio
	]{Figures/figure1_architecture_diagram}
	\caption{SHIELD-GH Four-Layer System Architecture}
	\label{fig:shield-gh-architecture}
\end{figure}

Figure~\ref{fig:shield-gh-architecture} presents the four-layer SHIELD-GH system architecture. As shown in Figure~\ref{fig:shield-gh-architecture}, SHIELD-GH is built upon a four-layer architecture, each serving a distinct security role. At the bottom, the Data Plane hosts legitimate, malicious, and emergency vehicles that generate ZKP forwarding proofs and V2V/V2I traffic, which are monitored by RSUs. The Control Plane contains the SDN controller, explicitly designated as untrusted: it may only install flow rules that have been pre-verified by the RSU-maintained Flow Rule Whitelist, preventing a compromised controller from injecting malicious dropping rules undetected. The Blockchain Plane hosts the Dual-Evidence Blockchain Smart Contract (DEBSC) on a Hyperledger Fabric ledger maintained exclusively by RSU nodes via BFT consensus; the SDN controller is excluded from this consensus group and cannot write to the ledger or trigger mitigation independently. The RSU Network Monitor collects traffic metadata and suspicious event logs, committing signed records to the ledger. At the top, the Intelligence Layer's LLM Agent, Fusion Engine, and FL Aggregator generate threat scores that are sent directly to the DEBSC, bypassing the controller. Isolation commands flow from the DEBSC to OpenFlow switches via threshold-signed FlowMods, and group re-keying is performed via $\log N$ Kyber messages broadcast by RSUs.

% --- SUPERVISOR REVISION BLOCK 2: multi-controller, controller trust, RSU bootstrap, whitelist governance, DKG ---
\mypar
\textbf{Multi-Controller Flat Architecture.}
$M$ SDN controllers are registered on the blockchain at bootstrap,
each holding a unique Dilithium key pair and an initial trust score
$\mathcal{T}_c(0)=1$. The active controller for each segment $z$
is the first eligible entry in segment $z$'s blockchain-stored
ordered list (Assumption~A4). Because different segments have
different ordered lists, no single controller accumulates all
routing authority. Controller compromise is therefore not a
single point of failure: segment authority distributes across
the registered fleet by design.

\mypar
\textbf{Controller Trust Module.}
A per-controller trust score $\mathcal{T}_c(t)$ is maintained on
the blockchain by RSU consensus, analogous to the vehicle trust
model. The score decrements on confirmed CP signature triggers
(S4--S6) and on aggregate behavioural anomaly breaches defined in
Section~\ref{sec:signatures}. The RSU Network Monitor serves as
the distributed monitoring entity: each RSU independently
evaluates controller flow rules and submits signed evidence to the
blockchain. No single RSU assessment is sufficient; the trust
score aggregates evidence from all consensus-group RSUs.

\mypar
\textbf{RSU Bootstrap Authentication.}
Each RSU holds a Dilithium certificate issued by the transport
authority at deployment. Participation in the blockchain consensus
group is conditional on presenting a valid, non-revoked certificate.
RSU certificate revocation requires a threshold of other RSUs to
submit signed evidence, preventing any single authority from
silencing a legitimate RSU.

\mypar
\textbf{Flow Rule Whitelist Governance.}
The Flow Rule Whitelist is a blockchain smart contract, not a static
list. Any controller may propose a rule addition, but the rule
becomes active only after $k_{wl}$ independent RSU co-signatures
(Assumption~A5). Proposed-but-unapproved rules remain in a pending
state; OpenFlow switches reject pending rules. This prevents a
compromised controller from whitelisting its own malicious flow
rules to bypass signatures S4--S6.

\mypar
\textbf{Distributed Key Generation for PQC-LKH.}
The root Kyber key pair for the PQC-LKH tree is generated jointly
by the active controller and RSUs via Pedersen Distributed Key
Generation (DKG), as formalised in Section~\ref{sec:dkg}. No
single entity holds complete key material. All re-keying broadcast
messages require a threshold co-signature from the controller and
$k_{\text{key}}-1$ RSUs, so that neither the controller alone nor
any single RSU can unilaterally distribute session keys.
% --- END SUPERVISOR REVISION BLOCK 2 ---

\mypar
\textbf{Data Plane (Vehicular Layer) — Detection role}

\mypar
Vehicles communicate via V2V and V2I links and act as Federated Learning clients. Their sole detection contribution is to generate the raw telemetry from which forwarding statistics are computed: each vehicle records per-slot packet reception counts $n_i^{\text{rx}}(\tau)$ and forwarding counts $n_i^{\text{fwd}}(\tau)$, which are committed to the blockchain as cryptographic forwarding proofs (ZKP, Eq.~\eqref{eq:pedersen}--\eqref{eq:zkp_proof}). The data plane has \emph{no mitigation authority}; isolation commands originate exclusively from the control plane.

\mypar
\textbf{Control Plane (SDN Layer) — Mitigation role}

\mypar
The SDN controller is the sole mitigation authority. Upon receiving a confirmed malicious verdict from the Knowledge Plane, it issues a Dilithium-signed flow modification command (Eq.~\eqref{eq:dilithium_sign}) to the relevant OpenFlow switch to block or reroute traffic around the isolated node. The control plane does \emph{not} perform detection; it only acts on verdicts already produced by the detection pipeline.

\mypar
\textbf{Blockchain Plane (Trust Layer) — Detection and audit role}

\mypar
The blockchain serves two distinct roles that must not be conflated.
As a \textbf{detection baseline}, the immutable ledger accumulates per-vehicle forwarding proofs over all historical RSU interactions $\mathcal{H}_i$; the blockchain reputation score $\mathcal{R}_i(t)$ (Eq.~\eqref{eq:reputation}) is computed from this history and constitutes the statistical gate of the Dual-Evidence Blockchain Smart Contract (DEBSC, Eq.~\eqref{eq:debsc}). A node cannot escape this baseline by short-term honest behaviour because the ledger is append-only.
For \textbf{audit and smart contract enforcement}, smart contracts evaluate the dual-evidence isolation rule: isolation is authorized only when both the statistical reputation gate and the cryptographic ZKP gate fail simultaneously. This is a \emph{detection} decision executed autonomously on-chain, not a mitigation action.

\mypar
\textbf{Knowledge and Security Plane (Intelligence Layer) — Detection role}

\mypar
This plane contains the full-mode detection pipeline. Its components have strictly defined, non-overlapping roles.
The \textbf{LLM Agent} performs semantic threat scoring by processing tokenised blockchain forwarding logs $\mathbf{x}_i^{(t)}$ as a sequence classification task, producing the threat score $\mathcal{Q}_i(t)$ (Eq.~\eqref{eq:llm_score}). The LLM produces a \emph{detection verdict only}; it does not generate routing policies or issue flow rules — routing policy generation is the exclusive responsibility of the Mitigation Engine in the control plane, which receives the LLM verdict as input.
The \textbf{Federated Learning Aggregator} aggregates local vehicle model updates $\Delta\mathbf{w}_i$ via blockchain-verified FedAvg (Eq.~\eqref{eq:fedavg}--\eqref{eq:gradient_verify}) to maintain an up-to-date global detection model without centralising raw traffic data \cite{yurdem2024}.

\section{Threat Model and Assumptions}

\subsection{Attacker Capabilities}

In \textbf{fixed-rate selective dropping}, malicious nodes drop an agreed percentage (e.g., 30--50\%) of packets to cause sustained network throughput disruption.
\textbf{Time-based intermittent attacks} have attackers alternate between normal and malicious states to mimic natural network fluctuations, making statistical detection unreliable.
\textbf{Target-specific intelligent attacks} involve nodes discarding high-priority packets such as emergency broadcasts based on content, source, or destination analysis.
Adversaries may also attempt \textbf{node reputation tampering} by forging forwarding logs to escape detection \cite{sahu2024}.

\subsection{Threat Model Assumptions}

\mypar
These assumptions define the constraints placed upon the adversary to bound the scope of the security framework:

\textbf{T1 (Honest Majority)} assumes that the majority of blockchain consensus nodes are honest, which eliminates 51\% attacks on the chain.
\textbf{T2 (Cryptographic Bound)} presumes the adversary is computationally constrained and cannot forge digital signatures or compromise the blockchain's encryption.
\textbf{T3 (Independent Malice)} assumes that attackers do not cooperate to coordinate their intermittent dropping patterns so as to exactly replicate the traffic of a single honest node.

\mypar
\textbf{System Assumptions}

\textbf{A1 (Immutable Ledger)} assumes the blockchain provides a tamper-proof record of node behaviours.
\textbf{A2 (Privacy Preservation)} presumes Federated Learning is effective at pooling global intelligence without revealing individual vehicle data.
\textbf{A3 (Inference Constraints)} assumes the LLM can perform online inference on blockchain data within the latency and compute constraints of the vehicular environment.

% --- SUPERVISOR REVISION BLOCK 1: extended threat model (T4, T5, A4, A5) ---
\mypar
The following entries extend the threat model to cover coordinated
adversarial behaviour and infrastructure-level trust, which are
necessary for the multi-controller and distributed key architecture
introduced in Section~\ref{sec:methodology}.

\textbf{T4 (Bounded Collusion.)}
The adversary may simultaneously compromise the active SDN
controller and a strict minority of RSUs, provided the number of
compromised RSUs remains below the BFT threshold
$\lfloor(n-1)/3\rfloor$. Coordinated low-and-slow attacks — in
which individually sub-threshold malicious actions by the
controller and minority RSUs combine to degrade the network — are
explicitly within scope. No assumption of independent adversaries
is made between the controller and RSU compromise events.

\textbf{T5 (Infrastructure Bootstrap Authenticity.)}
RSUs are authenticated at system bootstrap via a transport
authority PKI certificate. Physical compromise of an RSU after
bootstrap is detectable through behavioural anomalies in its
signed blockchain submissions, evaluated by the remaining RSU
consensus group.

\textbf{A4 (Zone-Based Controller Assignment.)}
At bootstrap, each geographic network segment $z$ is assigned an
ordered controller list $[c_1^z, c_2^z, \ldots, c_M^z]$
determined by geographic proximity and capacity. This list is
stored immutably on the blockchain and is modifiable only by RSU
consensus. Trust acts as a disqualifier within the list, not as
the ranking criterion.

\textbf{A5 (Whitelist Governance.)}
The Flow Rule Whitelist is a blockchain smart contract requiring
$k_{wl}$ RSU co-signatures to update. No single entity,
including any registered SDN controller, may unilaterally add or
remove whitelist entries.
% --- END SUPERVISOR REVISION BLOCK 1 ---

\clearpage % Force next section to start on a new page


% SUPERVISOR COMMENT — Cryptographic Analysis Section:
Cryptography alone cannot mitigate grey hole attacks.
Reason: A grey hole node receives authenticated packets legitimately, then
silently drops them. Authentication prevents forgery of packets in transit
but does not prevent a legitimate recipient from discarding received packets.
Therefore, cryptography must be used in two distinct phases:
  (1) Pre-detection: as input evidence to improve detection performance —
      specifically by generating unforgeable forwarding proofs (ZKP) that
      make silent dropping detectable via blockchain discrepancy.
  (2) Post-detection: to enforce mitigation — node isolation, group
      re-keying (Kyber), and route blocking (Dilithium-signed FlowMod).

\section{Attack Scenarios in Gray Hole attack}

\mypar
In this project, we consider three types of Grey Hole attacks in Software Defined Vehicular Networks under both controller plane and data plane attack scenarios.

\mypar
\subsection{Data Plane version attack scenarios}

\mypar
\textbf{Fixed Rate Selective Dropping Grey Hole Attack in Data Plane}

\begin{figure}[H]
	\centering
	\includegraphics[width=0.7\textwidth,height=0.45\textheight,keepaspectratio]{d_fixed.png}
	\caption{Fixed-rate selective dropping Grey Hole attack in the data plane}
	\label{fig:data-plane-fixed-rate}
\end{figure}

\mypar
Figure \ref{fig:data-plane-fixed-rate} shows a fixed rate selectively dropping grey hole attack of the data plane. According to that figure there are four vehicles and one RSU in that network. In this network, vehicle 2 act as a attacker.

\mypar
As shown in Figure~\ref{fig:data-plane-fixed-rate}, the attacker begins packet dropping by 50\%: data packets arrive from vehicle 1 to vehicle 2 (attacker), which forwards only the remaining 2 packets to the RSU, which then forwards them to vehicle 3.

\mypar
\textbf{Time Based Intermittent Grey Hole Attack in Data Plane}

\mypar
Figure \ref{fig:data-plane-time-normal} and \ref{fig:data-plane-time-malicious} show the time based intermittent grey hole attack scenario in the data plane at two different times. In this scenario there are three vehicles, one RSU, and an SDN controller.

\begin{figure}[H]
	\centering
	\includegraphics[width=0.8\textwidth,height=0.45\textheight,keepaspectratio]{d_time_normal.png}
	\caption{Normal data-plane behavior at time $t_1$ with no packet dropping}
	\label{fig:data-plane-time-normal}
\end{figure}

\mypar
Figure \ref{fig:data-plane-time-normal} shows the normal behavior of the data plane at time $t = t_1$. All vehicles behave correctly: vehicle V1 sends 4 packets to vehicle 2 (attacker), which forwards all four packets to the RSU, which in turn forwards them to vehicle 3. Communication proceeds normally without any dropping.

\begin{figure}[H]
	\centering
	\includegraphics[width=0.8\textwidth,height=0.45\textheight,keepaspectratio]{d_time_malicous.png}
	\caption{Malicious data-plane behavior at time $t_2$ with intermittent packet dropping by V2}
	\label{fig:data-plane-time-malicious}
\end{figure}

\mypar
As shown in Figure~\ref{fig:data-plane-time-malicious}, at time $t = t_2$, the attacker switches to malicious behaviour: vehicle V1 sends 4 packets to vehicle 2, which drops two packets (50\%) and forwards only the remaining two to the RSU, which forwards them to vehicle 3. The alternation between normal and malicious epochs makes this attack difficult to detect using threshold-based methods.

\mypar
\textbf{Target Specific Intelligent Grey Hole Attack in Data Plane}

\mypar
Figure \ref{fig:d-target} shows the target specific intelligent grey hole attack scenario in the data plane. In this scenario there are four vehicles, one RSU, and a controller. The controller installs flow rules (shown as dashed lines) that vehicle 2 (attacker) selectively enforces.

\mypar
As shown in Figure~\ref{fig:d-target}, vehicle 3 sends three packets to vehicle 2 (attacker), which forwards all three to the RSU without dropping — that path operates correctly according to the flow rule. On a second path, vehicle 1 sends four packets to vehicle 2, which drops 50\% and forwards only two packets to the RSU, which then delivers them to vehicle 4. The attacker thus targets only specific vehicles for dropping, leaving other traffic flows unaffected and making the attack difficult to detect with per-flow aggregates.

\begin{figure}[H]
	\centering
	\includegraphics[width=0.8\textwidth,height=0.45\textheight,keepaspectratio]{d_target.png}
	\caption{Target-specific intelligent Grey Hole attack in the data plane}
	\label{fig:d-target}
\end{figure}

\mypar
\subsection{Controller Plane version attack scenarios}

\mypar
\textbf{Fixed Rate Selective Dropping Grey Hole Attack in Controller Plane}

\mypar
Figure~\ref{fig:controller_fixed_rate} show the  fixed rate selectively dropping  grey hole attack scenario in the controller plane. So, in this controller install malicious flow mod rule in the vehicular network act as a attacker.

\mypar
As shown in Figure~\ref{fig:controller_fixed_rate}, first controller install malicious flow mod rule on RSU and the vehicle 2. So, controller configure on the RSU to drop 50\% of all incoming flows and configures on vehicle 2 to drop 60\% of its received packets. 2nd step show vehicel 1 send 5 packets to the vehicel 2 and 3rd step show, vehicle 2 drop three packets according to the malicious configuration. In 4th step send two packets to the vehicle 3 and vehicle forward all two packets to the RSU as a 5th step. Since RSU already configured malicious flow mod rule RSU drop 50\% packets and forward remain packet to the vehicle 4 as 6th step.Since packets are only partially dropped, the communication appears normal, while the overall network performance is gradually degraded also making the attack difficult to detect in like these dynamic vehicular network.

\begin{figure}[H]
	\centering
	\includegraphics[width=0.8\textwidth,height=0.45\textheight,keepaspectratio]{c_fixed_rate.png}
	\caption{Controller plane fixed rate selective dropping Grey Hole attack}
	\label{fig:controller_fixed_rate}
\end{figure}

\mypar
\textbf{Time Based Intermittent Grey Hole Attack in Controller Plane}

\mypar
Figure  \ref{fig:c_time_normal} and \ref{fig:c_time_malicious} show the time based intermittent gery hole attack scenario in control plane two different times.

\mypar
As shown in Figure~\ref{fig:c_time_normal}, all steps working without packet dropping according to the flow rule. at time $t = t_1$, the SDN controller behaves normally and does not perform any packet dropping. The controller installs correct flow rules on all vehicles and the RSU.

\begin{figure}[H]
	\centering
	\includegraphics[width=0.85\textwidth,height=0.45\textheight,keepaspectratio]{c_time_normal.png}
	\caption{Normal controller behavior at time $t_1$ with no packet dropping.}
	\label{fig:c_time_normal}
\end{figure}

\mypar
As shown in Figure~\ref{fig:c_time_malicious}, at time $t = t_2$, SDN controller install malicious flow mod rule in the network. So, controller configure on the RSU to drop 50\% all flows and on the vehicle 2 to drop 60\% all flows. Then as the 2nd step vehicle 1 send 5 packets to the vehicle 2 and vehicle 2 drop 60\% packets and forward remain two packets to the vehicle 3 as the 4nd step. In 5th step show vehicle 3 forward that two packets to the RSU. According to the malicious flow mode rule on RSU, RSU drop one packet and forward remain packet to the vehicle 4 as the 6th step.This intermittent behavior allows the controller to hide the attack by switching between normal and malicious operation.

\begin{figure}[H]
	\centering
	\includegraphics[width=0.85\textwidth,height=0.45\textheight,keepaspectratio]{c_time_malicous.png}
	\caption{Malicious controller behavior at time $t_2$ with selective packet dropping.}
	\label{fig:c_time_malicious}
\end{figure}

\mypar
\textbf{Target Specific Intelligent Grey Hole Attack in Controller Plane}

\mypar
Figure \ref{fig:controller-target-greyhole} shows the target specific intelligent grey hole attack in controller plane. So, in this first controller install flow rule. So, in this figure only install malicious flow mod rule on RSU. So, RSU drop 50\% packets only if source is vehicle 1. All packets came from other sources, forward without any dropping.

\mypar
As shown in Figure~\ref{fig:controller-target-greyhole}, the first step installs the malicious flow mod rule on RSU. 2nd step shows, vehicle 1 send 4 packets to the RSU and 3 rd step shows, RSU drop two packets since source is vehicle 1. After RSU forward remain 2 packets to the vehicle 4. In step 5 and 6 show the another path packet routing, in that steps there is no any packets dropping since in that path not contain source vehicle 1. In conclusion of this attack scenario, by selectively targeting emergency traffic while forwarding normal traffic, the controller degrades safety critical communication while remaining difficult to detect.

\begin{figure}[H]
	\centering
	\includegraphics[width=0.85\textwidth,height=0.45\textheight,keepaspectratio]{c_target.png}
	\caption{Target-specific intelligent Grey Hole attack in the controller plane}
	\label{fig:controller-target-greyhole}
\end{figure}

\section{Mobility as a Grey Hole Attack Enabler in SDVN}
\label{sec:mobility_analysis}

Existing grey hole detection literature treats the vehicular network as either a
static graph (generic SDN) or a mobile graph with fixed routing protocols
(VANET / AODV). Neither model captures the compound effect of high vehicle mobility
on an SDN-controlled network, where the control plane has centralised visibility
but decentralised physical topology. In SDVN, mobility does not merely add noise
to forwarding statistics — it actively creates four structural conditions that a
grey hole attacker can exploit to evade detection. This section analyses each
condition, quantifies the exploitation mechanism, and identifies the specific
framework component that neutralises it. This analysis is the primary
architectural distinction between the proposed framework and prior VANET or
generic SDN approaches \cite{islam2021sdvn, hsieh2018sdn}.

\subsection{RSU Handoff Windows Create Undetectable Dropping Epochs}

In SDVN, as a vehicle moves between Road Side Unit (RSU) coverage zones, the SDN
controller must tear down flow rules on the departing RSU and re-install them on
the arriving RSU. During this handoff interval $\Delta t_{ho}$, transient packet
loss is \emph{topologically expected} and is indistinguishable from intentional
dropping by a grey hole attacker. A rational attacker exploits this by
synchronising its dropping epochs with handoff events. Formally, if the expected
handoff-induced loss rate at vehicle speed $s_i$ is
$\rho_{ho}(v_i,t) = {s_i(t) \cdot \Delta t_{ho}}/{R_{\text{RSU}}} \cdot
\rho_{\max}$ (Eq.~\eqref{eq:handoff_loss}), the attacker sets its drop rate
$\rho_{att} \leq \rho_{ho}$ during handoff periods, achieving zero detection
probability against any threshold-based detector that does not model handoff loss.

This exploitation is impossible in static SDN (no handoffs) and is less severe
in VANET (no centralised flow rule re-installation, so there is no discrete
handoff signal for the attacker to detect). SDVN uniquely exposes this
vulnerability because the SDN controller broadcasts flow modification messages
whose timing reveals the handoff event to any observing node on the network.

\mypar
\textbf{Framework response:} The Mobility-Aware Trust Decay (MATD) mechanism
(Eq.~\eqref{eq:matd}) corrects the observed PDR by adding back the expected
handoff-induced loss $\rho_{ho}(v_i,t)$ before any signature test is applied
(Eq.~\eqref{eq:mobility_pdr}). The dual-evidence DEBSC gate
(Eq.~\eqref{eq:debsc}) further ensures that a node suffering genuine handoff
loss — whose PDR temporarily drops but whose ZKP forwarding proofs remain valid
for packets it did forward — cannot be falsely isolated.

\subsection{Trust Score Volatility Allows Reputation Reset Through Speed}

Reputation-based and blockchain-based detection schemes accumulate trust scores
over a sequence of forwarding observations. A vehicle that moves through an RSU
coverage zone faster than the trust accumulation horizon $T_{\text{trust}}$
produces too few observations for the score to be statistically significant. An
attacker that intentionally maintains high speed — or that naturally operates in
highway scenarios — can perpetually reset its per-RSU trust score before
incriminating evidence accumulates, while still causing substantial packet loss
across the broader network.

This vulnerability is specific to mobile networks. In static SDN, a node
accumulates observations indefinitely at the same switch. In VANET without a
central reputation store, trust is local and similarly short-lived, but there is
no global baseline to expose the pattern. SDVN with a centralised blockchain
reputation store creates a false sense of security: the ledger is persistent, but
the per-RSU observation window is short, making the long-term reputation score
a lagging indicator that an attacker can game.

\mypar
\textbf{Framework response:} The Bayesian trust model (Eq.~\eqref{eq:trust_update})
uses Beta-distribution priors $(\alpha, \beta)$ to express uncertainty about a
node's behaviour when observations are few. The MATD exponential penalty
(Eq.~\eqref{eq:matd}) explicitly reduces the trust contribution of high-speed
vehicles, so a fast-moving attacker's short observations carry lower weight in
the global reputation score $\mathcal{R}_i(t)$ (Eq.~\eqref{eq:reputation}).
The blockchain aggregates MATD-corrected scores across \emph{all} historical RSU
interactions $\mathcal{H}_i$, so speed-induced resets at individual RSUs do not
erase the global behaviour record.

\subsection{Non-IID Traffic Patterns Corrupt FL-Based Detection Models}

Federated Learning-based intrusion detection assumes that local vehicle models
contribute useful gradient updates toward a globally accurate detection model.
In SDVN, vehicle mobility produces severely non-IID local datasets: a vehicle
on a congested urban street sees dense, high-PDR traffic; a vehicle on a sparse
highway segment sees sparse, high-variance traffic with frequent legitimate loss.
A grey hole attacker that positions itself in a low-density network segment
operates in a region where the global FL model has little local representation,
reducing its detection sensitivity precisely where the attacker is active.

Furthermore, a sophisticated attacker can \emph{exploit} the FL update mechanism
by submitting poisoned gradient updates — gradients computed on manipulated local
data — to degrade the global model's ability to detect its own dropping pattern.
This attack vector is absent in generic SDN (no FL) and underexplored in VANET
FL work, which typically assumes IID traffic \cite{ahsan2024flbert}.

\mypar
\textbf{Framework response:} The blockchain-anchored gradient integrity
mechanism (Eq.~\eqref{eq:gradient_commit}--\eqref{eq:gradient_verify}) requires
each vehicle to commit a hash of its gradient update to the blockchain
\emph{before} transmission. The aggregator rejects any update whose on-chain
hash does not match the received gradient, blocking poisoning without requiring a
trusted aggregator. The weighted FedAvg aggregation (Eq.~\eqref{eq:fedavg})
uses dataset-size weights to reduce the influence of low-observation vehicles in
sparse segments, partially mitigating non-IID skew.

\subsection{Controller Re-Association Enables Stealthy Flow Rule Injection}

SDVN introduces a vulnerability with no equivalent in pure VANET: a vehicle
transitioning between SDN controller domains must re-associate with the new
controller and obtain authenticated flow rules. During this re-association phase,
the vehicle has not yet established a verified session with the new controller
domain. A compromised controller segment can exploit this window to install
malicious flow modification rules that enforce selective dropping before the
vehicle's integrity-checking mechanisms are active.

This attack is controller-plane specific and cannot occur in VANET, which has no
centralised controller. It is also distinct from data-plane grey hole attacks
because the malicious behaviour is encoded in flow rules, not in the vehicle's own
forwarding logic, making it invisible to PDR-based detectors that only observe
forwarding statistics.

\mypar
\textbf{Framework response:} All SDN controller flow modification commands are
authenticated with CRYSTALS-Dilithium post-quantum signatures
(Eq.~\eqref{eq:dilithium_sign}--\eqref{eq:dilithium_verify}). A vehicle receiving
a flow rule during re-association verifies the Dilithium signature against the
controller's public key before applying the rule. The controller-plane attack
signatures S4--S6 (Eqs.~\eqref{eq:sig_cpfr}--\eqref{eq:sig_cpts}) detect
anomalous drop rules directly from the installed flow rule set, independent of
forwarding statistics, catching injected rules that would be invisible to
data-plane-only detectors.

\section{Detection Methodology}
\label{sec:detection_methodology}

This section establishes the mathematical foundations for grey hole detection,
modelling packet forwarding behaviour, correcting for mobility-induced loss,
and defining the attack signature conditions used by the detection engine.

\subsection{Packet Forwarding Behaviour Modelling}

To establish a quantitative basis for grey hole detection, we first model the
forwarding behaviour of each vehicle over a sliding observation window. The
The \emph{observed packet delivery ratio} of vehicle $v_i$ over a window of $W$
time slots ending at time $t$ is defined in Equation~\eqref{eq:pdr}.

\begin{equation}
\label{eq:pdr}
\text{PDR}_i(t, W) = \frac{\displaystyle\sum_{\tau=t-W}^{t} n_i^{\text{fwd}}(\tau)}
                           {\displaystyle\sum_{\tau=t-W}^{t} n_i^{\text{rx}}(\tau)}
\end{equation}

As shown in Eq.~\eqref{eq:pdr}, the packet delivery ratio $\text{PDR}_i(t, W)$ is the ratio of the total packets forwarded $\sum_\tau n_i^{\text{fwd}}(\tau)$ to the total packets received $\sum_\tau n_i^{\text{rx}}(\tau)$ by vehicle $v_i$ over a sliding window of $W$ time slots ending at time $t$. Here, $n_i^{\text{fwd}}(\tau)$ and $n_i^{\text{rx}}(\tau)$ are the per-slot forwarded and received packet counts respectively. This ratio is the primary forwarding behaviour metric evaluated by all six detection signatures S1--S6.

The \emph{instantaneous drop rate} at a single slot is defined in Equation~\eqref{eq:droprate}.

\begin{equation}
\label{eq:droprate}
\delta_i(t) = 1 - \text{PDR}_i(t, 1)
\end{equation}

As shown in Eq.~\eqref{eq:droprate}, and building on Eq.~\eqref{eq:pdr}, the instantaneous drop rate $\delta_i(t)$ is defined as one minus the single-slot PDR of vehicle $v_i$ at time $t$. A value of $\delta_i(t) = 0$ indicates all received packets were forwarded, while $\delta_i(t) = 1$ indicates all were dropped. This per-slot measure is used in the autocorrelation test of Signature S2 to detect periodic on/off dropping patterns.

Because grey hole attackers exploit the similarity between their behaviour and
natural channel loss, we also track the \emph{variance of PDR} over the window
to separate consistently low forwarding (fixed-rate attack) from fluctuating
forwarding (intermittent attack or legitimate loss), as defined in Equation~\eqref{eq:pdr_var}.

\begin{equation}
\label{eq:pdr_var}
\sigma^2_i(W) = \frac{1}{W} \sum_{\tau=t-W}^{t}
\!\left(\text{PDR}_i(\tau,1) - \overline{\text{PDR}}_i(W)\right)^2
\end{equation}

As shown in Eq.~\eqref{eq:pdr_var}, the PDR variance $\sigma_i^2(W)$ is the mean squared deviation of the per-slot PDR from the window mean $\overline{\text{PDR}_i}(W)$ over window $W$. Here, $\text{PDR}_i(\tau, 1)$ is the single-slot PDR at slot $\tau$ and $\overline{\text{PDR}_i}(W)$ is the mean PDR over the full window. Low variance combined with a low mean PDR triggers Signature S1 (fixed-rate attack), while high variance suggests either intermittent dropping or legitimate channel loss.

% \noindent where $\overline{\text{PDR}}_i(W)$ is the mean PDR over the window.

\subsection{Mobility-Aware Packet Loss Correction}

In SDVN, RSU handoff events cause transient packet loss that is topologically
expected and must be separated from malicious dropping before any signature
test is applied. We model the \emph{expected handoff-induced loss rate}
$\rho_{ho}(v_i, t)$ as a linear function of vehicle speed $s_i(t)$ normalised
by the RSU coverage radius $R_{\text{RSU}}$, as defined in Equation~\eqref{eq:handoff_loss}.

\begin{equation}
\label{eq:handoff_loss}
\rho_{ho}(v_i, t) = \frac{s_i(t) \cdot \Delta t_{ho}}{R_{\text{RSU}}}
\cdot \rho_{\max}
\end{equation}

As shown in Eq.~\eqref{eq:handoff_loss}, the expected handoff-induced loss rate $\rho_{\text{ho}}(v_i, t)$ is modelled as a linear function of vehicle speed $s_i(t)$. Here, $\Delta t_{\text{ho}}$ is the average duration of an RSU handoff transition, $R_{\text{RSU}}$ is the RSU coverage radius, and $\rho_{\text{max}}$ is the worst-case handoff loss rate calibrated from simulation. The ratio $s_i(t) \cdot \Delta t_{\text{ho}} / R_{\text{RSU}}$ approximates the fraction of time a vehicle spends in a handoff transition, and multiplying by $\rho_{\text{max}}$ yields the expected loss contribution.

The \emph{mobility-corrected PDR} adds this expected loss back so that the
signature engine operates on behaviour net of topology effects, as defined in Equation~\eqref{eq:mobility_pdr}.

\begin{equation}
\label{eq:mobility_pdr}
\widehat{\text{PDR}}_i(t, W) = \text{PDR}_i(t, W) + \rho_{ho}(v_i, t)
\end{equation}

As shown in Eq.~\eqref{eq:mobility_pdr}, the mobility-corrected PDR $\widehat{\text{PDR}}_i(t, W)$ is obtained by adding the expected handoff-induced loss rate $\rho_{\text{ho}}(v_i, t)$ from Eq.~\eqref{eq:handoff_loss} back to the observed PDR from Eq.~\eqref{eq:pdr}. This correction ensures that all six attack signatures S1--S6 are evaluated on forwarding behaviour net of topology-induced effects. Without this correction, a legitimate high-speed vehicle undergoing RSU handoff would exhibit a PDR drop that could incorrectly trigger Signature S1.

\subsection{Attack Signature Conditions}
\label{sec:signatures}

The six formal attack signatures are stated as binary indicator functions
over the mobility-corrected statistics derived above. Each signature is
evaluated independently by the lightweight rule-based detector.

\textbf{Signature S1 (DP-FR).} A node exhibiting a consistently low
mobility-corrected PDR with small variance satisfies the fixed-rate data-plane
signature, since a deliberate constant drop rate produces a stable
low-forwarding region is shown in Equation~\eqref{eq:sig_dpfr}.

\begin{equation}
\label{eq:sig_dpfr}
\mathcal{S}_{\text{DP-FR}}(v_i) =
\mathbf{1}\!\left[\,
\widehat{\text{PDR}}_i(t, W) < \tau_f \;\wedge\;
\sigma^2_i(W) < \varepsilon_f
\,\right]
\end{equation}

As shown in Eq.~\eqref{eq:sig_dpfr}, $\tau_f$ is the minimum 
acceptable forwarding ratio and $\varepsilon_f$ is the maximum 
allowable variance for a steady-state drop pattern. The combination 
of low mean and low variance distinguishes a deliberate constant-rate 
dropper from a vehicle experiencing variable legitimate loss.

\textbf{Signature S2 (DP-IT).} An intermittent attacker alternates between
normal and malicious periods, producing a periodic binary pattern in the
instantaneous drop signal. We detect this via the empirical autocorrelation
$R_{m_i}(T) = \frac{1}{W}\sum_\tau m_i(\tau)\,m_i(\tau{-}T)$ of the binary
malicious indicator $m_i(\tau)=\mathbf{1}[\text{PDR}_i(\tau,1)<\tau_{it}]$ is shown in Equation~\eqref{eq:sig_dpit}.

\begin{equation}
\label{eq:sig_dpit}
\mathcal{S}_{\text{DP-IT}}(v_i) =
\mathbf{1}\!\left[\,
\exists\,T^* \in [T_{\min}, T_{\max}] : R_{m_i}(T^*) > \gamma_{it}
\,\right]
\end{equation}

As shown in Eq.~\eqref{eq:sig_dpit}, $\tau_{\text{it}}$ is the 
per-slot drop threshold, $T^*$ is the candidate period, 
$[T_{\min}, T_{\max}]$ is the expected attack cycle range, and 
$\gamma_{\text{it}}$ is the autocorrelation significance threshold. 
The signature fires when a statistically significant periodic 
component exists, identifying an attacker whose malicious epochs 
repeat at a fixed cycle.

\textbf{Signature S3 (DP-TS).} A target-specific attacker drops packets
selectively by source, so its per-source PDR distribution deviates
significantly from the uniform baseline expected of a benign forwarder.
We quantify this deviation using the KL-divergence between the empirical
per-source PDR distribution $P_{\text{PDR}_i}^{(s)}$ and the uniform
reference $\mathcal{U}$ is shown in Equation~\eqref{eq:sig_dpts}.

\begin{equation}
\label{eq:sig_dpts}
\mathcal{S}_{\text{DP-TS}}(v_i) =
\mathbf{1}\!\left[\,
D_{\mathrm{KL}}\!\left(P_{\text{PDR}_i}^{(s)} \;\|\; \mathcal{U}\right)
> \tau_{ts}
\,\right]
\end{equation}

As shown in Eq.~\eqref{eq:sig_dpts}, $P_{\text{PDR}_i}^{(s)}$ 
captures how forwarding rates vary across different traffic sources, 
$\mathcal{U}$ is the uniform reference distribution, and $\tau_{ts}$ 
is the divergence threshold. A high KL-divergence indicates that 
vehicle $v_i$ forwards traffic from some sources normally while 
dropping from others, which is the defining characteristic of a 
target-specific grey hole attacker.

\textbf{Signature S4 (CP-FR).} In the controller plane, a malicious controller
injects flow rules that unconditionally drop a fixed fraction of all flows.
The signature fires when any installed flow rule carries a drop action with
probability exceeding threshold $\tau_c$ is shown in Equation~\eqref{eq:sig_cpfr}.

\begin{equation}
\label{eq:sig_cpfr}
\mathcal{S}_{\text{CP-FR}}(c) =
\mathbf{1}\!\left[\,
\exists\,f \in \mathcal{F}_c(t) :
\mathrm{action}(f) = \texttt{drop} \;\wedge\;
p_{\text{drop}}(f) > \tau_c
\,\right]
\end{equation}

As shown in Eq.~\eqref{eq:sig_cpfr}, $\mathcal{F}_c(t)$ is the set 
of all flow rules currently installed by controller $c$, 
$\text{action}(f)$ is the action field of rule $f$, and 
$p_{\text{drop}}(f)$ is the associated drop probability exceeding 
threshold $\tau_c$. Unlike data-plane signatures, this check operates 
directly on the flow rule set and does not require historical PDR data.

\textbf{Signature S5 (CP-IT).} An intermittent controller attacker alternates
between benign and malicious flow-rule epochs. Using the same autocorrelation
test applied to the time series of malicious flow-rule counts
$\mathcal{F}_c^{\text{mal}}(t)$, the signature is is shown in Equation~\eqref{eq:sig_cpit}.

\begin{equation}
\label{eq:sig_cpit}
\mathcal{S}_{\text{CP-IT}}(c) =
\mathbf{1}\!\left[\,
R_{\mathcal{F}_c^{\text{mal}}}(T^*) > \gamma_c \;\wedge\;
\overline{\mathcal{F}_c^{\text{mal}}}(W) > 0
\,\right]
\end{equation}

As shown in Eq.~\eqref{eq:sig_cpit}, $R_{\mathcal{F}_c^{\text{mal}}}(T^*)$ 
is the autocorrelation of malicious flow-rule counts at candidate 
period $T^*$, $\gamma_c$ is the significance threshold, and 
$\overline{\mathcal{F}_c^{\text{mal}}}(W) > 0$ confirms at least one 
malicious rule was observed in the window. The signature detects a 
compromised controller alternating between benign and malicious flow 
rules at a fixed period.

% --- SUPERVISOR REVISION BLOCK 3: revised S6 with whitelist + safety-class gating ---
\textbf{Signature S6 (CP-TS).}
A target-specific controller attack installs drop rules conditioned
on specific match fields. However, standard SDN quality-of-service
(QoS) policies also use non-wildcard drop rules to deprioritise
low-priority traffic on specific match fields. A signature that
fires on any non-wildcard drop rule would produce continuous false
positives in production networks. S6 is therefore tightened to
fire only when three joint conditions hold: the rule has a
non-wildcard drop action, the rule is absent from the blockchain
Flow Rule Whitelist $\mathcal{W}_{\text{BC}}$, and the targeted
traffic class belongs to the safety-critical class set
$\mathcal{C}_{\text{safety}}$. The revised signature is given in
Equation~\eqref{eq:sig_cpts}.

\begin{equation}
\label{eq:sig_cpts}
\mathcal{S}_{\text{CP-TS}}(c) =
\mathbf{1}\!\left[\,
\exists\,f \in \mathcal{F}_c(t) :
\;\mathrm{action}(f) = \texttt{drop}
\;\wedge\;
\mathrm{match}(f) \neq \texttt{WILDCARD}
\;\wedge\;
f \notin \mathcal{W}_{\text{BC}}
\;\wedge\;
\mathrm{class}(f) \in \mathcal{C}_{\text{safety}}
\,\right]
\end{equation}

As given in Equation~\eqref{eq:sig_cpts}, $\mathcal{F}_c(t)$ is
the active flow rule set of controller $c$, $\mathcal{W}_{\text{BC}}$
is the RSU-governed blockchain whitelist of approved rules, and
$\mathcal{C}_{\text{safety}}$ is the set of safety-critical traffic
classes (emergency broadcasts, collision warnings, intersection
management signals). The whitelist cross-reference ensures that
legitimate QoS drop rules approved at bootstrap by RSU consensus
and entered into $\mathcal{W}_{\text{BC}}$ do not trigger S6.
Only unapproved rules targeting safety-critical traffic are flagged,
eliminating false positives from routine traffic management.
% --- END SUPERVISOR REVISION BLOCK 3 ---

% --- SUPERVISOR REVISION BLOCK 4: aggregate controller behavioural anomaly score ---
\mypar
\textbf{Aggregate Controller Behavioural Anomaly Score.}
Signatures S4--S6 operate at the individual rule level. A
sophisticated controller can evade them by installing many rules
each with a drop probability just below the hard threshold
$\tau_c$ — none of which individually triggers S4, but whose
combined effect substantially degrades the network. To detect
this low-and-slow evasion, the aggregate controller behavioural
anomaly score $\Psi_c(t)$ counts sub-threshold suspicious rules
within a rolling window. The aggregate score is given in
Equation~\eqref{eq:controller_agg}.

\begin{equation}
\label{eq:controller_agg}
\Psi_c(t) =
\sum_{\tau=t-W_c}^{t}
\sum_{f \in \mathcal{F}_c(\tau)}
\mathbf{1}\!\left[\,
p_{\text{drop}}(f) \in (\tau_{\text{sub}},\, \tau_c)
\,\right]
\end{equation}

As given in Equation~\eqref{eq:controller_agg}, $W_c$ is the
controller observation window, $\tau_{\text{sub}}$ is the lower
bound of the suspicious sub-threshold range, and $\tau_c$ is
the S4 hard threshold. A rule with $p_{\text{drop}}(f) \in
(\tau_{\text{sub}}, \tau_c)$ is individually below the detection
threshold but collectively indicative of coordinated degradation.
When $\Psi_c(t)$ exceeds $\psi_{\text{thresh}}$, the controller
trust score is decremented even if no individual rule triggered
S4--S6.

\mypar
The per-controller trust score integrating both hard-threshold and
aggregate-anomaly evidence is given in
Equation~\eqref{eq:controller_trust}.

\begin{equation}
\label{eq:controller_trust}
\mathcal{T}_c(t) = \mathcal{T}_c(t-1)
\;-\;
\delta_{\text{sig}}\cdot
\mathbf{1}\!\left[\,\bigvee_{k=4}^{6}
\mathcal{S}_k(c,t)\,\right]
\;-\;
\delta_{\text{agg}}\cdot
\mathbf{1}\!\left[\,\Psi_c(t)
> \psi_{\text{thresh}}\,\right]
\end{equation}

As given in Equation~\eqref{eq:controller_trust},
$\delta_{\text{sig}}$ is the decrement weight for a confirmed
S4--S6 trigger, $\delta_{\text{agg}}$ is the decrement weight for
an aggregate anomaly breach, and $\bigvee_{k=4}^{6}\mathcal{S}_k
(c,t)$ is the logical OR of the three CP signatures. The score
is initialised at $\mathcal{T}_c(0)=1$, maintained on the
blockchain by RSU consensus, and is the basis for the controller
failover condition in Section~\ref{sec:failover}.
% --- END SUPERVISOR REVISION BLOCK 4 ---

\section{Detection and Mitigation Framework}
\label{sec:design}

The proposed framework goes beyond applying existing technologies by
introducing three design-level innovations that address SDVN-specific
constraints not resolvable by straightforward technology deployment.





\subsection{Dual-Mode Detection Framework}

The framework operates in two modes that trade detection latency for detection depth.
In \textbf{Lightweight mode}, the rule-based detector evaluates the six formal attack signatures S1--S6 (Eqs.~\eqref{eq:sig_dpfr}--\eqref{eq:sig_cpts}) directly against mobility-corrected PDR statistics from the blockchain log. No LLM inference is required. This mode is the primary detector for high-confidence, low-ambiguity attack patterns such as fixed-rate dropping and clear target-specific rules, and it operates within the latency budget of real-time vehicular environments.
In \textbf{Full mode}, when no lightweight signature fires or when the edge LLM confidence falls below threshold $\varepsilon_u$ (Eq.~\eqref{eq:llm_routing}), the full-mode pipeline engages the LLM semantic scorer and the FL-updated global model. The final verdict fuses all three evidence sources — rule-based, LLM, and blockchain reputation — via the learned linear combination in Eq.~\eqref{eq:fusion}.

Both modes feed into the same DEBSC isolation gate (Eq.~\eqref{eq:debsc}), ensuring that neither mode can trigger isolation on statistical evidence alone without cryptographic confirmation.

% \begin{figure}[htbp]
% 	\centering
% 	\includegraphics[width=\textwidth,height=0.45\textheight,keepaspectratio]{Figures/figure2_dual_mode_architecture}
% 	\caption{SHIELD-GH Dual-Mode Detection Framework}
% 	\label{fig:dual-mode-architecture}
% \end{figure}

\begin{figure}[H]
	\centering
	\includegraphics[
		width=\textwidth,
		height=0.7\textheight,
		keepaspectratio
	]{Figures/figure2_dual_mode_architecture}
	\caption{SHIELD-GH Dual-Mode Detection Framework}
	\label{fig:dual-mode-architecture}
\end{figure}

Figure~\ref{fig:dual-mode-architecture} illustrates the SHIELD-GH dual-mode detection framework. As shown in Figure~\ref{fig:dual-mode-architecture}, SHIELD-GH employs two detection pipelines operating in parallel, as described below. Both pipelines share the Mobility-Aware Trust Decay (MATD) module at the top, which supplies mobility-corrected PDR values to eliminate false positives caused by RSU handoff packet loss. In Lightweight Mode, the OpenFlow Counter Collector gathers per-vehicle packet statistics per time window, which are passed to the Attack Signature Engine that evaluates the six formal signatures (S1: Fixed-Rate DP, S2: Intermittent DP, S3: Target-Specific DP, S4: Fixed-Rate CP, S5: Intermittent CP, S6: Target-Specific CP). If a signature fires, Lightweight Mitigation is triggered using HMAC Authentication combined with an RSU Threshold-signed FlowMod, all verified against the blockchain baseline. In Full Mode, the Blockchain Log Tokenizer converts forwarding records into sequences fed simultaneously to the LLM Agent for semantic threat scoring and to the FL Global Model for pattern classification. The Fusion Engine combines the LLM score, FL score, and Blockchain Reputation Deficit into a single combined detection decision. If the decision is positive, the PQC Mitigation Engine performs ZKP verification, DEBSC-triggered isolation, PQC-LKH group re-keying, and Dilithium FlowMod installation. Both modes read from and write to the shared Blockchain Trust Layer maintained by RSUs.

















\subsection{Dual-Evidence Blockchain Smart Contract (DEBSC)}

Naive application of blockchain to vehicular security simply logs node
behaviours as immutable records and relies on a single reputation
threshold to trigger isolation. This design fails in SDVN for two
reasons: (i) mobility causes legitimate PDR volatility that crosses
reputation thresholds, producing false isolations; and (ii) a
sophisticated attacker that is aware of the threshold can calibrate
drop rates just below it while still maintaining effective attack
throughput.

The DEBSC resolves both issues through a dual-evidence gate: isolation
is only authorized when a statistical anomaly (low MATD-corrected
reputation score) AND a cryptographic failure (invalid or absent ZKP
forwarding proof) are detected simultaneously, as formalised in
Eq.~\eqref{eq:debsc}.

The design insight is that a legitimate vehicle with genuine packet loss
(e.g., during RSU handoff) will have a temporarily low reputation but
will still produce valid ZKP forwarding proofs for the packets it did
forward, so only the statistical gate fires and isolation is blocked.
A grey hole attacker that drops packets cannot produce valid proofs for
the dropped packets, so both gates fire and isolation proceeds. This
dual-gating design reduces false isolations to near zero while maintaining
detection sensitivity, a property not achievable by either mechanism alone.

The DEBSC also implements a \emph{graduated response} rather than a
binary isolate/accept decision. Define the \emph{suspicion level}
$\Lambda_i(t)$ as the number of consecutive rounds in which the statistical
gate has fired is shown in Equation~\eqref{eq:suspicion}.

\begin{equation}
\label{eq:suspicion}
\Lambda_i(t) = \sum_{\tau=t-W_s}^{t}
\mathbf{1}\bigl[(1-\mathcal{R}_i(\tau)) > \theta_R\bigr]
\end{equation}

As shown in Eq.~\eqref{eq:suspicion}, $W_s$ is the sliding window 
length, $\theta_R$ is the reputation threshold, and 
$\mathcal{R}_i(\tau)$ is the reputation score at round $\tau$. A 
higher $\Lambda_i(t)$ triggers progressively stricter graduated 
responses from the smart contract.



The smart contract applies three graduated responses, (1) for
$\Lambda_i < \lambda_1$, rate-limit traffic from $v_i$; (2) for
$\lambda_1 \leq \Lambda_i < \lambda_2$, require ZKP proof verification
for every packet batch, (3) for $\Lambda_i \geq \lambda_2$ and
$\Pi_{\mathrm{ZKP}}(v_i) \in \{\texttt{FAIL},\,\texttt{ABSENT}\}$,
trigger full isolation.
This design eliminates the abrupt false-isolation problem of
threshold-based single-evidence systems.

% --- SUPERVISOR REVISION BLOCK 7: controller failover protocol subsection ---
\subsection{Controller Failover Protocol}
\label{sec:failover}

When the controller trust score $\mathcal{T}_c(t)$ falls below
threshold $\theta_c$, the DEBSC triggers a controller switch event
for every segment $z$ where $c$ is currently active. Selecting the
globally highest-trust controller at every failover event would
cause one controller to accumulate all fallback sessions,
reintroducing a dominant single point of failure. Instead, failover
selects the first eligible controller from the segment-specific
ordered list stored on the blockchain (Assumption~A4), using trust
only as a disqualifier. The failover selection condition is given
in Equation~\eqref{eq:failover_select}.

\begin{equation}
\label{eq:failover_select}
c_{\text{new}}^z =
\underset{j}{\arg\min}\; j
\quad \text{s.t.} \quad
c_j^z \neq c_{\text{failed}},\quad
\mathcal{T}_{c_j^z}(t) > \theta_c
\end{equation}

As given in Equation~\eqref{eq:failover_select},
$[c_1^z, c_2^z, \ldots, c_M^z]$ is the blockchain-stored ordered
list for segment $z$, $c_{\text{failed}}$ is the controller being
replaced, and $\theta_c$ is the minimum eligibility threshold.
The selected replacement is the lowest-ranked eligible entry, not
the globally highest-trust controller. Because different segments
carry different ordered lists established at bootstrap by
geographic proximity and capacity, failover load distributes
naturally across the fleet.

Upon selection, the DEBSC executes three actions atomically:
(1)~the failed controller's Dilithium public key is appended to the
blockchain revocation list via RSU consensus; (2)~all OpenFlow
switches in segment $z$ are notified to reject signatures from
the revoked key; (3)~the new controller re-establishes flow rules
for segment $z$ exclusively through the whitelist-governed
installation procedure, ensuring no unverified rules take effect
during the transition.
% --- END SUPERVISOR REVISION BLOCK 7 ---

\subsection{Blockchain-Anchored FL Gradient Integrity Verification}

Standard federated learning is vulnerable to model poisoning: a malicious
vehicle submits corrupted gradient updates to degrade the global detection
model. Existing defences (e.g., robust aggregation, anomaly filtering on
gradient norms) are computationally intensive and require a trusted
aggregator which does not exist in a decentralised SDVN.

The proposed design anchors each gradient update to the blockchain
\emph{before} the aggregation round. Each vehicle $v_i$ submits a
commitment $\mathcal{C}_i = \texttt{Hash}(\Delta\mathbf{w}_i \| t \| \mathrm{id}_i)$
to the blockchain at the start of the round, and transmits
$\Delta\mathbf{w}_i$ to the aggregator. The aggregator verifies is shown in Equation~\eqref{eq:gradient_commit}.

\begin{equation}
\label{eq:gradient_commit}
\texttt{Valid}_i = \mathbf{1}\bigl[
\texttt{Hash}(\Delta\mathbf{w}_i \| t \| \mathrm{id}_i) =
\mathcal{C}_i^{\mathrm{BC}}
\bigr]
\end{equation}

As shown in Eq.~\eqref{eq:gradient_commit}, a gradient update from vehicle $v_i$ is accepted only if the hash of the received gradient $\Delta w_i$ matches the pre-committed hash $C_i^{\text{BC}}$ stored on the blockchain. Here, $\Delta w_i$ is the local model update, $t$ is the round timestamp, and $id_i$ is the vehicle identifier. Any tampered or poisoned gradient will produce a hash mismatch and is rejected before aggregation.

\noindent where $\mathcal{C}_i^{\mathrm{BC}}$ is the commitment retrieved
from the blockchain. Only verified updates enter the FedAvg aggregation
(Eq.~\eqref{eq:fedavg}).

\subsection{Edge-LLM Architecture for Real-Time Grey Hole Detection}

Deploying a full LLM on resource-constrained vehicle OBUs is infeasible.
The naive solution deploying the LLM only on the cloud introduces
unacceptable latency for real-time safety-critical decisions. Motivated
by the edge-cloud design of MistralBSM~\cite{hamhoum2024mistralbsm},
the proposed framework employs a two-tier LLM deployment. They are,

\textbf{Tier 1 (Edge LLM at RSU):} A 4-bit NF4 quantised,
LoRA-adapted Qwen2.5-7B-Instruct model (7.6\,B parameters,
LoRA rank $r_{\text{edge}}=8$) processes short blockchain log
windows $\mathbf{x}_i^{(t)}$ of $n_e = 96$ tokens at the RSU
in real time. Each token encodes one forwarding event per slot
(\texttt{FWD}/\texttt{DRP}/\texttt{HOF}); controller-plane
variants prepend a \texttt{RULE} prefix token. The model
produces a seven-class probability vector (BENIGN $+$ S1--S6)
from which the binary preliminary verdict
$\hat{y}_i^{\mathrm{edge}} \in \{0,1\}$ is derived.

\textbf{Tier 2 (Cloud LLM):} For cases where the edge LLM
confidence falls below threshold $\varepsilon_u = 0.15$, the
same Qwen2.5-7B-Instruct base model operating under a
higher-rank LoRA configuration ($r_{\text{cloud}}=16$) processes
a longer historical window $n_c \gg n_e$ for a definitive
verdict, providing greater representational capacity without
maintaining a separate model family.

The routing decision is shown in Equation~\eqref{eq:llm_routing}.

\begin{equation}
\label{eq:llm_routing}
\text{Use Tier 2} =
\mathbf{1}\bigl[
\max_c \,\mathrm{softmax}(\mathrm{LLM}_{\mathrm{edge}}(\mathbf{x}_i))_c < \varepsilon_u
\bigr]
\end{equation}

As shown in Eq.~\eqref{eq:llm_routing}, the routing condition for the two-tier LLM is defined such that if the maximum softmax confidence of the edge LLM over all classes $c$ falls below uncertainty threshold $\varepsilon_u$, the input $x_i$ is escalated to the cloud-tier LLM for deeper analysis. Here, $\text{LLM}_{\text{edge}}(x_i)$ is the edge model output logit vector and $\varepsilon_u$ is the confidence threshold tuned on the validation set.



The final threat score $\mathcal{Q}_i(t)$ (Eq.~\eqref{eq:llm_score}) is
taken from whichever tier produces the verdict. This design reduces average
inference latency to $\bar{L}_{\mathrm{edge}}$ ms for the majority of cases
(high-confidence edge decisions), while reserving the full cloud LLM for
ambiguous borderline cases a latency-accuracy trade-off not achieved by
either single-tier deployment. The window size $n_e$ is a design parameter
tuned jointly with $\varepsilon_u$ on the validation split to balance
latency and accuracy.

\subsection{Mobility-Aware Trust Scoring and Blockchain Reputation}

Trust accumulation in high-mobility environments must account for the fact that
short RSU dwell times yield statistically insufficient observations. We adopt a
Bayesian Beta-distribution trust model in which forwarding successes and failures
update the posterior mean, giving the \emph{instantaneous trust score} is shown in Equation~\eqref{eq:trust_update}.

\begin{equation}
\label{eq:trust_update}
\mathcal{T}_i(t) = \frac{\alpha + n_i^{\text{fwd}}(t)}
                         {\alpha + n_i^{\text{fwd}}(t)
                          + \beta + n_i^{\text{drop}}(t)}
\end{equation}

As shown in Eq.~\eqref{eq:trust_update}, $\alpha$ and $\beta$ are 
Beta prior parameters representing initial trust assumptions, 
$n_i^{\text{fwd}}(t)$ is the number of packets forwarded, and 
$n_i^{\text{drop}}(t) = n_i^{\text{rx}}(t) - n_i^{\text{fwd}}(t)$ 
is the number of packets dropped.

A high-speed vehicle accumulates few observations per RSU segment, so blindly
trusting this score risks false isolation. The \emph{Mobility-Aware Trust
Decay (MATD)} applies an exponential penalty proportional to current speed,
reducing the effective trust contribution of high-speed vehicles for which
the observation basis is shallow is shown in Equation~\eqref{eq:matd}.

\begin{equation}
\label{eq:matd}
\mathcal{T}_i^{\text{mob}}(t) = \mathcal{T}_i(t) \cdot
e^{-\lambda_s\, s_i(t)\, \Delta t}
\end{equation}

As shown in Eq.~\eqref{eq:matd}, $\lambda_s > 0$ is the mobility 
decay coefficient calibrated to the RSU coverage radius and vehicle 
density, $s_i(t)$ is the current speed of vehicle $v_i$, and 
$\Delta t$ is the observation slot duration.

 The \emph{blockchain reputation
score} then aggregates MATD-corrected trust over all historical RSU interactions
$\mathcal{H}_i$ stored in the immutable ledger, providing a long-term behavioural
baseline immune to short-term manipulation is shown in Equation~\eqref{eq:reputation}.

\begin{equation}
\label{eq:reputation}
\mathcal{R}_i(t) = \frac{1}{|\mathcal{H}_i|}
\sum_{h \in \mathcal{H}_i} \mathcal{T}_i^{\text{mob}}(h)
\end{equation}

As shown in Eq.~\eqref{eq:reputation}, $\mathcal{H}_i$ is the set 
of historical RSU interactions stored on the immutable blockchain 
ledger and $|\mathcal{H}_i|$ is the total number of such interactions. 
Since the ledger is append-only, $\mathcal{R}_i(t)$ cannot be reset 
by short-term honest behaviour.

% --- SUPERVISOR REVISION BLOCK 5: three-state ZKP, RSU cross-reference, revised DEBSC ---
To prevent a single corrupted evidence source from triggering false
isolation, the \emph{Dual-Evidence Blockchain Smart Contract (DEBSC)}
requires both the statistical reputation score and the cryptographic
forwarding proof to fail simultaneously before an isolation command
is authorised.

\mypar
\textbf{ZKP Evidence State.}
The original two-state ZKP model (PASS or FAIL) contains a critical
gap: a malicious vehicle can simply refuse to submit a forwarding
proof, or submit it after the observation window closes. Under the
two-state model, an absent proof is treated as neither PASS nor
FAIL, so the cryptographic gate never fires. The complete three-state
ZKP evidence model is given in Equation~\eqref{eq:zkp_state}.

\begin{equation}
\label{eq:zkp_state}
\Pi_{\text{ZKP}}(v_i, t) =
\begin{cases}
\texttt{PASS}
  & \text{if valid } \pi_i \text{ received within } T_{\text{zkp}}
    \text{ and RSU cross-reference consistent} \\[4pt]
\texttt{FAIL}
  & \text{if } \pi_i \text{ received but verification fails,}
    \text{ or cross-reference inconsistent} \\[4pt]
\texttt{ABSENT}
  & \text{if no } \pi_i \text{ received within } T_{\text{zkp}}
\end{cases}
\end{equation}

As given in Equation~\eqref{eq:zkp_state}, $T_{\text{zkp}}$ is
the proof submission deadline per observation window. The ABSENT
state closes the self-reporting evasion: a vehicle that withholds
its proof is treated identically to one that submits a verifiably
false proof.

\mypar
\textbf{RSU-Observed Cross-Reference.}
Forwarding counts in the vehicle's Pedersen commitment
(Equation~\eqref{eq:pedersen}) are independently corroborated
against RSU observations at the coverage boundary. The consistency
condition is given in Equation~\eqref{eq:rsu_crossref}.

\begin{equation}
\label{eq:rsu_crossref}
\mathrm{Consistent}(v_i, t) =
\mathbf{1}\!\left[\,
\left|
n_i^{\text{fwd}} - \widehat{n}_i^{\text{fwd}}(t)
\right| \leq \varepsilon_{\text{obs}}
\,\right]
\end{equation}

As given in Equation~\eqref{eq:rsu_crossref}, $n_i^{\text{fwd}}$
is the forwarded count in the vehicle's commitment,
$\widehat{n}_i^{\text{fwd}}(t)$ is the RSU-independently-observed
forwarded count at time $t$, and $\varepsilon_{\text{obs}}$ is the
tolerance accounting for measurement uncertainty. A ZKP proof that
passes cryptographic verification but violates the consistency
condition is classified as FAIL, closing the fabricated commitment
attack vector.

\mypar
The revised DEBSC isolation condition incorporating all three ZKP
states is given in Equation~\eqref{eq:debsc}.

\begin{equation}
\label{eq:debsc}
\mathrm{Isolate}(v_i) =
\mathbf{1}\!\left[\,
\bigl(1 - \mathcal{R}_i(t)\bigr) > \theta_R
\;\wedge\;
\Pi_{\text{ZKP}}(v_i, t) \in
\left\{\,\texttt{FAIL},\; \texttt{ABSENT}\,\right\}
\,\right]
\end{equation}

As given in Equation~\eqref{eq:debsc}, $\theta_R$ is the
reputation isolation threshold and $\Pi_{\text{ZKP}}(v_i,t)$
now covers three states per Equation~\eqref{eq:zkp_state}. The
ABSENT state ensures that withholding a proof triggers the same
isolation pathway as submitting a false one. A legitimate vehicle
experiencing handoff-induced packet loss will retain valid ZKP
proofs for the packets it did forward and will not satisfy the
cryptographic gate, preventing false isolation.
% --- END SUPERVISOR REVISION BLOCK 5 ---

% --- SUPERVISOR REVISION BLOCK 6: new vehicle probationary mode ---
\mypar
\textbf{New Vehicle Probationary Mode.}
New vehicles have no blockchain reputation history, so
$|\mathcal{H}_i|=0$ and $\mathcal{R}_i(t)$ is initialised at the
Beta prior mean $\alpha/(\alpha+\beta)$, which may lie above
$\theta_R$, permanently blocking the statistical gate. An attacker
exploiting pseudonym rotation under IEEE~1609.2 can continuously
register fresh identities, always presenting as a new vehicle with
no adverse history and evading the statistical gate indefinitely.
The operating mode for vehicle $v_i$ is given in
Equation~\eqref{eq:vehicle_mode}.

\begin{equation}
\label{eq:vehicle_mode}
\mathrm{Mode}(v_i, t) =
\begin{cases}
\texttt{PROBATIONARY}
  & \text{if } |\mathcal{H}_i| < N_{\text{min}} \\[4pt]
\texttt{ACTIVE}
  & \text{otherwise}
\end{cases}
\end{equation}

As given in Equation~\eqref{eq:vehicle_mode}, $N_{\text{min}}$
is the minimum number of blockchain-recorded RSU interactions
before $\mathcal{R}_i(t)$ is considered statistically reliable.
In PROBATIONARY mode the reputation-based statistical gate is
suspended and replaced by a single-strike ZKP rule: any
$\Pi_{\text{ZKP}}(v_i,t) \in \{\texttt{FAIL}, \texttt{ABSENT}\}$
immediately triggers graduated response level~2 (mandatory
per-batch ZKP verification and rate-limiting) without requiring
a low reputation score. This prevents a freshly registered
attacker from freely exploiting the unestablished reputation
window, while avoiding false isolation of genuinely new
legitimate vehicles that are still building their history.
% --- END SUPERVISOR REVISION BLOCK 6 ---

\subsection{Federated Learning with Blockchain-Verified Gradient Integrity}

Each vehicle $v_i$ trains a local detection model $\mathbf{w}_i$ by minimising
the cross-entropy loss over its local dataset $\mathcal{D}_i$ of labelled
forwarding-log sequences is shown in Equation~\eqref{eq:local_loss}.

\begin{equation}
\label{eq:local_loss}
\mathcal{L}_i(\mathbf{w}) = \frac{1}{|\mathcal{D}_i|}
\sum_{(\mathbf{x},y)\in\mathcal{D}_i}
\ell\!\left(f(\mathbf{x};\mathbf{w}),\, y\right)
\end{equation}

As shown in Eq.~\eqref{eq:local_loss}, $\mathbf{w}$ is the model 
weight vector, $\mathbf{x}$ is an input forwarding-log feature vector, 
$y$ is the ground-truth label, $f(\mathbf{x};\mathbf{w})$ is the model 
prediction, and $\ell(\cdot)$ is the cross-entropy loss function.

The global model is updated each round $r$ via weighted averaging, where
weights are proportional to local dataset size to account for heterogeneous
vehicular observation volumes is shown in Equation~\eqref{eq:fedavg}.

\begin{equation}
\label{eq:fedavg}
\mathbf{w}^{(r+1)} =
\sum_{i \in \mathcal{A}} \frac{|\mathcal{D}_i|}{|\mathcal{D}_{\mathcal{A}}|}
\mathbf{w}_i^{(r)}
\end{equation}

As shown in Eq.~\eqref{eq:fedavg}, $\mathcal{A}$ is the set of 
accepted vehicles whose gradients passed the blockchain integrity 
check in Eq.~\eqref{eq:gradient_verify}, $|\mathcal{D}_i|$ is the 
local dataset size of vehicle $v_i$, and 
$|\mathcal{D}_{\mathcal{A}}| = \sum_{i \in \mathcal{A}}|\mathcal{D}_i|$ 
is the total accepted data volume.

\begin{equation}
\label{eq:gradient_verify}
\mathrm{Accept}(\Delta\mathbf{w}_i) =
\mathbf{1}\!\left[\,
\mathcal{H}_{\mathrm{BC}}(\Delta\mathbf{w}_i) =
\mathtt{Hash}(\Delta\mathbf{w}_i)
\,\right]
\end{equation}

As shown in Eq.~\eqref{eq:gradient_verify}, 
$\mathtt{Hash}(\Delta\mathbf{w}_i)$ is the locally computed hash of 
the received gradient and $\mathcal{H}_{\mathrm{BC}}(\Delta\mathbf{w}_i)$ 
is the pre-submitted on-chain commitment. A mismatch indicates 
tampering or gradient poisoning, and the update is discarded.

\subsection{LLM-Based Semantic Threat Scoring and Fusion}

In the full detection mode, the LLM processes tokenised blockchain forwarding
logs $\mathbf{x}_i^{(t)}$ as a sequence classification task. The threat score
is the softmax probability assigned to the malicious class by the LLM with
parameters $\boldsymbol{\theta}$ is shown in Equation~\eqref{eq:llm_score}.

\begin{equation}
\label{eq:llm_score}
\mathcal{Q}_i(t) =
\sum_{k \in \{S1,\ldots,S6\}}
\mathrm{softmax}\!\Bigl(
\mathrm{LLM}\!\left(\mathbf{x}_i^{(t)};\boldsymbol{\theta}\right)
\Bigr)_{k}
\end{equation}

As shown in Eq.~\eqref{eq:llm_score}, $\mathbf{x}_i^{(t)}$ is the
tokenised blockchain forwarding log of vehicle $v_i$ at time $t$,
$\boldsymbol{\theta}$ are the fine-tuned Qwen2.5-7B-Instruct LoRA
parameters, and $\mathcal{Q}_i(t)$ is the total softmax probability
mass assigned to the six malicious variant classes S1--S6 by a
seven-class classifier (BENIGN $+$ S1--S6). This equals $1 -
P(\mathrm{BENIGN}\mid\mathbf{x}_i^{(t)})$ and lies in $[0,1]$. The
seven-class formulation additionally yields per-variant class
probabilities that directly support per-variant MCC (M1,
Eq.~\eqref{eq:m1_mcc}) and AVCR (M3, Eq.~\eqref{eq:m3_avcr})
evaluation without requiring a separate classification head.

Because no single evidence source is individually sufficient for reliable
detection under all six attack variants, the final detection decision fuses
the rule-based aggregate signature score, the LLM semantic score, and the
blockchain reputation deficit through a learned linear combination is shown in Equation~\eqref{eq:fusion}.

\begin{equation}
\label{eq:fusion}
\hat{y}_i(t) = \mathbf{1}\!\left[\,
\mu_1\,\mathcal{S}_{\mathrm{total}}(v_i)
+ \mu_2\,\mathcal{Q}_i(t)
+ \mu_3\,\bigl(1 - \mathcal{R}_i(t)\bigr)
> \theta_{\mathrm{det}}
\,\right]
\end{equation}

As shown in Eq.~\eqref{eq:fusion}, $\mathcal{S}_{\mathrm{total}}(v_i)$ 
is the maximum binary signature score across all six attack signatures 
S1--S6, $\mathcal{Q}_i(t)$ is the LLM threat score from 
Eq.~\eqref{eq:llm_score}, $(1 - \mathcal{R}_i(t))$ is the blockchain 
reputation deficit from Eq.~\eqref{eq:reputation}, and 
$\theta_{\mathrm{det}}$ is the detection threshold. A positive verdict 
$\hat{y}_i(t) = 1$ triggers the mitigation pipeline.

\noindent where $\mu_1 + \mu_2 + \mu_3 = 1$ are fusion weights
optimised on the validation set.

\subsection{Post-Quantum Cryptographic Mitigation Formulations}

\textbf{CRYSTALS-Kyber Key Encapsulation.}
After a malicious node is isolated, all remaining vehicles must be re-keyed to
exclude the compromised node from future sessions. Rather than classical ECDH,
which is vulnerable to quantum attacks, we use CRYSTALS-Kyber whose security
reduces to the Module Learning With Errors (M-LWE) problem. Given recipient
public key $\mathbf{pk}$, a fresh session key and its ciphertext are produced as shown in Equation~\eqref{eq:kyber_enc}.

\begin{equation}
\label{eq:kyber_enc}
(K,\, c) = \mathrm{Kyber.Enc}(\mathbf{pk},\, m),
\quad m \xleftarrow{\$} \{0,1\}^{256}
\end{equation}

As shown in Eq.~\eqref{eq:kyber_enc}, $K$ is the encapsulated session 
key, $c$ is its ciphertext, and $m \xleftarrow{\$} \{0,1\}^{256}$ is 
the uniformly random encapsulation seed applied to recipient public key 
$\mathbf{pk}$.

\begin{equation}
\label{eq:kyber_dec}
K = \mathrm{Kyber.Dec}(\mathbf{sk},\, c)
\end{equation}

As shown in Eq.~\eqref{eq:kyber_dec}, $\mathbf{sk}$ is the private key 
corresponding to $\mathbf{pk}$ in Eq.~\eqref{eq:kyber_enc}, and $c$ is 
the ciphertext produced there. Only the holder of the correct 
$\mathbf{sk}$ can recover $K$, ensuring isolated vehicles excluded from 
key refresh cannot derive the new group session key.

\textbf{CRYSTALS-Dilithium Flow Rule Authentication.}
SDN controller isolation commands and smart contract transactions must be
authenticated to prevent replay and forgery. The controller signs each
flow modification message $\mathcal{M}$ with its Dilithium private key
and recipients verify before applying the rule as shown in Equation~\eqref{eq:dilithium_sign}.

\begin{equation}
\label{eq:dilithium_sign}
\sigma = \mathrm{Dilithium.Sign}(\mathbf{sk}_c,\, \mathcal{M})
\end{equation}

As shown in Eq.~\eqref{eq:dilithium_sign}, $\mathbf{sk}_c$ is the 
controller's Dilithium signing key and $\mathcal{M}$ is the isolation 
FlowMod command. This prevents a spoofed or replayed controller message 
from installing false blocking rules on OpenFlow switches.

\begin{equation}
\label{eq:dilithium_verify}
b = \mathrm{Dilithium.Verify}(\mathbf{pk}_c,\, \mathcal{M},\, \sigma)
\in \{0,1\}
\end{equation}

As shown in Eq.~\eqref{eq:dilithium_verify}, $\mathbf{pk}_c$ is the 
controller's public key, and $b \in \{0,1\}$ indicates authenticity: 
$b = 1$ installs the block rule; $b = 0$ causes the command to be 
rejected.

\textbf{ZKP Forwarding Proof via Pedersen Commitment.}
A grey hole node that drops packets cannot produce a valid forwarding proof for
the dropped packets. Each vehicle commits to its forwarded packet count
$n_i^{\text{fwd}}$ using a Pedersen commitment with random blinding factor
$r_i$, so that the commitment is binding (count cannot be changed) and
hiding (count is not revealed) as shown in Equation~\eqref{eq:pedersen}.

\begin{equation}
\label{eq:pedersen}
C_i = g^{n_i^{\text{fwd}}} \cdot h^{r_i} \pmod{p}
\end{equation}

As shown in Eq.~\eqref{eq:pedersen}, $g$ and $h$ are public group 
generators, $r_i$ is a random blinding factor chosen by $v_i$, and 
$p$ is a large prime modulus.

The vehicle then generates a zero-knowledge proof $\pi_i$ demonstrating
knowledge of a valid opening $(n_i^{\text{fwd}}, r_i)$ consistent with $C_i$
and the independently observable packet counts on the blockchain, without
disclosing packet content as shown in Equation~\eqref{eq:zkp_proof}.

\begin{equation}
\label{eq:zkp_proof}
\pi_i = \mathrm{ZKP.Prove}\!\left(C_i,\, n_i^{\text{fwd}},\, r_i\right)
\end{equation}

As shown in Eq.~\eqref{eq:zkp_proof}, $C_i$ is the Pedersen commitment 
from Eq.~\eqref{eq:pedersen}, $n_i^{\text{fwd}}$ is the claimed 
forwarded count, and $r_i$ is the blinding factor. A grey hole attacker 
that dropped packets cannot produce a valid $\pi_i$, as its committed 
$n_i^{\text{fwd}}$ will not match the observable receipt count on the 
blockchain.

\textbf{Threshold Signature for Collective Blacklisting.}
To prevent any single RSU from unilaterally blacklisting a vehicle — which
could be exploited by a compromised RSU — node isolation requires a
$(k,n)$-threshold co-signature from $k$ independent RSUs. Each participating
RSU $j$ produces a partial signature over the blacklist message $\mathcal{B}(v_i)$ as shown in Equation~\eqref{eq:threshold_partial}.

\begin{equation}
\label{eq:threshold_partial}
\sigma_j = \mathrm{TS.PartialSign}\!\left(\mathbf{sk}_j,\, \mathcal{B}(v_i)\right),
\quad j = 1,\ldots,k
\end{equation}

As shown in Eq.~\eqref{eq:threshold_partial}, $\mathbf{sk}_j$ is the 
signing key of RSU $j$, $\mathcal{B}(v_i)$ is the blacklist message 
identifying suspect vehicle $v_i$, and $k$ is the minimum co-signature 
threshold.

The partial signatures are combined off-chain and a single compact group
signature is produced for on-chain verification as shown in Equation~\eqref{eq:threshold_combine}.

\begin{equation}
\label{eq:threshold_combine}
\sigma^* = \mathrm{TS.Combine}\!\left(\{\sigma_j\}_{j=1}^{k}\right)
\end{equation}

As shown in Eq.~\eqref{eq:threshold_combine}, the $k$ partial 
signatures from Eq.~\eqref{eq:threshold_partial} are combined 
off-chain into a single compact aggregate signature $\sigma^*$ for 
on-chain submission.

\begin{equation}
\label{eq:threshold_verify}
b = \mathrm{TS.Verify}\!\left(\mathbf{pk}_{\mathrm{group}},\,
\mathcal{B}(v_i),\, \sigma^*\right) \in \{0,1\}
\end{equation}

As shown in Eq.~\eqref{eq:threshold_verify}, $\mathbf{pk}_{\mathrm{group}}$ 
is the collective public key of the $k$-of-$n$ RSU threshold scheme, 
and $b = 1$ confirms that at least $k$ independent RSUs have endorsed 
the blacklisting decision.

\noindent Isolation is executed by the smart contract only when $b = 1$,
ensuring that the decision has the collective cryptographic endorsement of
$k$ independent infrastructure nodes — a design that remains secure as long
as fewer than $k$ RSUs are compromised simultaneously.

% --- SUPERVISOR REVISION BLOCK 9: degraded mode for sparse RSU deployments ---
\mypar
\textbf{Degraded Mode for Sparse RSU Deployments.}
The $(k,n)$-threshold co-signature requirement assumes at least
$k$ RSUs are reachable within timeout $T_q$. In sparse or rural
deployments fewer than $k$ RSUs may be available, permanently
blocking full isolation and leaving an active attacker
unmitigated. The degraded mode activation condition is given in
Equation~\eqref{eq:degraded_mode}.

\begin{equation}
\label{eq:degraded_mode}
\mathrm{Degraded}(t) =
\mathbf{1}\!\left[\,
\left|\left\{
RSU_j : \text{reachable within } T_q \text{ at } t
\right\}\right| < k
\,\right]
\end{equation}

As given in Equation~\eqref{eq:degraded_mode}, the flag fires
when the count of reachable RSUs falls below the blacklisting
quorum $k$ within timeout $T_q$. When degraded mode is active
the DEBSC applies graduated response level~2 (mandatory
per-batch ZKP verification and traffic rate-limiting of the
suspect vehicle) as temporary containment without full isolation.
The degraded mode event is recorded on the blockchain so that the
full isolation decision is re-evaluated automatically when quorum
is restored. The attacker is rate-limited rather than allowed to
operate freely during sparse-coverage periods.
% --- END SUPERVISOR REVISION BLOCK 9 ---

% --- SUPERVISOR REVISION BLOCK 15: VRF-based dynamic endorser selection ---
\subsection{VRF-Based Dynamic Endorser Selection}
\label{sec:endorser_selection}

The threshold co-signature scheme in
Eqs.~\eqref{eq:threshold_partial}--\eqref{eq:threshold_combine}
originally selects whichever $k$ RSUs respond first within timeout
$T_q$ — a static, predictable arrangement. An adversary satisfying
T4 who has compromised a known fixed subset of RSUs can reliably
participate in or block every endorsement round. To eliminate this
predictability, endorsers are selected per-transaction at chaincode
level via a Verifiable Random Function (VRF), filtered by the
per-RSU trust score maintained on the blockchain, with a graduated
bypass mechanism ensuring liveness when the eligible pool is
depleted.

\mypar
\textbf{Step 1 — Eligible Endorser Pool.}
Each RSU $RSU_j \in \mathcal{R}$ holds a trust score
$\mathcal{T}_{r_j}(t)$ maintained on the blockchain by BFT
consensus, computed by the same Bayesian model applied to vehicles
(Eq.~\eqref{eq:trust_update}). RSUs below the eligibility threshold
$\theta_{\text{RSU}}$ or in probationary mode ($|\mathcal{H}_j| <
N_{\min}$) are excluded. The eligible endorser pool is given in
Equation~\eqref{eq:endorser_pool}.

\begin{equation}
\label{eq:endorser_pool}
\mathcal{E}(t) =
\bigl\{\,RSU_j \in \mathcal{R} :
\mathcal{T}_{r_j}(t) \geq \theta_{\text{RSU}}
\;\wedge\;
|\mathcal{H}_j| \geq N_{\min}
\,\bigr\}
\end{equation}

As given in Equation~\eqref{eq:endorser_pool}, $\theta_{\text{RSU}}$
is the minimum trust score required for endorser candidacy and
$\mathcal{H}_j$ is the set of blockchain-recorded interactions
of RSU $j$. Only RSUs satisfying both the trust threshold and the
minimum observation count enter the candidate pool, preventing
newly joined or recently compromised RSUs from influencing
endorsement decisions.

\mypar
\textbf{Step 2 — Per-Transaction VRF Seed.}
For each chaincode transaction $tx$ requiring endorsement, a
deterministic but unpredictable seed is committed to the blockchain
before any RSU evaluates it, given in
Equation~\eqref{eq:vrf_seed}.

\begin{equation}
\label{eq:vrf_seed}
s_{tx} =
\texttt{Hash}\!\left(
\text{tx\_id}
\;\|\;
\text{block\_height}
\;\|\;
t
\;\|\;
\mathcal{C}_{\text{prev}}
\right)
\end{equation}

As given in Equation~\eqref{eq:vrf_seed}, $\text{tx\_id}$ is the
unique transaction identifier, $\text{block\_height}$ is the
current ledger height, $t$ is the wall-clock timestamp, and
$\mathcal{C}_{\text{prev}}$ is the hash of the immediately
preceding block header. Binding the seed to the immutable chain
state ensures no single party can predict or bias it before
commitment.

\mypar
\textbf{Step 3 — VRF Evaluation and Endorser Ranking.}
Each RSU $RSU_j \in \mathcal{E}(t)$ independently evaluates the
VRF using its Dilithium private key $\mathbf{sk}_j$ and the
published seed. The VRF output and its proof are given in
Equation~\eqref{eq:vrf_eval}.

\begin{equation}
\label{eq:vrf_eval}
(\beta_j,\;\pi_j^{\text{VRF}}) =
\texttt{VRF}(\mathbf{sk}_j,\; s_{tx})
\end{equation}

As given in Equation~\eqref{eq:vrf_eval}, $\beta_j \in
\{0,1\}^{256}$ is the pseudorandom VRF output and
$\pi_j^{\text{VRF}}$ is the verifiable proof that
$\beta_j$ was computed correctly from $\mathbf{sk}_j$ and
$s_{tx}$. Any party may verify correctness without knowledge
of $\mathbf{sk}_j$, as given in
Equation~\eqref{eq:vrf_verify}.

\begin{equation}
\label{eq:vrf_verify}
b =
\texttt{VRF.Verify}
\!\left(
\mathbf{pk}_j,\;
s_{tx},\;
\beta_j,\;
\pi_j^{\text{VRF}}
\right)
\in \{0,1\}
\end{equation}

As given in Equation~\eqref{eq:vrf_verify}, $\mathbf{pk}_j$ is
RSU $j$'s public key and $b = 1$ confirms the output is genuine.
RSUs in $\mathcal{E}(t)$ are ranked by $\beta_j$ ascending; the
selected endorser set is the top $k_{\text{end}}$ entries with
verified proofs, as given in Equation~\eqref{eq:endorser_select}.

\begin{equation}
\label{eq:endorser_select}
\Omega(t) =
\underset{j \in \mathcal{E}(t)}
{\operatorname{top-}k_{\text{end}}}
\;\text{ranked by } \beta_j \text{ ascending,}
\quad
\text{s.t.}
\quad
\texttt{VRF.Verify}
\!\left(
\mathbf{pk}_j,\; s_{tx},\;
\beta_j,\; \pi_j^{\text{VRF}}
\right) = 1
\end{equation}

As given in Equation~\eqref{eq:endorser_select}, only RSUs whose
VRF proof verifies successfully are eligible for selection,
and the lowest-output entries are preferred, producing a
per-transaction endorser set that is unpredictable to any
adversary who does not control the seed.

\mypar
\textbf{Step 4 — Minimum Endorser Set Size.}
The minimum endorser count must guarantee at least one honest
RSU in $\Omega(t)$ under the T4 bounded collusion assumption,
where at most $f_{\max} = \lfloor(|\mathcal{E}(t)|-1)/3\rfloor$
RSUs in the eligible pool may be Byzantine. The minimum is given
in Equation~\eqref{eq:kend}.

\begin{equation}
\label{eq:kend}
k_{\text{end}} =
\max\!\Bigl(
k_{\min},\;
\Bigl\lceil
|\mathcal{E}(t)| \cdot \alpha_{\text{end}}
\Bigr\rceil
\Bigr),
\quad
k_{\min} = 3,
\quad
\alpha_{\text{end}} \in
\!\left(\tfrac{1}{3},\; 1\right]
\end{equation}

As given in Equation~\eqref{eq:kend}, $k_{\min} = 3$ is the
hard practical floor below which endorsement evidence is
considered statistically insufficient, and
$\alpha_{\text{end}} > 1/3$ ensures the selected set exceeds
the maximum Byzantine fraction of the eligible pool, guaranteeing
at least one honest endorser in $\Omega(t)$ with high probability.
The value of $\alpha_{\text{end}}$ is a deployment parameter
calibrated to the RSU population density.

\mypar
\textbf{Step 5 — Bypass Mechanism.}
When the eligible pool is depleted the framework degrades
gracefully through four ordered modes, given in
Equation~\eqref{eq:endorser_mode}.

\begin{equation}
\label{eq:endorser_mode}
\mathrm{EndorserMode}(t) =
\begin{cases}
\texttt{NORMAL}
  & \text{if } |\mathcal{E}(t)| \geq k_{\text{end}} \\[4pt]
\texttt{RELAXED}
  & \text{if } k_{\min} \leq |\mathcal{E}(t)| < k_{\text{end}} \\[4pt]
\texttt{DEFERRED}
  & \text{if } 0 < |\mathcal{E}(t)| < k_{\min} \\[4pt]
\texttt{EMERGENCY}
  & \text{if } \mathcal{E}(t) = \emptyset
\end{cases}
\end{equation}

As given in Equation~\eqref{eq:endorser_mode}, the four modes
operate as follows. In \texttt{NORMAL} mode, $k_{\text{end}}$
endorsers are selected from $\mathcal{E}(t)$ via
Eq.~\eqref{eq:endorser_select} and full isolation proceeds.
In \texttt{RELAXED} mode, all $|\mathcal{E}(t)|$ available
trusted RSUs act as endorsers; the transaction is committed
with a reduced-confidence flag and the RSU cross-reference
tolerance $\varepsilon_{\text{obs}}$ (Eq.~\eqref{eq:rsu_crossref})
is tightened to compensate for the smaller set. In
\texttt{DEFERRED} mode, the eligible pool is below the hard
floor: the DEBSC applies graduated response level~2 (mandatory
per-batch ZKP verification and rate-limiting,
Eq.~\eqref{eq:degraded_mode}) as temporary containment and
queues the isolation decision for re-evaluation at the next
observation window. In \texttt{EMERGENCY} mode, no RSU
satisfies $\theta_{\text{RSU}}$: the RSU holding the globally
highest $\mathcal{T}_{r_j}(t)$ regardless of threshold is
used as a single emergency endorser, the event is committed
to the blockchain as an emergency alert record, and a
network-wide RSU trust audit is triggered via the transport
authority PKI.

\mypar
\textbf{Integration with threshold signature scheme.}
The partial signatures in Eq.~\eqref{eq:threshold_partial} are
now collected exclusively from $\Omega(t)$, replacing the
static first-$k$-to-respond behaviour. Each RSU in $\Omega(t)$
attaches its VRF proof $\pi_j^{\text{VRF}}$ alongside its
partial signature $\sigma_j$, allowing the chaincode to verify
both endorser eligibility and signature authenticity atomically
before aggregation into $\sigma^*$ via
Eq.~\eqref{eq:threshold_combine}. The seed $s_{tx}$ is
committed to the ledger before any RSU evaluates the VRF,
ensuring no RSU can predict selection in advance and
selectively misbehave only when chosen.
% --- END SUPERVISOR REVISION BLOCK 15 ---

% --- SUPERVISOR REVISION BLOCK 8: DKG for PQC-LKH root key (new subsection) ---
\subsection{Distributed Key Generation for PQC-LKH Root Key}
\label{sec:dkg}

In the original design the SDN controller alone signed and broadcast
PQC-LKH re-keying messages using its Dilithium key. A compromised
controller holding a valid Dilithium key could therefore broadcast
fraudulent group re-keying messages and disrupt all vehicle sessions
unilaterally. To eliminate this single point of cryptographic
failure, the root Kyber key pair of the PQC-LKH tree is generated
jointly by the active controller and RSUs via Pedersen Distributed
Key Generation (DKG), and all re-keying broadcasts require a
threshold co-signature from the controller and RSUs.

\mypar
\textbf{DKG Key Share Generation.}
Each participant $p \in \mathcal{P} = \{c\} \cup
\{RSU_1, \ldots, RSU_n\}$ selects a secret polynomial
$f_p(x) = s_p + a_{p,1}x + \cdots + a_{p,t-1}x^{t-1}$
over $\mathbb{Z}_q$ and broadcasts commitments to its
coefficients. The individual key share held by participant $p$ is
given in Equation~\eqref{eq:dkg_share}.

\begin{equation}
\label{eq:dkg_share}
\mathrm{sk}_p^{\text{share}} =
\sum_{q \in \mathcal{P}} f_q(p)
\end{equation}

As given in Equation~\eqref{eq:dkg_share}, $\mathcal{P}$ is the
set of all DKG participants and $f_q(p)$ is participant $q$'s
secret polynomial evaluated at participant $p$'s index. No single
participant holds the complete private key; at least $t$
participants must cooperate to perform any root key operation.

The aggregate root public key reconstructed from participant
commitments is given in Equation~\eqref{eq:dkg_pubkey}.

\begin{equation}
\label{eq:dkg_pubkey}
\mathbf{pk}_{\text{root}} =
\sum_{p \in \mathcal{P}} \mathbf{pk}_p^{\text{share}}
\end{equation}

As given in Equation~\eqref{eq:dkg_pubkey},
$\mathbf{pk}_p^{\text{share}}$ is participant $p$'s contribution
to the root public key. The root public key is publicly verifiable
from committed values without revealing any private share.

\mypar
\textbf{Threshold Re-Keying Broadcast Signature.}
All PQC-LKH re-keying broadcast messages require a combined
Dilithium threshold signature from the active controller
(contributing one share) and $k_{\text{key}}-1$ RSU participants.
The aggregate re-keying authorisation signature is given in
Equation~\eqref{eq:rekey_sig}.

\begin{equation}
\label{eq:rekey_sig}
\sigma_{\text{rekey}}^* =
\mathrm{TS.Combine}\!\left(
\sigma_c^{\text{rekey}},\;
\left\{\sigma_j^{\text{rekey}}\right\}_{j=1}^{k_{\text{key}}-1}
\right)
\end{equation}

As given in Equation~\eqref{eq:rekey_sig},
$\sigma_c^{\text{rekey}}$ is the controller's Dilithium partial
signature over the re-keying message and
$\{\sigma_j^{\text{rekey}}\}$ are the RSU partial signatures.
A vehicle accepts a re-keying broadcast only when
$\mathrm{TS.Verify}(\mathbf{pk}_{\text{group}}, m_{\text{rekey}},
\sigma_{\text{rekey}}^*)$ passes. This ensures neither the
controller alone nor any single RSU can unilaterally distribute
fraudulent session keys; the controller is involved in key
operations but holds no monopoly.
% --- END SUPERVISOR REVISION BLOCK 8 ---

% ============================================================
% INSERT AT END OF SECTION 3.3, BEFORE ALGORITHM SECTION
% ============================================================

\subsection{PQC-LKH: Post-Quantum Logical Key Hierarchy for Group Re-Keying}
\label{sec:pqc_lkh}

When a vehicle $v_i$ is isolated, all remaining $N-1$ vehicles must
receive an updated group session key $K_{\text{grp}}$ that excludes
$v_i$. Unicast Kyber re-keying requires $O(N-1)$ KEM operations,
which is prohibitive in dense SDVNs. We adapt the Logical Key
Hierarchy (LKH)~\cite{wallner1999lkh} to PQC, reducing re-keying
cost to $O(\log N)$.

% --- SUPERVISOR REVISION BLOCK 10: dynamic PQC-LKH membership (lazy policy) ---
Vehicles are organised into a binary tree $\mathcal{T}$ of depth
$\lceil \log_2 N \rceil$. Each internal node $u \in \mathcal{T}$
holds a Kyber-1024 key pair $(\mathbf{pk}_u, \mathbf{sk}_u)$. The
original design assumed a static group of $N$ vehicles, which is
unrealistic in a high-mobility SDVN. SHIELD-GH adopts a
\emph{lazy membership policy} that preserves the $O(\log N)$
re-keying complexity under continuous vehicle joins and departures
without restructuring the tree.

\textbf{Lazy Departure.}
When vehicle $v_i$ departs (leaves RSU coverage or is isolated),
its leaf node $\ell_i$ is marked \texttt{VACANT} but is not
removed from the tree. A vacant leaf holds no private key
material and does not participate in subsequent re-keying
operations. The tree structure is fully preserved.

\textbf{Lazy Join.}
A new vehicle $v_{\text{new}}$ is assigned to the
lowest-indexed \texttt{VACANT} leaf. The path from that leaf
to the root is refreshed with new Kyber key pairs, and
$v_{\text{new}}$ receives the refreshed path keys via
RSU-threshold-signed broadcast (Equation~\eqref{eq:rekey_sig}).
This requires at most $\lceil \log_2 N \rceil$ Kyber operations,
identical in cost to an isolation re-keying event. The tree depth
increases only when all leaves are occupied.

Each vehicle $v_j$ holds the keys of all non-vacant nodes on its
root-to-leaf path. The key set for $v_j$ under the lazy policy is
given in Equation~\eqref{eq:lkh_keys}.

\begin{equation}
\label{eq:lkh_keys}
\mathcal{K}_j =
\bigl\{(\mathbf{pk}_u, \mathbf{sk}_u) :
u \in \mathrm{path}(v_j \to \mathrm{root}),\;
\ell_u \neq \texttt{VACANT}
\bigr\}
\end{equation}

As given in Equation~\eqref{eq:lkh_keys}, the vacant leaf
exclusion condition $\ell_u \neq \texttt{VACANT}$ ensures that
departed vehicles' key material is excluded from all active path
key sets. Path membership determines which re-keying broadcasts
in Equation~\eqref{eq:lkh_grp} a vehicle can decrypt.
% --- END SUPERVISOR REVISION BLOCK 10 ---

\begin{equation}
\label{eq:lkh_grp}
(K_{\text{grp}},\, c_{\mathrm{root}}) =
\mathrm{Kyber.Enc}(\mathbf{pk}_{\mathrm{root}},\, m),
\quad m \xleftarrow{\$} \{0,1\}^{256}
\end{equation}

As shown in Eq.~\eqref{eq:lkh_grp} encapsulates the new group session key under 
the refreshed root public key using the same KEM as 
Eq.~\eqref{eq:kyber_enc}, where $m \xleftarrow{\$} \{0,1\}^{256}$ 
is a fresh random seed. Only vehicles holding a valid $\mathbf{sk}_u$ 
for at least one refreshed path node can decapsulate and derive 
$K_{\mathrm{grp}}$.

Upon isolating $v_i$ (leaf node $\ell_i$), only the keys along
the path $\mathrm{path}(\ell_i \to \mathrm{root})$ are refreshed.
For each affected node $u$ on this path, a new key pair is generated
and the updated key is re-encrypted under the sibling subtree keys,
requiring exactly $\lceil \log_2 N \rceil$ Kyber operations, as shown in Equation~\eqref{eq:lkh_rekey}.

\begin{equation}
\label{eq:lkh_rekey}
(K_u^{\mathrm{new}},\, c_u) =
\mathrm{Kyber.Enc}(\mathbf{pk}_{u}^{\mathrm{sib}},\, k_u^{\mathrm{new}}),
\quad \forall\, u \in \mathrm{path}(\ell_i \to \mathrm{root})
\end{equation}

As shown in Eq.~\eqref{eq:lkh_rekey} shows that at each node $u$ on the isolated 
vehicle's root path, a new key $k_u^{\mathrm{new}}$ is encapsulated 
under the sibling subtree's public key $\mathbf{pk}_u^{\mathrm{sib}}$, 
producing ciphertext $c_u$. One operation per tree level yields 
exactly $\lceil \log_2 N \rceil$ total Kyber operations.

RSUs broadcast $\lceil \log_2 N \rceil$ update messages, each carrying
one re-encrypted node key and authenticated by the threshold re-keying
signature $\sigma_{\text{rekey}}^*$ (Eq.~\eqref{eq:rekey_sig}), which
requires a combined Dilithium contribution from the active controller
and $k_{\text{key}}-1$ RSUs before any vehicle accepts the update. Vehicles not on the affected path
derive the new $K_{\text{grp}}$ by decapsulating the updated ancestor
key using their held $\mathbf{sk}_u$. The isolated vehicle $v_i$,
lacking any valid $\mathbf{sk}_u$ for the refreshed path, is
cryptographically excluded from the new group session without
contacting it directly.

% \begin{figure}[htbp]
% 	\centering
% 	\includegraphics[width=\textwidth,height=0.45\textheight,keepaspectratio]{Figures/figure6_pqc_lkh.png}
% 	\caption{PQC-LKH Binary Tree for Post-Quantum Group Re-Keying}
% 	\label{fig:pqc-lkh}
% \end{figure}

\begin{figure}[H]
	\centering
	\includegraphics[
		width=\textwidth,
		height=0.7\textheight,
		keepaspectratio
	]{Figures/figure6_pqc_lkh.png}
	\caption{PQC-LKH Binary Tree for Post-Quantum Group Re-Keying}
	\label{fig:pqc-lkh}
\end{figure}

As shown in figure~\ref{fig:pqc-lkh} illustrates the PQC-LKH binary tree structure for post-quantum group re-keying. As shown in Figure~\ref{fig:pqc-lkh}, the PQC-LKH scheme adopts a binary tree structure for an example network of $N=8$ vehicles. Each vehicle occupies a leaf node (V1--V8) and each internal node holds a Kyber-1024 key pair. Vehicle V3 has been identified as a grey hole attacker and isolated. The re-keying procedure refreshes only the nodes along the path from V3's leaf to the root — the K\_sub parent node, the K\_left subtree node, and K\_root — totalling $\lceil\log_2 8\rceil = 3$ Kyber operations. For each refreshed node, the updated key is encrypted under the sibling subtree's existing public key and broadcast, so only vehicles on the unaffected side of each split (e.g., V1 and V2 receive an encrypted update via their K\_sub) can decrypt and derive the new group session key. V3, lacking the private keys for any refreshed node on the path, is permanently and cryptographically excluded from the new group session without the RSU needing to contact it directly. The embedded efficiency comparison table in the figure quantifies the benefit: naïve unicast re-keying requires $N-1 = 7$ messages and $7$ KEM operations at $O(N)$ complexity, while PQC-LKH achieves the same exclusion guarantee with $\log_2 N = 3$ messages and operations at $O(\log N)$ — a factor of 2.33$\times$ reduction for $N=8$ that grows logarithmically with network size.

% ============================================================
% INSERT AS SECTION 3.5 — ALGORITHMS
% ============================================================

\section{Detection and Mitigation Algorithms}
\label{sec:algorithms}

This section presents the four algorithms that operationalise the detection and
mitigation framework described in the preceding sections. Each algorithm is a
self-contained procedure with clearly defined inputs, outputs, and decision
logic. Algorithms~\ref{alg:lw_dp_det} and~\ref{alg:lw_cp_det} implement the
lightweight rule-based detection path for data-plane and controller-plane
variants respectively. Algorithm~\ref{alg:fv_det} implements the full-mode
pipeline combining LLM semantic scoring, federated learning inference, and
blockchain reputation fusion. Algorithm~\ref{alg:pqc_mit} executes the
post-quantum cryptographic mitigation sequence — ZKP verification, threshold
blacklisting, flow modification, and PQC-LKH group re-keying — once a vehicle
is confirmed as a grey hole attacker.

\subsection{Lightweight Rule-Based Grey Hole Detection — Data Plane}

% --- SUPERVISOR REVISION BLOCK 11: Algorithm 1 revised with probationary mode ---
\begin{algorithm}[H]
\caption{Lightweight Rule-Based Grey Hole Detection --- Data Plane
(\texttt{LW-DP-Det})}
\begin{algorithmic}[1]
\small
\Procedure{\texttt{LW-DP-Det}}{$v_i, t, W, \tau_f,
\varepsilon_f, \tau_{it}, \gamma_{it}, \tau_{ts},
N_{\text{min}}$}
  \State \textbf{// Probationary mode: applies to new vehicles
    with insufficient blockchain history}
  \If{$|\mathcal{H}_i| < N_{\text{min}}$}
    \Comment{Eq.~\eqref{eq:vehicle_mode}}
    \State $\pi_{\text{state}} \gets
      \texttt{ZKP-CROSSCHECK}(v_i, t, T_{\text{zkp}},
      \varepsilon_{\text{obs}})$
    \Comment{Algorithm~\ref{alg:zkp_crosscheck}}
    \If{$\pi_{\text{state}} \in
      \{\texttt{FAIL},\, \texttt{ABSENT}\}$}
      \State \textbf{return}
        $(\texttt{PROBATIONARY\_ALERT}, v_i)$
      \Comment{Triggers graduated response level 2}
    \Else
      \State \textbf{return} $(\texttt{BENIGN}, v_i)$
    \EndIf
  \EndIf
  \State \textbf{// Active mode: standard signature evaluation}
  \State $\text{PDR\_hist} \gets
    \texttt{GET\_BLOCKCHAIN\_LOG}(v_i, t, W)$
  \State $\widehat{\text{PDR}} \gets
    \texttt{MOBILITY\_CORRECT}
    (\text{PDR\_hist},\, s_i(t))$
  \Comment{Eq.~\eqref{eq:mobility_pdr}}
  \State $\sigma^2 \gets
    \texttt{VARIANCE}(\widehat{\text{PDR}})$
  \Comment{Eq.~\eqref{eq:pdr_var}}
  \If{$\overline{\widehat{\text{PDR}}} < \tau_f$
      \textbf{and} $\sigma^2 < \varepsilon_f$}
    \State \textbf{return} $(\texttt{DP-FR}, v_i)$
    \Comment{Signature S1}
  \EndIf
  \State $R_m \gets
    \texttt{AUTOCORR}(\text{PDR\_hist}, \tau_{it})$
  \If{$\max_{T} R_m(T) > \gamma_{it}$}
    \State \textbf{return} $(\texttt{DP-IT}, v_i)$
    \Comment{Signature S2}
  \EndIf
  \State $D_{KL} \gets
    \texttt{PER\_SOURCE\_PDR\_KL}(v_i, t, W)$
  \Comment{Eq.~\eqref{eq:sig_dpts}}
  \If{$D_{KL} > \tau_{ts}$}
    \State \textbf{return} $(\texttt{DP-TS}, v_i)$
    \Comment{Signature S3}
  \EndIf
  \State \textbf{return} $(\texttt{BENIGN}, v_i)$
\EndProcedure
\end{algorithmic}
\label{alg:lw_dp_det}
\end{algorithm}
% --- END SUPERVISOR REVISION BLOCK 11 ---

Algorithm~\ref{alg:lw_dp_det} presents the \texttt{LW-DP-Det} lightweight rule-based grey hole detection procedure for the data plane. As shown in Algorithm~\ref{alg:lw_dp_det}, \texttt{LW-DP-Det} performs lightweight rule-based grey hole detection on the data plane. It takes as input the vehicle
identifier $v_i$, the current time $t$, observation window $W$, and five
threshold parameters: $\tau_f$ and $\varepsilon_f$ govern fixed-rate detection;
$\tau_{it}$ and $\gamma_{it}$ govern intermittent detection; and $\tau_{ts}$
governs target-specific detection. The algorithm first retrieves the vehicle's
historical forwarding log from the blockchain (line~1) and applies the mobility
correction of Eq.~\eqref{eq:mobility_pdr} to remove expected handoff-induced
loss from the PDR signal (line~2). It then evaluates the three data-plane attack
signatures in priority order. The fixed-rate test (Signature S1,
Eq.~\eqref{eq:sig_dpfr}) fires when mean corrected PDR falls below $\tau_f$ and
variance is below $\varepsilon_f$, identifying a node that consistently drops a
fixed fraction of traffic. If this test does not fire, the intermittent test
(Signature S2, Eq.~\eqref{eq:sig_dpit}) computes the autocorrelation of the
binary drop indicator and fires if a periodic pattern is found — the hallmark of
an on/off attacker. Finally, the target-specific test (Signature S3,
Eq.~\eqref{eq:sig_dpts}) computes the KL-divergence of the per-source PDR
distribution and fires when the distribution is non-uniform, indicating
selective per-source dropping. The algorithm returns a two-element tuple: one
of \texttt{DP-FR}, \texttt{DP-IT}, or \texttt{DP-TS} paired with $v_i$ on
detection, or \texttt{BENIGN} if all three signatures pass.

% \begin{figure}[htbp]
% 	\centering
% 	\includegraphics[width=\textwidth,height=0.45\textheight,keepaspectratio]{Figures/figure3_lightweight_sequence}
% 	\caption{Lightweight Mode UML Sequence Diagram}
% 	\label{fig:lightweight-sequence}
% \end{figure}

\begin{figure}[H]
	\centering
	\includegraphics[
		width=\textwidth,
		height=0.7\textheight,
		keepaspectratio
	]{Figures/figure3_lightweight_sequence}
	\caption{Lightweight Mode UML Sequence Diagram}
	\label{fig:lightweight-sequence}
\end{figure}

Figure~\ref{fig:lightweight-sequence} presents the complete UML message sequence for the lightweight detection and mitigation path, traced step by step. The malicious vehicle selectively drops packets (step 1) while legitimate vehicles forward normally (step 2). The RSU/Network Monitor computes mobility-corrected PDR per vehicle (step 3), queries the Blockchain Ledger for the historical baseline (step 4), and receives the historical reputation score (step 5). The six formal attack signatures S1--S6 are evaluated against the corrected PDR (step 6), and if a signature fires alongside a low reputation score (step 7), a signed suspicious event log is submitted to the ledger (step 8). The DEBSC dual-evidence gate is then checked (step 9): if the ZKP forwarding proof is valid or reputation is insufficient, the sequence aborts and no isolation occurs. On a positive verdict, $k$ partial threshold signatures are collected from RSUs (step 10), combined, and a threshold-signed FlowMod is broadcast directly to the OpenFlow Switch to block the attacker (step 11). Finally, the reputation score is updated to zero and committed to the ledger (step 13). At no point does the SDN controller issue or approve an isolation command — this is enforced by design to eliminate the controller as a single point of compromise.

\subsection{Lightweight Rule-Based Grey Hole Detection — Controller Plane}

Algorithm~\ref{alg:lw_cp_det} presents the \texttt{LW-CP-Det} lightweight rule-based grey hole detection procedure for the controller plane. As shown in Algorithm~\ref{alg:lw_cp_det},
\texttt{LW-CP-Det} performs lightweight rule-based grey hole
detection on the controller plane. It takes as input the
controller identifier $c$, current time $t$, drop probability threshold
$\tau_c$, autocorrelation threshold $\gamma_c$, and observation window $W$.
Unlike the data-plane algorithm, which analyses per-vehicle PDR statistics,
this algorithm inspects the OpenFlow flow rules installed by the controller.
It first retrieves all active flow rules from the controller (line~1) and
scans each rule for a drop action: if a rule unconditionally drops traffic
with probability above $\tau_c$, the fixed-rate controller signature fires
(Signature S4, Eq.~\eqref{eq:sig_cpfr}); if a rule drops traffic conditioned
on a specific non-wildcard match field (e.g., a particular source IP or
priority class), the target-specific controller signature fires (Signature S6,
Eq.~\eqref{eq:sig_cpts}). These two checks are applied per-rule and return
immediately on the first match, as their evidence is present in the current
flow table without requiring historical data. If neither instantaneous check
fires, the algorithm retrieves the historical flow-rule log from the blockchain
and applies the autocorrelation test to the time series of malicious flow-rule
counts to detect intermittent attack epochs (Signature S5,
Eq.~\eqref{eq:sig_cpit}). The algorithm returns (\texttt{CP-FR},~$c$),
(\texttt{CP-IT},~$c$), or (\texttt{CP-TS},~$c$) on detection, and
(\texttt{BENIGN},~$c$) if all tests pass.

\begin{algorithm}[H]
\caption{Lightweight Rule-Based Grey Hole Detection — Controller Plane (\texttt{LW-CP-Det})}
\begin{algorithmic}[1]
\small
\Procedure{\texttt{LW-CP-Det}}{$c, t, \tau_c, \gamma_c, W$}
  \State $\mathcal{F}_c \gets \texttt{GET\_FLOW\_RULES}(c, t)$
  \For{each $f \in \mathcal{F}_c$}
    \If{$\texttt{action}(f) = \texttt{DROP}$ \textbf{and} $p_{\text{drop}}(f) > \tau_c$}
      \State \textbf{return} $(\texttt{CP-FR}, c)$
      \Comment{Signature S4}
    \EndIf
    \If{$\texttt{action}(f) = \texttt{DROP}$ \textbf{and} $\texttt{match}(f) \neq \texttt{WILDCARD}$}
      \State \textbf{return} $(\texttt{CP-TS}, c)$
      \Comment{Signature S6}
    \EndIf
  \EndFor
  \State $\mathcal{F}_c^{\text{hist}} \gets \texttt{GET\_BLOCKCHAIN\_LOG}(c, t, W)$
  \State $R_{\mathcal{F}} \gets \texttt{AUTOCORR}(\mathcal{F}_c^{\text{hist}})$
  \If{$\max_T R_{\mathcal{F}}(T) > \gamma_c$ \textbf{and} $\overline{\mathcal{F}_c^{\text{mal}}} > 0$}
    \State \textbf{return} $(\texttt{CP-IT}, c)$
    \Comment{Signature S5}
  \EndIf
  \State \textbf{return} $(\texttt{BENIGN}, c)$
\EndProcedure
\end{algorithmic}
\label{alg:lw_cp_det}
\end{algorithm}





\subsection{Full-Version LLM + Federated Learning Detection Pipeline}

Algorithm~\ref{alg:fv_det} presents the \texttt{FV-Det} full-mode LLM and Federated Learning detection pipeline. As shown in Algorithm~\ref{alg:fv_det},
\texttt{FV-Det} implements the full-mode detection pipeline,
engaged when the lightweight rule-based detectors either do not
fire or produce uncertain results. It takes as input the vehicle
identifier $v_i$, current time $t$, the current global federated model weights
$\mathbf{w}^{(r)}$ from round $r$, and the LLM parameters $\boldsymbol{\theta}$.
The algorithm operates in two interleaved roles: detection and federated model
maintenance.

\begin{algorithm}[!ht]
\caption{Full-Version LLM + Federated Learning Detection Pipeline (\texttt{FV-Det})}
\begin{algorithmic}[1]
\small
\Procedure{\texttt{FV-Det}}{$v_i, t, \mathbf{w}^{(r)}, \boldsymbol{\theta}$}
  \State $\mathbf{x}_i \gets \texttt{TOKENIZE\_LOG}(\texttt{GET\_BLOCKCHAIN\_LOG}(v_i, t, W))$
  \State $\mathcal{Q}_i \gets
    \sum_{k \in \{S1,\ldots,S6\}}
    \mathrm{softmax}\!\bigl(
    \mathrm{LLM}(\mathbf{x}_i;\boldsymbol{\theta})
    \bigr)_{k}$
  \Comment{Eq.~\eqref{eq:llm_score}; 7-class,
           see \S\ref{sec:ai_pipeline_impl}}
  \State $\hat{y}_i^{\text{FL}} \gets f(\mathbf{x}_i^{\text{feat}}; \mathbf{w}^{(r)})$
  \Comment{FL inference: same Qwen2.5-7B base,
           collaborative training role ---
           see \S\ref{sec:ai_pipeline_impl}}
  \State $\mathcal{R}_i \gets \texttt{GET\_REPUTATION}(v_i)$
  \Comment{Eq.~\eqref{eq:reputation}}
  \State $\hat{y}_i \gets \mathbf{1}[\mu_1 \mathcal{S}_{\text{total}} + \mu_2 \mathcal{Q}_i + \mu_3(1-\mathcal{R}_i) > \theta_{\text{det}}]$
  \Comment{Eq.~\eqref{eq:fusion}}
  \If{$\hat{y}_i = 1$}
    \State $\texttt{TRIGGER\_MITIGATION}(v_i, t)$
  \Else
    \State $\Delta\mathbf{w}_i \gets \texttt{LOCAL\_TRAIN}(\mathbf{w}^{(r)}, \mathcal{D}_i)$
    \Comment{Eq.~\eqref{eq:local_loss}}
    \State $\texttt{SUBMIT\_GRADIENT}(\Delta\mathbf{w}_i)$
    \Comment{Blockchain-verified}
  \EndIf
  \State \textbf{return} $\hat{y}_i$
\EndProcedure

\vspace{0.5em}
\Procedure{\texttt{FL-AGGREGATE}}{$\{\Delta\mathbf{w}_i\}_{i=1}^{N}, \mathbf{w}^{(r)}$}
  \State $\mathcal{A} \gets \emptyset$
  \For{each $\Delta\mathbf{w}_i$}
    \If{$\texttt{BLOCKCHAIN\_VERIFY}(\Delta\mathbf{w}_i)$}
      \Comment{Eq.~\eqref{eq:gradient_verify}}
      \State $\mathcal{A} \gets \mathcal{A} \cup \{i\}$
    \EndIf
  \EndFor
  \State $\mathbf{w}^{(r+1)} \gets \sum_{i \in \mathcal{A}} \frac{|\mathcal{D}_i|}{|\mathcal{D}_\mathcal{A}|} \mathbf{w}_i^{(r)}$
  \Comment{Eq.~\eqref{eq:fedavg}}
  \State \textbf{return} $\mathbf{w}^{(r+1)}$
\EndProcedure
\end{algorithmic}
\label{alg:fv_det}
\end{algorithm}

\noindent For detection, the algorithm tokenises the vehicle's blockchain forwarding log
and passes it to the LLM to produce the semantic threat score $\mathcal{Q}_i$
(line~2, Eq.~\eqref{eq:llm_score}), which captures complex temporal and
semantic patterns — such as brief intermittent drop bursts — that evade the
rule-based signatures. Simultaneously, the current FL model produces a binary
classification score $\hat{y}_i^{\text{FL}}$ from the same feature vector
(line~3). The long-term blockchain reputation score $\mathcal{R}_i$ is
retrieved from the ledger (line~4, Eq.~\eqref{eq:reputation}), reflecting the
vehicle's historical behaviour across all past RSU interactions. These three
evidence streams are then fused via the weighted linear decision rule
(line~5, Eq.~\eqref{eq:fusion}): when the weighted sum exceeds the threshold
$\theta_{\text{det}}$, the vehicle is flagged as a grey hole attacker and the
mitigation procedure (\texttt{PQC-Mit}, Algorithm~\ref{alg:pqc_mit}) is
triggered (line~7).

\noindent For model maintenance, vehicles not flagged as attackers contribute to the
federated learning round by training the local model on their labelled
forwarding-log dataset $\mathcal{D}_i$ (line~9, Eq.~\eqref{eq:local_loss})
and submitting the resulting gradient update $\Delta\mathbf{w}_i$ with a
blockchain hash commitment (line~10). The companion \texttt{FL-AGGREGATE}
sub-procedure then collects all submitted updates, rejects any whose received
gradient does not match the pre-committed hash on the blockchain
(Eq.~\eqref{eq:gradient_verify}), and produces the updated global model via
weighted FedAvg (Eq.~\eqref{eq:fedavg}), where each vehicle's contribution is
weighted by its local dataset size to account for heterogeneous observation
volumes across vehicles.

% \begin{figure}[htbp]
% 	\centering
% 	\includegraphics[width=\textwidth,height=0.45\textheight,keepaspectratio]{Figures/figure4_full_version_sequence}
% 	\caption{Full-Mode UML Sequence Diagram}
% 	\label{fig:full-version-sequence}
% \end{figure}

% \vspace{10pt}
\mypar

\noindent Figure~\ref{fig:full-version-sequence} presents the complete UML message sequence for the full-mode detection pipeline, engaged when lightweight signatures do not fire or LLM confidence falls below threshold $\varepsilon_u$. The attacker performs intermittent packet drops (step 1). The legitimate vehicle generates ZKP forwarding proofs (step 2), which are logged to the Blockchain Ledger (step 3). The ledger tokenises the forwarding log and forwards the sequence to the LLM Agent (step 4), which returns a semantic threat score (step 5). In parallel, vehicles submit local FL model updates to the FL Aggregator (step 6), which returns a globally verified model (step 7). Below the dashed detection/mitigation boundary, the Blockchain executes a fusion of LLM score, FL score, and reputation (step 8), after which $k$ partial threshold signatures are collected from RSUs and combined into sigma* (step 9). A threshold-signed FlowMod is applied directly to the OpenFlow switch (step 10), blocking the attacker without involving the SDN controller. PQC-LKH group re-keying is then performed in $\log N$ Kyber broadcast messages (step 11), and the attacker's reputation is updated to zero.

\begin{figure}[H]
	\centering
	\includegraphics[
		width=\textwidth,
		height=0.7\textheight,
		keepaspectratio
	]{Figures/figure4_full_version_sequence}
	\caption{Full-Mode UML Sequence Diagram}
	\label{fig:full-version-sequence}
\end{figure}

\subsection{Post-Quantum Cryptographic Mitigation and Node Isolation}

Algorithm~\ref{alg:pqc_mit} presents the \texttt{PQC-Mit} post-quantum cryptographic mitigation and node isolation procedure. As shown in Algorithm~\ref{alg:pqc_mit},
\texttt{PQC-Mit} executes the complete post-quantum cryptographic
mitigation sequence once a grey hole attacker has been confirmed
by either detection algorithm. It takes as input the suspect
vehicle $v_i$, the SDN controller identifier $c$, the set of $n$ available
RSUs $\{RSU_j\}_{j=1}^{n}$, and the threshold parameter $k$ specifying the
minimum number of RSU co-signatures required to authorise blacklisting. The
algorithm enforces a strict four-step pipeline that prevents any single
compromised component from triggering false isolation or blocking legitimate
isolation.

\begin{algorithm}[!htbp]
\caption{Post-Quantum Cryptographic Mitigation and Node Isolation (\texttt{PQC-Mit})}
\begin{algorithmic}[1]
\small
\Procedure{\texttt{PQC-Mit}}{$v_i, c, \{RSU_j\}_{j=1}^{n}, k$}
  \State \textbf{// Step 1: ZKP Forwarding Proof Verification}
  \State $\pi_i \gets \texttt{GET\_ZKP\_PROOF}(v_i)$
  \Comment{From blockchain log}
  \If{$\texttt{ZKP.Verify}(C_i, \pi_i) = \texttt{FAIL}$}
    \Comment{Eq.~\eqref{eq:zkp_proof}}
    \State $\texttt{FLAG}(v_i, \texttt{CONFIRMED\_ATTACKER})$
  \EndIf
  \State \textbf{// Step 2: Threshold-Signed Blacklisting}
  \State $\mathcal{B}(v_i) \gets \texttt{CREATE\_BLACKLIST\_MSG}(v_i, t)$
  \For{$j = 1$ \textbf{to} $k$}
    \State $\sigma_j \gets \text{Dilithium.Sign}(\mathbf{sk}_j, \mathcal{B}(v_i))$
    \Comment{Eq.~\eqref{eq:threshold_partial}}
  \EndFor
  \State $\sigma^* \gets \texttt{TS.Combine}(\{\sigma_j\}_{j=1}^{k})$
  \Comment{Eq.~\eqref{eq:threshold_combine}}
  \If{$\texttt{TS.Verify}(\mathbf{pk}_{\text{group}}, \mathcal{B}(v_i), \sigma^*)$}
    \Comment{Eq.~\eqref{eq:threshold_verify}}
% --- SUPERVISOR REVISION BLOCK 12: Algorithm 4 Steps 3-4 with threshold key operations ---
    \State \textbf{// Step 3: Controller + RSU Threshold-Signed
      Flow Modification}
    \State $\mathcal{M}_{\text{iso}} \gets
      \texttt{CREATE\_FLOWMOD}(v_i, \texttt{BLOCK})$
    \For{$j = 1$ \textbf{to} $k_{\text{key}}-1$}
      \State $\sigma_j^{\text{iso}} \gets
        \mathrm{Dilithium.Sign}
        (\mathbf{sk}_j,\, \mathcal{M}_{\text{iso}})$
      \Comment{RSU $j$ partial signature}
    \EndFor
    \State $\sigma_c^{\text{iso}} \gets
      \mathrm{Dilithium.Sign}
      (\mathbf{sk}_c,\, \mathcal{M}_{\text{iso}})$
    \Comment{Controller partial signature (one share)}
    \State $\sigma_{\text{iso}}^* \gets
      \mathrm{TS.Combine}\!\left(
      \sigma_c^{\text{iso}},\,
      \left\{\sigma_j^{\text{iso}}\right\}_{j=1}^{k_{\text{key}}-1}
      \right)$
    \Comment{Eq.~\eqref{eq:rekey_sig}}
    \State $\texttt{BROADCAST\_FLOWMOD}
      (\mathcal{M}_{\text{iso}},\, \sigma_{\text{iso}}^*)$

    \State \textbf{// Step 4: PQC-LKH Group Re-Keying via
      DKG Threshold Broadcast}
    \State $\mathrm{path} \gets
      \texttt{GET\_PATH}(\ell_i, \mathrm{root}, \mathcal{T})$
    \For{each $u \in \mathrm{path}$}
      \State $(K_u^{\mathrm{new}}, c_u) \gets
        \mathrm{Kyber.Enc}
        (\mathbf{pk}_u^{\mathrm{sib}},\,
        k_u^{\mathrm{new}})$
      \Comment{Eq.~\eqref{eq:lkh_rekey}}
      \State $\sigma_u^* \gets
        \texttt{THRESHOLD-SIGN-REKEY}
        (c_u,\, \mathbf{sk}_c,\,
        \{\mathbf{sk}_j\}_{j=1}^{k_{\text{key}}-1})$
      \Comment{Eq.~\eqref{eq:rekey_sig}: controller + RSUs}
      \State $\texttt{BROADCAST}(c_u,\, \sigma_u^*)$
    \EndFor
    \State $(K_{\text{grp}},\, c_{\mathrm{root}}) \gets
      \mathrm{Kyber.Enc}
      (\mathbf{pk}_{\mathrm{root}}^{\mathrm{new}},\, m)$
    \Comment{Eq.~\eqref{eq:lkh_grp}}
% --- END SUPERVISOR REVISION BLOCK 12 ---
  
  
  
  \EndIf




\EndProcedure
\end{algorithmic}
\label{alg:pqc_mit}
\end{algorithm}

\textbf{Step 1 — ZKP proof verification (lines 2--4):} The algorithm retrieves
the suspect vehicle's ZKP forwarding proof $\pi_i$ from the blockchain
(Eq.~\eqref{eq:zkp_proof}) and verifies it against the Pedersen commitment
$C_i$. A verification failure confirms that the node dropped packets it
received — this is the cryptographic gate that distinguishes malicious dropping
from legitimate channel loss, because a vehicle that merely lost packets due to
channel conditions will still hold valid proofs for the packets it did forward.

\textbf{Step 2 — Threshold-signed blacklisting (lines 6--12):} A blacklist
message $\mathcal{B}(v_i)$ is constructed and circulated to $k$ independent
RSUs. Each RSU $j$ produces a Dilithium partial signature over the message
(Eq.~\eqref{eq:threshold_partial}); the $k$ partial signatures are combined
into a single compact aggregate signature $\sigma^*$
(Eq.~\eqref{eq:threshold_combine}), and verified against the group public key
(Eq.~\eqref{eq:threshold_verify}). This co-signing requirement ensures that
isolation cannot proceed unless $k$ independent infrastructure nodes concur,
preventing any single compromised RSU from triggering false isolation.

\textbf{Step 3 — Threshold-signed flow modification (lines 13--17):}
On a valid $\sigma^*$, the controller contributes one Dilithium partial
signature and $k_{\text{key}}-1$ RSUs each contribute a partial signature
over the flow modification command $\mathcal{M}_{\text{iso}}$. These are
combined into the aggregate signature $\sigma_{\text{iso}}^*$
(Eq.~\eqref{eq:rekey_sig}), which is broadcast to the affected OpenFlow
switch. The switch verifies $\sigma_{\text{iso}}^*$ against the group
public key before installing the block rule, ensuring that neither the
controller alone nor any single RSU can unilaterally issue a valid
isolation command.

\textbf{Step 4 — PQC-LKH group re-keying (lines 16--22):} The algorithm
traverses the path from the isolated vehicle's leaf node $\ell_i$ to the root
of the PQC-LKH binary tree. For each node $u$ on this path, a new Kyber key
pair is generated and the updated key is re-encrypted under the sibling
subtree's existing public key (Eq.~\eqref{eq:lkh_rekey}), signed with
Dilithium, and broadcast. Finally, a new group session key $K_{\text{grp}}$ is
encapsulated under the refreshed root key (Eq.~\eqref{eq:lkh_grp}). This
requires exactly $\lceil \log_2 N \rceil$ Kyber operations — as opposed to
$O(N)$ for naïve unicast re-keying — and cryptographically excludes the
isolated vehicle without requiring direct contact with it, since it no longer
holds valid private keys for any refreshed node on the path.

% --- SUPERVISOR REVISION BLOCK 13: Algorithm 5 controller failover and key revocation ---
\subsection{Controller Failover and Key Revocation}
\label{sec:alg_failover}

Algorithm~\ref{alg:controller_failover} presents the
\texttt{Controller-Failover} procedure, which is triggered by
the DEBSC when $\mathcal{T}_c(t) < \theta_c$. As shown in
Algorithm~\ref{alg:controller_failover}, the procedure iterates
the segment-specific blockchain-stored ordered list to find the
first eligible replacement controller, revokes the failed
controller's signing key, and re-establishes flow rules through
the whitelist-governed installation process. If no eligible
replacement exists for a segment, degraded mode is activated
(Equation~\eqref{eq:degraded_mode}) to provide partial
containment until a qualified controller becomes available.
The ordered-list traversal (rather than highest-trust selection)
ensures that failover load distributes across the controller
fleet and that no single controller accumulates all fallback
sessions.

\begin{algorithm}[H]
\caption{Controller Failover and Key Revocation
(\texttt{Controller-Failover})}
\begin{algorithmic}[1]
\small
\Procedure{\texttt{Controller-Failover}}
  {$c_{\text{failed}},\;
    \{[c_j^z]\}_z,\;
    \theta_c$}
  \State $\mathcal{Z} \gets
    \texttt{GET\_ACTIVE\_SEGMENTS}(c_{\text{failed}})$
  \For{each segment $z \in \mathcal{Z}$}
    \State $\mathrm{selected} \gets \texttt{false}$
    \For{each $c_j^z$ in ordered list for segment $z$}
      \If{$c_j^z \neq c_{\text{failed}}$
          \textbf{and}
          $\mathcal{T}_{c_j^z}(t) > \theta_c$}
        \Comment{Eq.~\eqref{eq:failover_select}}
        \State $c_{\text{new}}^z \gets c_j^z$
        \State $\mathrm{selected} \gets \texttt{true}$
        \State \textbf{break}
      \EndIf
    \EndFor
    \If{\textbf{not} $\mathrm{selected}$}
      \State $\texttt{ACTIVATE\_DEGRADED\_MODE}(z)$
      \Comment{Eq.~\eqref{eq:degraded_mode}: no eligible
        replacement}
    \Else
      \State $\texttt{BLOCKCHAIN\_REVOKE\_KEY}
        (\mathbf{pk}_{c_{\text{failed}}})$
      \Comment{Via RSU consensus}
      \State $\texttt{NOTIFY\_SWITCHES}
        (z,\;
        \mathbf{pk}_{c_{\text{failed}}},\;
        \texttt{REVOKED})$
      \State $\texttt{REINSTALL\_FLOW\_RULES}
        (c_{\text{new}}^z,\; z,\;
        \mathcal{W}_{\text{BC}})$
      \Comment{Whitelist-governed; no unverified rules}
      \State $\texttt{REFRESH\_DKG\_SHARE}
        (c_{\text{new}}^z,\;
        \{RSU_j\})$
      \Comment{Eq.~\eqref{eq:dkg_share}: new controller
        joins DKG}
    \EndIf
  \EndFor
\EndProcedure
\end{algorithmic}
\label{alg:controller_failover}
\end{algorithm}
% --- END SUPERVISOR REVISION BLOCK 13 ---

% --- SUPERVISOR REVISION BLOCK 14: Algorithm 6 ZKP cross-check with RSU observation ---
\subsection{ZKP Cross-Check with RSU Observation}
\label{sec:alg_zkp_crosscheck}

Algorithm~\ref{alg:zkp_crosscheck} presents the
\texttt{ZKP-Crosscheck} procedure, invoked by both the lightweight
and full-mode detectors at each observation window. As shown in
Algorithm~\ref{alg:zkp_crosscheck}, the procedure enforces the
three-state ZKP model: it first checks proof presence within
deadline $T_{\text{zkp}}$ (returning ABSENT if missing), then
verifies cryptographic validity, and finally cross-references the
committed forwarding count against the RSU-independently-observed
count within tolerance $\varepsilon_{\text{obs}}$. A proof that
passes cryptographic verification but fails the cross-reference
is classified as FAIL, closing the fabricated commitment attack
vector. The RSU observation $\widehat{n}_i^{\text{fwd}}$ is
collected at the RSU coverage boundary independently of any data
provided by vehicle $v_i$, ensuring the cross-reference cannot
be manipulated by the vehicle under evaluation.

\begin{algorithm}[H]
\caption{ZKP Cross-Check with RSU Observation
(\texttt{ZKP-Crosscheck})}
\begin{algorithmic}[1]
\small
\Procedure{\texttt{ZKP-Crosscheck}}
  {$v_i,\; t,\; T_{\text{zkp}},\;
    \varepsilon_{\text{obs}}$}
  \State $\pi_i,\; t_{\text{recv}} \gets
    \texttt{GET\_PROOF\_FROM\_BLOCKCHAIN}(v_i, t)$
  \If{$\pi_i = \texttt{null}$
      \textbf{or}
      $(t_{\text{recv}} - t) > T_{\text{zkp}}$}
    \State \textbf{return} \texttt{ABSENT}
    \Comment{Eq.~\eqref{eq:zkp_state}: proof absent or late}
  \EndIf
  \If{\textbf{not}
      $\texttt{ZKP.Verify}(C_i,\, \pi_i)$}
    \Comment{Eq.~\eqref{eq:zkp_proof}}
    \State \textbf{return} \texttt{FAIL}
    \Comment{Cryptographic verification failed}
  \EndIf
  \State $\widehat{n}_i \gets
    \texttt{GET\_RSU\_OBSERVED\_COUNT}(v_i, t)$
  \Comment{RSU-independently observed count}
  \If{$\left|n_i^{\text{fwd}} - \widehat{n}_i\right|
      > \varepsilon_{\text{obs}}$}
    \Comment{Eq.~\eqref{eq:rsu_crossref}}
    \State \textbf{return} \texttt{FAIL}
    \Comment{Commitment inconsistent with RSU observation}
  \EndIf
  \State \textbf{return} \texttt{PASS}
\EndProcedure
\end{algorithmic}
\label{alg:zkp_crosscheck}
\end{algorithm}
% --- END SUPERVISOR REVISION BLOCK 14 ---


\begin{figure}[H]
	\centering
		\includegraphics[
		width=\textwidth,
		height=0.7\textheight,
		keepaspectratio
	]{Figures/figure5_cryptographic_mitigation.png}
	\caption{Cryptographic Mitigation Flowchart}
	\label{fig:crypto-mitigation-flow}
\end{figure}

Figure~\ref{fig:crypto-mitigation-flow} illustrates the cryptographic decision pipeline for post-quantum mitigation and node isolation. As shown in Figure~\ref{fig:crypto-mitigation-flow}, the cryptographic decision pipeline executes following a malicious detection verdict from either the lightweight or full-mode detector. The first gate is the ZKP forwarding proof: if the suspect node's proof retrieved from the blockchain is valid, the statistical detection is classified as a false positive and the sequence terminates with no action — this prevents false isolation of legitimate nodes that experience high packet loss due to mobility or channel conditions. If the ZKP proof fails, the node is flagged as a confirmed attacker. The second gate requires collecting $k$ partial threshold signatures from independent RSUs, which are combined into an aggregate signature sigma*: if insufficient RSUs concur, the process aborts, preventing single-RSU false isolation. On a valid sigma*, RSUs broadcast a Dilithium-signed FlowMod directly to the affected OpenFlow switch, installing a block rule without routing the command through the potentially compromised SDN controller. Finally, PQC-LKH group re-keying broadcasts exactly $\log N$ Kyber-encrypted key updates along the refresh path, ensuring the isolated attacker cannot derive any updated group session key without requiring unicast contact with every remaining vehicle.

\section{Cryptographic Analysis}
\label{sec:crypto_analysis}

\subsection{Why Cryptography Alone Cannot Mitigate Grey Hole Attacks}

The fundamental property of a grey hole attack is that the attacker is a
\emph{legitimate participant} in the network. It receives authenticated packets
through properly established sessions, holds valid certificates, and produces
correct message authentication codes. Classical cryptographic mechanisms --- HMAC,
ECDSA, TLS --- verify that a packet was not forged or tampered with in transit.
They say nothing about what the receiving node does with the packet after
authentication succeeds.

Formally, let $v_i$ receive packet $p$ authenticated under key $k_i$. Classical
authentication confirms $\texttt{HMAC}(p, k_i) = \texttt{tag}$, establishing
that $p$ arrived intact and from the claimed sender. However, the post-receipt
forwarding decision $\mathbf{1}[v_i \text{ forwards } p]$ is entirely within
$v_i$'s control and is \emph{not constrained} by any authentication mechanism.
A grey hole attacker simply sets this decision to zero for selected packets.
No amount of stronger authentication --- including post-quantum signatures --- changes
this: authentication secures the channel, not the forwarding behaviour of
the channel endpoint.

This is the precise reason why grey hole attacks are harder to detect than
black hole or wormhole attacks: the attacker's traffic is cryptographically
valid at every point it does appear, making statistical anomaly detection the
only viable primary signal.

\mypar
\textbf{Consequence for system design.} Cryptography must be repurposed from
channel authentication to \emph{behaviour attestation}. Instead of proving
that a packet was sent correctly, the cryptographic layer must prove that a
received packet was subsequently forwarded --- or expose the absence of that proof
as a detection signal. This requires a two-phase cryptographic deployment:

\textbf{Phase 1 --- Pre-detection (evidence generation).}
Cryptography generates unforgeable forwarding proofs that serve as input evidence to the detection engine. A grey hole attacker that drops packets \emph{cannot} generate valid proofs for those packets, creating a verifiable discrepancy on the blockchain ledger. This transforms silent dropping into a detectable cryptographic failure without requiring the network to observe the drop directly.

\textbf{Phase 2 --- Post-detection (mitigation enforcement).}
Once the detection engine issues a malicious verdict, cryptography enforces isolation: post-quantum-signed flow modification commands block the attacker's traffic paths; threshold signatures require multi-RSU consensus before blacklisting, preventing single-point false isolation; and group re-keying cryptographically excludes the attacker from future session keys, making continued participation impossible even if the node rejoins under the same identity.

This two-phase design ensures that cryptographic mechanisms contribute to both
detection \emph{and} mitigation, rather than leaving detection entirely to
statistical methods that a sophisticated attacker can evade by calibrating drop
rates below detection thresholds.

\subsection{Selected Cryptographic Techniques by Phase}

The candidate set considered is: \{homomorphic encryption, attribute-based
encryption (ABE), broadcast encryption, homomorphic signatures, ring signatures,
threshold signatures, ZKP, multi-party computation (MPC)\}, with post-quantum
cryptography (PQC) mandatory for quantum-resistance. Techniques are selected
and assigned to phases as follows.

\mypar
\textbf{Phase 1 --- Pre-Detection Evidence Generation}

\textbf{(1) Zero-Knowledge Proofs (ZKP):}
ZKP is the sole technique capable of proving forwarding behaviour without
revealing packet content. Each vehicle $v_i$ commits to its forwarded packet
count via a Pedersen commitment $C_i = g^{n_i^{\text{fwd}}} \cdot h^{r_i}$
(Eq.~\eqref{eq:pedersen}) and generates a proof $\pi_i$
(Eq.~\eqref{eq:zkp_proof}) demonstrating knowledge of a valid opening
consistent with the independently observable blockchain receipt counts. A grey
hole attacker that drops packets cannot produce $\pi_i$ for the dropped batch
because $n_i^{\text{fwd}}$ in its commitment will not match the observable
receipt count on the blockchain --- the discrepancy is the detection signal fed
into the DEBSC gate (Eq.~\eqref{eq:debsc}). ZKP is selected over homomorphic
signatures because it is binding, hiding, and compatible with PQC lattice
constructions, while homomorphic signatures require linear-size proof structures
that scale poorly with batch size.

\textbf{(2) Attribute-Based Encryption (ABE) --- Full mode, target-specific attacks:}
In target-specific attacks, the attacker selectively drops packets from emergency
vehicles while forwarding normal traffic. ABE contributes to detection by
encrypting emergency traffic under the policy \texttt{priority = EMERGENCY}: only
nodes holding valid attribute credentials can decrypt and forward emergency packets.
An attacker without the emergency-traffic credential cannot perform
\emph{content-aware} selective dropping, forcing it to either drop all traffic
(detectable as fixed-rate attack) or forward emergency traffic (attack fails).
ABE is pre-detection because it degrades the attacker's targeting capability,
making its behaviour detectable by the standard PDR signatures S1--S3.

\mypar
\textbf{Phase 2 --- Post-Detection Mitigation Enforcement}

\textbf{(3) Threshold Signatures --- Blacklisting authorisation:}
A $(k,n)$-threshold signature scheme (Eqs.~\eqref{eq:threshold_partial}--\eqref{eq:threshold_verify})
requires $k$ of $n$ RSUs to co-sign the blacklisting decision before the SDN
controller executes isolation. This is critical in high-mobility SDVN where any
single RSU has a short observation window and limited statistical confidence.
Threshold signatures prevent a single compromised or mistaken RSU from unilaterally
isolating a legitimate vehicle. Selected over ring signatures because ring
signatures do not provide individual accountability --- a requirement for auditable
blacklisting in a regulated vehicular network.

\textbf{(4) CRYSTALS-Dilithium --- Flow modification authentication:}
Dilithium (EUF-CMA secure under Module-LWE) authenticates all SDN controller flow
modification commands (Eqs.~\eqref{eq:dilithium_sign}--\eqref{eq:dilithium_verify}).
An OpenFlow switch receiving an isolation command verifies the Dilithium signature
before installing the block rule, preventing a compromised or spoofed controller
from injecting false isolation commands during the re-association window identified
in Section~\ref{sec:mobility_analysis}.

\textbf{(5) CRYSTALS-Kyber (KEM) --- Group re-keying:}
Kyber (IND-CCA2 secure under Module-LWE) encapsulates new group session keys after
isolation (Eqs.~\eqref{eq:kyber_enc}--\eqref{eq:kyber_dec}). The PQC-LKH tree
structure (Section~\ref{sec:pqc_lkh}) reduces re-keying cost from $O(N)$ to
$O(\log N)$ KEM operations. ML-KEM-1024 (Kyber-1024, NIST
security level~5) is used in both lightweight and full modes,
satisfying the mandatory post-quantum security requirement
throughout the SHIELD-GH framework.

\mypar
\textbf{Techniques not selected:} Homomorphic encryption and MPC impose
prohibitive computational overhead on vehicle OBU hardware and are excluded.
Broadcast encryption does not address the forwarding proof requirement. Homomorphic
signatures duplicate ZKP functionality at higher proof size and without the
hiding property needed to protect packet content privacy.

\section{Data Collection and Simulation Environment}

The limited real world datasets to use in vehicular Grey Hole attack makes this study use a controlled data generation process in a simulated environment. The area of data collection is centered on the dynamic interactions between the high-mobility vehicular nodes and the multi-layered security framework.


\subsection{Data Generation Process}

The simulation environment is set to simulate high-mobility vehicular situations and the standard communication protocols needed by Software-Defined Vehicular Networks. To produce a robust dataset for training the detection components, the environment injects three grey hole attack variants as programmed malicious behaviours.
\textbf{Selective dropping data} is generated by introducing traffic flows with varying packet-loss percentages to model fixed-rate attacks.
\textbf{Temporal discrepancy data} is produced by simulating intermittent malicious states, capturing time-based attack patterns.
\textbf{Semantic traffic data} is created by assigning priority to specific packet types such as emergency broadcasts, emulating target-specific intelligent attacks.


\subsection{Functional Feature Extraction}

Raw simulation logs are analyzed to capture the features required by the decentralized security modules.
\textbf{Behavioral logs} record packet transmission and reception outcomes to enable node integrity tracking over time.
\textbf{Network contextual features} capture temporal and semantic metadata — including packet timestamps and priority headers — to support identification of complex behavioural patterns.
\textbf{Stochastic model gradients} represent locally computed model updates extracted from each vehicle for use in collaborative learning without exposing raw traffic data.


\begin{comment}
\section{Performance Evaluation}

To prove the efficiency of the suggested multi-layered security framework, it is compared to three state-of-the-art (SOTA) techniques that are the representatives of various paradigms of the Grey Hole detection in the vehicular networks. The choice of these comparison methods is due to their recent publication and the relevance to individual elements of our framework (Neural Networks, Blockchain, and Machine Learning):

\begin{itemize}
	\item \textbf{ANN-ABC based IDS (2023):} A smart Intrusion Detection System, based on an Artificial Neural Network (ANN) optimized using the Artificial Bee Colony (ABC) algorithm to detect packet-dropping patterns in AODV-based vehicular networks.
	
	\item \textbf{VCBC Smart Contract Protocol (2024):} A security mechanism that is based on the use of Voting-Classification Blockchain-Based Contracts (VCBC) to categorize the nodes as white, black, or grey according to their ratios of packet delivery that is documented on a public ledger (based on blockchain).
	
	\item \textbf{RF-SVM Anomaly Detection (2023):} An enhanced Machine Learning-based IDS utilizing a Random Forest (RF) classifier to identify subtle Grey Hole signatures through an analysis of 20 network traffic features derived out of simulation logs.
\end{itemize}


\section{Performance Evaluation Metrics}

The proposed and baseline techniques are tested with the help of a set of quantitative metrics that will measure the security performance and the operational performance on the network. The metrics undertaken are the following:

\subsection{Network Performance Metrics}

\begin{itemize}
	\item \textbf{Packet Delivery Ratio (PDR):} The proportion of the data packets that have been effectively sent to the target destination nodes relative to the total amount of packets sent. This is the main measure of the capability of the framework to counter the attacks of packet-dropping.
	
	\item \textbf{Throughput:} The rate of successful data transmission over the network, measured in bits per second (bps) or bytes per second (Bps). This evaluates the overall network efficiency under attack conditions.
	
	\item \textbf{End-to-End (E2E) Delay:} The average time taken for a packet to travel from the source vehicle to the destination node. T This measure evaluates the delay caused by the multi-layered analysis (Blockchain consensus and LLM inference).
\end{itemize}

\subsection{Detection Reliability Metrics}

\begin{itemize}
	\item \textbf{Detection Accuracy:} The proportion of the number of correctly identified malicious and legitimate behaviors to the total number of test samples, which gives a general picture of the model performance.
	\item \textbf{False Positive Rate (FPR):} The percentage of legitimate nodes wrongly identified as malicious to the overall legitimate number of nodes. This is an important metric in SDVNs to make sure that security measures do not introduce an unwarranted fragmentation to the network.
	\item \textbf{Matthews Correlation Coefficient (MCC):} An effective statistical rate that yields a high score can only be achieved when the detection results are good in all the four categories of the confusion matrix (True Positives, True Negatives, False Positives, and False Negatives). This is used to assess the performance of the LLM and FL modules, namely in unbalanced datasets where malicious nodes are not common as the legitimate nodes.
\end{itemize}
\end{comment}






\section{Simulation Settings}
\label{sec:sim_settings}

This section specifies the complete simulation configuration of the proposed
SHIELD-GH framework, covering the mobility, radio, blockchain, cryptography and
detection pipelines, together with the rationale for each chosen value, since
these values drive the results reported in Section~\ref{sec:results_progress}.
Every parameter is either taken from the literature, backed by an explicit
engineering justification, or fixed by a sensitivity sweep
(Section~\ref{subsec:gh_sensitivity}). Following standard practice, the internal
hyper-parameters of the state-of-the-art baselines are not re-specified here;
each baseline is reproduced with the configuration reported in its original
paper (Section~\ref{sec:evaluation}).

\subsection{Urban Mobility Map and Scenario}

We selected a dense urban district of Galle, Sri Lanka (centre
$6.0535^{\circ}$\,N, $80.2210^{\circ}$\,E) for the simulation. This district is
a dense street network with many overlapping and crossing roads, producing a
high density of intersections and therefore a worst-case multipath/contention
environment for an SDVN. Its mobility traces were imported using the SUMO and
OpenStreetMap (OSM) tools over a $2\,\text{km}\times2\,\text{km}$ region
measured directly in OSM; the OSM extract was cropped to that exact bounding
box and processed with \texttt{netconvert}, yielding a SUMO network spanning
$2097\,\text{m}\times2120\,\text{m}$ with $284$ real intersections. The network
was then simulated using Network Simulator~3 (version~3.35). The
$2\,\text{km}\times2\,\text{km}$ extent is chosen because it is tiled exactly by
the RSU grid deployed in Section~\ref{subsec:rsu_deployment}.

The SUMO scenario runs for $600$\,s and injects $200$ vehicles of five types
($100$ cars, $25$ buses, $25$ lorries, $25$ vans and $25$ trucks), a
light/heavy mix that approximates a real urban fleet so that mobility (speed,
acceleration) is heterogeneous rather than uniform. The maximum speed of every
type is capped at $150$\,km/h ($41.67$\,m/s); the effective speed is governed by
the OSM road-class limits, giving an observed mean network speed of
$\approx21$\,km/h with peaks above $30$\,km/h, which is representative of
dense-city driving. All $200$ vehicles are injected within the first $30$\,s
(warm-up), and each vehicle carries SUMO's auto-rerouting device together with
a city-wide rerouter, so that vehicles are continuously handed new destinations
and keep driving for the whole run rather than arriving and leaving the network.
As a result, at least $190$ vehicles are concurrently present for
$\approx94\%$ of the measurement window (observed mean $\approx195$ concurrent
vehicles). Because SUMO requires this warm-up to inject every vehicle, the NS-3
evaluation begins at $t=30$\,s, once all $200$ vehicles are present.

Vehicle positions, speeds and headings are exported at $10$\,Hz via SUMO's FCD
output and (i)~converted to an NS-2 mobility trace consumed by NS-3's
\texttt{Ns2MobilityHelper}, and (ii)~streamed to the SUMO--NS-3 TraCI bridge
(\texttt{sumo\_ns3\_bridge.py}), which updates NS-3 node positions each step and
writes per-vehicle speed logs for the MATD module (Eq.~\eqref{eq:matd}).
Figure~\ref{fig:sim_environment} shows the resulting simulation environment,
combining the SUMO microscopic simulation with the imported NetAnim mobility
trace.

\begin{figure}[H]
    \centering
    \begin{subfigure}[b]{0.49\textwidth}
        \centering
        \includegraphics[width=\textwidth,keepaspectratio]{Screenshots/galle_sumo.png}
        \caption{SUMO traffic simulation.}
        \label{fig:sumo_sim}
    \end{subfigure}
    \hfill
    \begin{subfigure}[b]{0.49\textwidth}
        \centering
        \includegraphics[width=\textwidth,keepaspectratio]{Screenshots/galle_netanim.png}
        \caption{NetAnim mobility trace.}
        \label{fig:netanim_sim}
    \end{subfigure}
    \caption{Simulation environment on the Galle map
        ($2\,\text{km}\times2\,\text{km}$): (a)~the SUMO microscopic simulation
        of $200$ vehicles, and (b)~the NetAnim view of the imported mobility
        trace with the $64$ ($8\times8$) RSUs overlaid.}
    \label{fig:sim_environment}
\end{figure}

\subsection{Network, Radio and RSU Deployment}
\label{subsec:rsu_deployment}

The network layer is simulated in NS-3~3.35. The vehicular data plane uses the
IEEE~802.11p (WAVE/DSRC) PHY/MAC at $5.9$\,GHz with a log-distance
propagation-loss model. We place $64$ RSUs on an $8\times8$ grid at $250$\,m
spacing (horizontally and vertically) overlaid on the mobility trace, giving a
$264$-node data plane ($200$ vehicles $+$ $64$ RSUs) in addition to the SDN
control plane. Under FCC regulations governing the DSRC service
(47~CFR~\S\,90.377~\cite{fcc2004dsrc}), the control channel is permitted a
maximum effective radiated power of $44$\,dBm EIRP. We set the transmission
power to $41$\,dBm, which stays within this regulatory ceiling and, under the
log-distance path-loss model, yields an $\approx270$\,m communication range;
this range is then used as the coverage radius $R_{\text{RSU}}$ input to the
MATD link-lifetime / handoff model (Eq.~\eqref{eq:handoff_loss}), so that
coverage is obtained by setting the transmit power rather than by
over-provisioning RSUs. The standard DSRC value for a shared multi-channel
deployment is $33$\,dBm, but because we operate a single (control) DSRC
channel under the higher $44$\,dBm ceiling, the $41$\,dBm setting is used, and
the $250$\,m grid spacing then gives continuous coverage at that range.

\subsection{Blockchain, Consensus and Cryptography}

The blockchain trust layer is realised in two consistent forms. During
prototyping it runs as an in-NS-3 in-memory ledger (\texttt{BlockchainLedger},
C++/OpenSSL SHA-256) that faithfully replicates the DEBSC dual-evidence gate
(Eq.~\eqref{eq:debsc}) and the reputation computation (Eq.~\eqref{eq:reputation})
without live-node overhead. For the permissioned-ledger experiments, each NS-3
node additionally acts as a Hyperledger Fabric client: a bridge replays the
node's on-chain events to a real Fabric test network over gRPC, where dynamic
VRF-based endorser selection (Eq.~\eqref{eq:endorser_select}) chooses a
per-transaction endorser set from the $64$ RSUs. Trust and reputation are
computed from forwarding records committed each $1$\,s observation slot
(Eqs.~\eqref{eq:trust_update}--\eqref{eq:reputation}). All cryptographic
material (hashes, signatures, ZKP proofs) is serialized to hexadecimal strings
and transmitted as bytes; the post-quantum primitives follow the NIST standards
(ML-KEM, FIPS~203~\cite{nist2024fips203}; ML-DSA, FIPS~204~\cite{nist2024fips204}).

\subsection{SHIELD-GH Detection Parameters}
\label{subsec:gh_detection_params}

The MATD module uses coverage radius $R_{\text{RSU}}=270$\,m (set by the
$41$\,dBm transmit power above), average handoff duration
$\Delta t_{ho}=0.5$\,s, worst-case handoff loss $\rho_{\max}=0.30$, and mobility
decay $\lambda_s=0.01$~\cite{islam2021sdvn}. The observation sliding window is
$W=10$ slots. The signature thresholds --- fixed-rate PDR threshold
$\tau_f=0.60$ and variance bound $\varepsilon_f=0.20$ (S1); intermittent
per-slot threshold $\tau_{it}=0.70$ and autocorrelation threshold
$\gamma_{it}=1.30$ (S2); target-specific KL-divergence threshold
$\tau_{ts}=0.50$ (S3); reputation isolation threshold $\theta_R=0.40$ --- are
each fixed by the sensitivity sweep of Section~\ref{subsec:gh_sensitivity}
rather than chosen arbitrarily. The DEBSC gate requires both the statistical
reputation gate and the ZKP gate to fail simultaneously before isolation
(Eq.~\eqref{eq:debsc}), and active mode needs a minimum blockchain history of
$N_{\text{min}}=5$ RSU interactions.

Table~\ref{tab:sim_settings} lists every proposed-method parameter together with
its value and the justification for that value, so that no setting is chosen
arbitrarily.

{\small
\setlength{\LTleft}{0pt}\setlength{\LTright}{\fill}
\begin{longtable}{| p{3.0cm} | p{4.3cm} | p{6.0cm} |}
\caption{Simulation settings and the justification for each value.}
\label{tab:sim_settings}\\
\hline
\textbf{Parameter} & \textbf{Value} & \textbf{Justification} \\
\hline
\endfirsthead
\multicolumn{3}{l}{\itshape Table~\ref{tab:sim_settings} (continued)}\\
\hline
\textbf{Parameter} & \textbf{Value} & \textbf{Justification} \\
\hline
\endhead
\hline
\endfoot
\hline
\endlastfoot
Network simulation tool & Network Simulator version 3.35 & Discrete-event simulator that is the de-facto standard for VANET/SDVN research. \\
\hline
Mobility trace extraction & SUMO, OpenStreetMap & Established open-source microscopic traffic simulator for deriving realistic mobility from real road maps~\cite{rathod2017survey}. \\
\hline
Map & Galle, Sri Lanka (dense urban; $2\,\text{km}\times2\,\text{km}$ measured, net $2097\times2120$\,m, $284$ intersections) & Dense grid with many crossing/overlapping roads gives a worst-case multipath/contention environment; the $2\,\text{km}^2$ extent matches the RSU grid footprint (see RSU deployment row). \\
\hline
Run time & SUMO trace $600$\,s; NS-3 measurement $300$\,s per data point (excl.\ $30$\,s warm-up) & Supervisor spec ($\geq300$\,s per point); the trace runs to $600$\,s so the $300$\,s window is fully warmed up. \\
\hline
NS-3 start time & $t=30$\,s & SUMO needs $\approx30$\,s to inject all $200$ vehicles; measurement begins only once all are present. \\
\hline
SUMO step length & $0.1$\,s ($6000$ steps / $600$\,s) & Sub-second step resolves per-vehicle mobility and matches the $10$\,Hz FCD export. \\
\hline
Number of nodes & $264$ (200 vehicles, 64 RSUs) plus SDN controllers & Vehicle count sized to a dense urban scenario; RSU count follows from the grid deployment (see RSU deployment row). \\
\hline
Vehicle composition & 100 cars, 25 buses, 25 lorries, 25 vans, 25 trucks & Light/heavy mix approximating real urban fleet composition, so mobility is heterogeneous rather than uniform. \\
\hline
Maximum vehicle speed & $150$\,km/h cap (observed mean $\approx21$\,km/h) & Cap set from the SUMO vehicle-class definitions; effective speed governed by OSM road limits, keeping it a realistic vehicular network ($>60$\,km/h attainable). \\
\hline
Concurrent vehicles & $\approx195$ (mean); $\geq190$ for $\sim94\%$ of the run & Auto-rerouting device $+$ city-wide rerouter keep $\sim200$ vehicles continuously in-network for the full run. \\
\hline
RSU deployment & $8\times8$ grid, $250$\,m spacing, $R=270$\,m coverage & Uniform grid overlay; $250$\,m gives continuous coverage at the $270$\,m radio range without over-provisioning RSUs. \\
\hline
Transmission power & $41$\,dBm ($\Rightarrow \approx270$\,m range) & Within the $44$\,dBm EIRP ceiling set by FCC 47~CFR~\S\,90.377 for the DSRC control channel~\cite{fcc2004dsrc}; resulting $270$\,m range then feeds the link-lifetime model (standard multi-channel value is $33$\,dBm). \\
\hline
Data-plane channel / propagation & $5.9$\,GHz DSRC IEEE~802.11p; log-distance loss & DSRC is the standard V2X data-plane radio; log-distance suits the urban macrocell path loss. \\
\hline
Location update frequency & $10$\,Hz (TraCI FCD export) & Cooperative-awareness beacon rate used for V2V safety messaging. \\
\hline
Blockchain (prototype / full) & In-memory C++ ledger (OpenSSL SHA-256) / Hyperledger Fabric via gRPC bridge, dynamic VRF endorsers & In-memory ledger replicates DEBSC logic cheaply during prototyping; real Fabric used for the permissioned-ledger experiments. \\
\hline
Consensus & rPBFT / BFT endorser set (Eq.~\eqref{eq:kend}) & Byzantine fault tolerance for multi-RSU verdict agreement; $k_{\text{end}}$ guarantees $\geq1$ honest endorser. \\
\hline
Post-quantum crypto & ML-KEM (FIPS~203), ML-DSA (FIPS~204), SHA-256 & NIST-standardized Level-5 PQC primitives, so the long-lived vehicular PKI resists future quantum adversaries~\cite{nist2024fips203,nist2024fips204}. \\
\hline
LLM base model & Qwen2.5-7B-Instruct (7.6\,B parameters) &
Satisfies the $>$4\,B, $<$15\,B moderately-large model requirement;
selected per model evaluation report. \\
\hline
LLM task head & 7-class sequence classification: BENIGN $+$ S1--S6 &
Aligns directly with the six-variant attack taxonomy; per-variant class
probabilities support M1 (MCC) and M3 (AVCR) evaluation without a
separate classification head. \\
\hline
LLM fine-tune method & LoRA/PEFT; target modules:
\texttt{q\_proj}, \texttt{k\_proj}, \texttt{v\_proj},
\texttt{o\_proj}; dropout 0.05 &
Federated-friendly: small $\Delta\mathbf{w}$ reduces gradient
commitment size and communication overhead
(Eq.~\eqref{eq:gradient_commit}). \\
\hline
LoRA rank & $r_{\text{edge}}=8$ (RSU inference),
$r_{\text{cloud}}=16$ (training / cloud tier) &
Lower rank for real-time OBU-class inference; higher rank for
training benchmark and cloud-tier re-evaluation. \\
\hline
LLM quantisation & 4-bit NF4, double quantisation, bf16 compute &
Fits RSU hardware memory footprint; consistent with OBU feasibility
discussion in Section~\ref{sec:lit_pqc_cav}. \\
\hline
LLM optimiser & AdamW, lr $= 2\times10^{-4}$, 3 local epochs
per FL round &
Standard LoRA training configuration; 3 epochs matches the FL
round budget. \\
\hline
LLM max token length ($n_e$) & 96 tokens &
Covers $W=10$ observation window at approximately 9 tokens per slot
(\texttt{FWD}/\texttt{DRP}/\texttt{HOF}) with one
\texttt{RULE} prefix token for controller-plane variants. \\
\hline
Edge escalation threshold $\varepsilon_u$ & 0.15 &
Validated on held-out split; below this maximum class confidence
the input is escalated to the cloud LLM tier
(Eq.~\eqref{eq:llm_routing}). \\
\hline
FL rounds & 5 &
Configurable; 5 rounds $\times$ 3 local LoRA epochs per round
produces the reported validation MCC. \\
\hline
Training hardware & NVIDIA RTX 5090 (32\,GB), CUDA 12.8,
PyTorch 2.11 &
Global LoRA training environment; RSU-side inference uses the
quantised $r=8$ edge configuration. \\
\hline
Attack model & Grey hole; default drop rate $\rho_a = 40\%$,
default attacker penetration $p = 40\%$; both swept
in Experiment~E1 (Eq.~\eqref{eq:e1_grid}) &
Default values match the fixed parameters of Experiments
E2--E5; $\rho_a = 40\%$ sits above the MATD noise floor
(Eq.~\eqref{eq:handoff_loss}) and below the S1 saturation
point, making it the canonical grey hole operating point
for controlled comparison. \\
\hline
Observation window $W$ & $10$ slots ($10$\,s) & Long enough to estimate per-flow PDR statistics, short enough for responsive detection (sweep, Table~\ref{tab:gh_sensitivity}). \\
\hline
MATD $\Delta t_{ho},\ \rho_{\max},\ \lambda_s$ & $0.5$\,s, $0.30$, $0.01$ & Handoff and mobility-decay constants from the SDVN mobility model~\cite{islam2021sdvn}. \\
\hline
S1 ($\tau_f,\ \varepsilon_f$) & $0.60,\ 0.20$ & Selected by grid search on validation traces (Table~\ref{tab:gh_sensitivity}) for best MCC. \\
\hline
S2 ($\tau_{it},\ \gamma_{it},\ [T_{\min},T_{\max}]$) & $0.70,\ 1.30,\ [3,30]$\,s & Thresholds from the sweep; the $[3,30]$\,s band spans plausible intermittent-attack periods. \\
\hline
S3 ($\tau_{ts}$) & $0.50$ & KL-divergence cut selected by the sweep to separate target-specific dropping from benign variation. \\
\hline
Reputation threshold $\theta_R$ & $0.40$ & Sweep-selected isolation cut balancing detection rate against false isolation. \\
\hline
Min.\ history $N_{\text{min}}$ & $5$ RSU interactions & Below this, reputation estimates are statistically insufficient to authorise isolation. \\
\hline
ZKP deadline $T_{\text{zkp}}$ & $2$\,s per window & Bounds proof latency to the observation cadence. \\
\hline
\end{longtable}
}

\subsection{Sensitivity Analysis and Threshold Selection}
\label{subsec:gh_sensitivity}

The detection thresholds are not set arbitrarily: each is chosen by a grid
search that inspects a candidate set of values on held-out validation traces and
retains the value maximising the Matthews Correlation Coefficient (MCC), the
balanced metric of Section~\ref{sec:evaluation}. Table~\ref{tab:gh_sensitivity}
lists the searched ranges and the retained values. The observation window $W$ is
swept jointly because it trades detection latency against the stability of the
per-flow PDR estimate. The retained operating point is the mid-grid,
literature-consistent configuration; the observed-MCC column is populated as the
full defence-on sweep is completed.

\begin{table}[H]
    \centering
    \small
    \caption{SHIELD-GH detection-threshold grid search. For each parameter the
        retained value maximises validation MCC; entries marked $\dagger$ are
        finalised when the full defence-on sweep is run.}
    \label{tab:gh_sensitivity}
    \begin{tabular}{| m{3.8cm} | m{5.2cm} | m{2.8cm} |}
        \hline
        \textbf{Parameter} & \textbf{Values searched (grid)} & \textbf{Selected} \\
        \hline
        Observation window $W$ (slots) & $\{5,\ 10,\ 20\}$ & $10$ \\
        \hline
        $\tau_f$ (S1 PDR threshold) & $\{0.50,\ 0.60,\ 0.70\}$ & $0.60$ \\
        \hline
        $\varepsilon_f$ (S1 variance bound) & $\{0.10,\ 0.20,\ 0.30\}$ & $0.20$ \\
        \hline
        $\tau_{it}$ (S2 per-slot threshold) & $\{0.60,\ 0.70,\ 0.80\}$ & $0.70$ \\
        \hline
        $\gamma_{it}$ (S2 autocorrelation) & $\{1.10,\ 1.30,\ 1.50\}$ & $1.30$ \\
        \hline
        $\tau_{ts}$ (S3 KL-divergence) & $\{0.30,\ 0.50,\ 0.70\}$ & $0.50$ \\
        \hline
        $\theta_R$ (reputation isolation) & $\{0.30,\ 0.40,\ 0.50\}$ & $0.40$ \\
        \hline
        Validation MCC at operating point & --- & $^{\dagger}$next phase \\
        \hline
    \end{tabular}
\end{table}










\chapter{Results}
\label{chap:results}

This chapter is divided into three parts. Section~\ref{sec:evaluation}
presents the performance evaluation methodology, comparing SHIELD-GH against
state-of-the-art baselines and reporting the internal ablation study and
benchmark experiment results. Section~\ref{sec:results_progress} presents
early simulation results obtained from the ns-3.35 prototype, demonstrating
the framework's behaviour under a fixed-rate grey hole attack.
Section~\ref{sec:expected_results} then presents the full set of expected
performance results based on analytical modelling of the six attack variants.

%% ============================================================
\section{Implemented AI Pipeline and Training Configuration}
\label{sec:ai_pipeline_impl}
%% ============================================================

This section describes the concrete AI pipeline as implemented and its end-to-end
data flow, realising the abstract full-mode detection logic of
Algorithm~\ref{alg:fv_det}, so that the reported results can be reproduced
exactly. The pipeline transforms raw blockchain forwarding records into a fused
detection verdict through five stages, all sharing the single selected language
model.

\mypar
\textbf{Stage~1 --- Blockchain log tokenisation.} Each vehicle's forwarding
record over the observation window $W = 10$ slots (Table~\ref{tab:sim_settings})
is rendered as a token sequence. Every slot emits one \texttt{ACTION:src} token,
where the action is \texttt{FWD} (forwarded), \texttt{DRP} (dropped), or
\texttt{HOF} (handoff slot, topologically-expected benign loss), and \texttt{src}
is the coarse source bucket of the packet batch. Controller-plane variants
prepend a \texttt{RULE} token describing the installed flow rule. This yields the
tokenised window $\mathbf{x}_i^{(t)}$ of Eq.~\eqref{eq:llm_score}, capped at
$n_e = 96$ tokens.

\mypar
\textbf{Stage~2 --- LLM semantic scoring.} The tokenised window is classified by
a \textbf{Qwen2.5-7B-Instruct} model fine-tuned with a LoRA adapter
(per the Edge-LLM architecture of Chapter~\ref{sec:methodology}) into the seven
classes $\{\mathrm{BENIGN}, \mathrm{S1}, \ldots, \mathrm{S6}\}$. The semantic
threat score $\mathcal{Q}_i(t)$ is the total softmax mass on the six malicious
classes (Eq.~\eqref{eq:llm_score}). When the maximum class confidence falls below
$\varepsilon_u = 0.15$, the input is escalated from the edge configuration
($r_{\text{edge}}=8$) to the higher-capacity cloud configuration
($r_{\text{cloud}}=16$) per the routing rule of Eq.~\eqref{eq:llm_routing}.

\mypar
\textbf{Stage~3 --- Federated training of the LLM.} The same Qwen2.5-7B model is
the federated learning model: each vehicle $v_i$ locally fine-tunes its LoRA
adapter on its private dataset $\mathcal{D}_i$ by minimising the cross-entropy
objective of Eq.~\eqref{eq:local_loss}, and transmits only the small adapter
update $\Delta\mathbf{w}_i$. Before transmission it commits the hash
$\mathcal{C}_i = \mathtt{Hash}(\Delta\mathbf{w}_i \Vert t \Vert \mathrm{id}_i)$ to
the blockchain (Eq.~\eqref{eq:gradient_commit}); the aggregator accepts an update
only if the received gradient's hash matches the on-chain commitment
(Eq.~\eqref{eq:gradient_verify}), and forms the new global model by dataset-size
weighted FedAvg (Eq.~\eqref{eq:fedavg}). The LLM and the FL model are therefore a
\emph{single} Qwen2.5-7B model, differentiated only by role: semantic scoring
(inference) versus collaborative training (federated fine-tuning).

\mypar
\textbf{Stage~4 --- Evidence fusion.} The LLM score $\mathcal{Q}_i(t)$ is combined
with the rule-based aggregate signature score $\mathcal{S}_{\mathrm{total}}(v_i)$
and the blockchain reputation deficit $(1 - \mathcal{R}_i(t))$ through the
validation-tuned linear rule of Eq.~\eqref{eq:fusion}, producing the binary
verdict $\hat{y}_i(t)$.

\mypar
\textbf{Stage~5 --- Verdict handoff.} A positive verdict is passed to the DEBSC
isolation gate (Eq.~\eqref{eq:debsc}); detection never triggers isolation on
statistical evidence alone.

\mypar
\textbf{Training configuration.} The LoRA adapter targets the attention
projections \texttt{q\_proj}, \texttt{k\_proj}, \texttt{v\_proj},
\texttt{o\_proj} with dropout $0.05$, optimised by AdamW at learning rate
$2\times10^{-4}$ for three local epochs per federated round over
five global FL rounds, batch size $16$, and
maximum sequence length $96$ tokens; every setting appears with its justification
in Table~\ref{tab:sim_settings}. The base model is loaded in 4-bit NF4
quantisation for deployment-footprint parity with resource-constrained road-side
units, consistent with the OBU feasibility discussion of
Section~\ref{sec:lit_pqc_cav}. All settings were fixed prior to the reported runs
and are held constant across the evaluation.

%% ============================================================
\section{Performance Evaluation}
\label{sec:evaluation}
%% ============================================================

\subsection{External Baselines}

Three post-2020 baselines targeting grey hole attacks are selected,
each representing a distinct methodological paradigm to isolate
the contribution of the proposed framework.

\textbf{B1}~\cite{malik2023greyhole} is an ML-only dynamic threshold
detector for grey hole attacks in NS-3-simulated VANETs — the same
simulator used here — enabling direct metric comparison. It carries
no blockchain, mobility correction, or cryptographic component.

\textbf{B2}~\cite{alabdulatif2024mitigating} employs VCBC blockchain
smart contracts for grey hole mitigation, achieving PDR of 78\% under
attack. It represents the blockchain-only paradigm with no AI detection
and no PQC.

\textbf{B3}~\cite{ahsan2024flbert} applies FL-BERT to grey hole detection
in software-defined VANET — the closest existing architecture to the
proposed full mode. It has no mobility correction, no controller-plane
model, and no cryptographic mitigation.

% --- SUPERVISOR REVISION BLOCK 16: Performance Evaluation Metrics ---
\subsection{Performance Evaluation Metrics}
\label{sec:metrics}

Twelve mutually exclusive metrics are defined to evaluate all formulated
design components of the SHIELD-GH framework. Metrics M1--M6 are
state-of-the-art comparable: they can be measured for external baselines
B1--B3 as well as for the proposed framework, enabling direct comparison.
Metrics M7--M12 are ablation-only: they evaluate design components absent
from all reviewed baselines and address aspects reviewers will question
regarding novel contributions. Together the twelve metrics cover detection
quality, attack impact reversal, false isolation safety, system response
timing, multi-dimensional overhead and scalability, ZKP evasion handling,
controller anomaly detection, re-keying complexity, FL poisoning robustness,
controller failover timing, and endorser pool resilience. Jitter and load
balance, reported in preliminary results tables, are removed from formal
evaluation as they are not causally linked to any formulated design component.

\vspace{0.6em}
\noindent\textbf{State-of-the-Art Comparable Metrics (M1--M6)}

\mypar
\textbf{M1 --- Attack Detection Matthews Correlation Coefficient.}
MCC is selected as the single classification quality metric in preference
to accuracy and FPR separately, as it accounts for all four confusion
matrix cells under the class imbalance inherent in vehicular attack
scenarios — malicious nodes constitute a strict minority of the vehicle
population. MCC is evaluated per attack variant $k \in \{S1,\ldots,S6\}$
and in aggregate across all variants, and is compared against baselines B1,
B3, and ablations A4 (LLM scorer disabled, $\mu_2 = 0$) and A5 (FL replaced
by centralised static model). The per-variant metric is given in
Equation~\eqref{eq:m1_mcc}.

\begin{equation}
\label{eq:m1_mcc}
\text{MCC}_k =
\frac{TP_k \cdot TN_k - FP_k \cdot FN_k}
{\sqrt{(TP_k + FP_k)(TP_k + FN_k)(TN_k + FP_k)(TN_k + FN_k)}}
\end{equation}

As given in Equation~\eqref{eq:m1_mcc}, $TP_k$, $TN_k$, $FP_k$, and
$FN_k$ are the true positive, true negative, false positive, and false
negative counts for attack variant $k$ over the full simulation run.
An aggregate MCC is computed by summing the confusion matrix cells across
all six variants before applying the formula. Evaluates: signatures S1--S6
(Eqs.~\eqref{eq:sig_dpfr}--\eqref{eq:sig_cpts}), LLM-FL fusion
(Eq.~\eqref{eq:fusion}), and aggregate controller anomaly score $\Psi_c(t)$
(Eq.~\eqref{eq:controller_agg}).

\mypar
\textbf{M2 --- Grey Hole Suppression Ratio (GHSR).}
GHSR measures the fraction of attack-induced PDR loss recovered after
detection and isolation completes, normalised against the attack impact
magnitude. Unlike raw PDR, GHSR is independent of attack intensity,
enabling fair comparison across different drop rates and attack variants.
The metric is given in Equation~\eqref{eq:m2_ghsr}.

\begin{equation}
\label{eq:m2_ghsr}
\text{GHSR} =
\frac{\overline{\text{PDR}}_{\text{post}} -
      \overline{\text{PDR}}_{\text{attack}}}
     {\overline{\text{PDR}}_{\text{baseline}} -
      \overline{\text{PDR}}_{\text{attack}}}
\end{equation}

As given in Equation~\eqref{eq:m2_ghsr},
$\overline{\text{PDR}}_{\text{baseline}}$ is the network-wide mean PDR
under no attack, $\overline{\text{PDR}}_{\text{attack}}$ is the mean PDR
during the attack epoch before mitigation, and
$\overline{\text{PDR}}_{\text{post}}$ is the mean PDR after isolation
completes. GHSR $= 1$ indicates complete PDR recovery; GHSR $= 0$ indicates
no recovery. Evaluates the full mitigation pipeline: DEBSC isolation gate
(Eq.~\eqref{eq:debsc}), threshold-signed FlowMod (Eq.~\eqref{eq:rekey_sig}),
and PQC-LKH re-keying (Eq.~\eqref{eq:lkh_rekey}). Compared against B2
(PDR 78\% reference point).

\mypar
\textbf{M3 --- Attack Variant Coverage Rate (AVCR).}
AVCR measures the proportion of the six formally defined attack variants
S1--S6 for which the framework achieves a per-variant true positive rate
above a minimum coverage threshold $\theta_{\text{cov}}$. Aggregate MCC
(M1) can mask variant-specific detection blindspots; AVCR exposes them.
The metric is given in Equation~\eqref{eq:m3_avcr}.

\begin{equation}
\label{eq:m3_avcr}
\text{AVCR} =
\frac{1}{6}
\sum_{k=1}^{6}
\mathbf{1}\!\left[
\frac{TP_k}{TP_k + FN_k} \geq \theta_{\text{cov}}
\right]
\end{equation}

As given in Equation~\eqref{eq:m3_avcr}, $TP_k / (TP_k + FN_k)$ is the
per-variant true positive rate and $\theta_{\text{cov}}$ is the minimum
acceptable detection rate for a variant to be counted as covered. AVCR $= 1$
indicates all six variants are detected above threshold; values below $1$
identify which variants remain underdetected. Evaluates all six attack
signatures Eqs.~\eqref{eq:sig_dpfr}--\eqref{eq:sig_cpts} individually.
Compared against B1 and B3, which cover at most two and three variants
respectively.

\mypar
\textbf{M4 --- False Isolation Rate (FIR).}
FIR is the fraction of legitimate vehicles incorrectly isolated during
the simulation run, measured directly from NS-3 isolation event logs as a
simulation-observable discrete count. It is mutually exclusive from M1
because it is evaluated exclusively on the confirmed-legitimate vehicle
subpopulation at varying vehicle densities and attacker ratios. The metric
is given in Equation~\eqref{eq:m4_fir}.

\begin{equation}
\label{eq:m4_fir}
\text{FIR} =
\frac{\left|\left\{
v_i \in \mathcal{V}_{\text{legit}} :
\mathrm{Isolate}(v_i) = 1
\right\}\right|}
{|\mathcal{V}_{\text{legit}}|}
\end{equation}

As given in Equation~\eqref{eq:m4_fir}, $\mathcal{V}_{\text{legit}}$ is
the set of all legitimate (non-attacking) vehicles in the simulation and
$\mathrm{Isolate}(v_i) = 1$ is the event that vehicle $v_i$ was incorrectly
isolated by the DEBSC during the measurement window. Evaluates: MATD
correction (Eq.~\eqref{eq:matd}), DEBSC dual-evidence gate
(Eq.~\eqref{eq:debsc}), and probationary mode (Eq.~\eqref{eq:vehicle_mode}).
Compared against B1 and B3 which lack MATD, directly revealing the
false-isolation reduction attributable to mobility correction.

\mypar
\textbf{M5 --- End-to-End Security Response Latency (ESRL).}
ESRL measures the total elapsed time from attack onset to completed node
isolation, evaluated separately for lightweight and full detection modes
and decomposed by pipeline stage. The aggregate metric is given in
Equation~\eqref{eq:m5_esrl} and the stage decomposition in
Equation~\eqref{eq:m5_esrl_decomp}.

\begin{equation}
\label{eq:m5_esrl}
\text{ESRL} = t_{\text{isolate}} - t_{\text{onset}}
\end{equation}

As given in Equation~\eqref{eq:m5_esrl}, $t_{\text{onset}}$ is the
simulation timestamp at which the attacker begins dropping packets and
$t_{\text{isolate}}$ is the timestamp at which the DEBSC FlowMod takes
effect at the OpenFlow switch. The stage decomposition is given in
Equation~\eqref{eq:m5_esrl_decomp}.

\begin{equation}
\label{eq:m5_esrl_decomp}
\text{ESRL} =
\underbrace{\Delta t_{\text{det}}}_{\text{Signature/LLM}}
\;+\;
\underbrace{\Delta t_{\text{ZKP}}}_{\text{ZKP verify}}
\;+\;
\underbrace{\Delta t_{\text{sig}}}_{\text{Threshold sign}}
\;+\;
\underbrace{\Delta t_{\text{FM}}}_{\text{FlowMod install}}
\end{equation}

As given in Equation~\eqref{eq:m5_esrl_decomp}, the four stages correspond
to: detection algorithm execution time
(Algorithms~\ref{alg:lw_dp_det}--\ref{alg:fv_det}); ZKP cross-check
latency (Algorithm~\ref{alg:zkp_crosscheck}); VRF endorser selection and
threshold co-signing latency
(Eqs.~\eqref{eq:vrf_eval}--\eqref{eq:endorser_select},
\eqref{eq:threshold_combine}); and FlowMod propagation and installation
time. PQC overhead relative to classical cryptography is isolated via
ablation A5.

\mypar
\textbf{M6 --- Multi-Dimensional Protocol Overhead and Scalability (MDPOS).}
MDPOS quantifies the additional resource burden introduced by SHIELD-GH
protocol components across three orthogonal resource dimensions — computation,
communication, and storage — each measured as a function of vehicle population
size $N$ and RSU count to provide a scalability profile. This single
composite metric addresses the scalability concerns reviewers will raise for
a PQC-heavy framework. Three sub-metrics are defined in
Equations~\eqref{eq:m6_comp}, \eqref{eq:m6_comm},
and~\eqref{eq:m6_store}.

\begin{equation}
\label{eq:m6_comp}
\Omega_{\text{comp}}(N) =
\sum_{op \in \mathcal{O}}
\bar{c}_{op} \cdot \bar{f}_{op}(N)
\end{equation}

As given in Equation~\eqref{eq:m6_comp},
$\mathcal{O} = \{\texttt{ZKP.Prove},\, \texttt{ZKP.Verify},\,
\texttt{VRF},\, \texttt{Dil.Sign},\, \texttt{Dil.Verify},\,
\texttt{Kyber.Enc},\, \texttt{Kyber.Dec},\, \texttt{DKG.Share}\}$
is the set of SHIELD-GH cryptographic operations, $\bar{c}_{op}$ is the
mean per-operation CPU time measured on OBU-class hardware, and
$\bar{f}_{op}(N)$ is the mean invocation frequency per vehicle per second
as a function of $N$. Evaluates: Eqs.~\eqref{eq:zkp_proof},
\eqref{eq:vrf_eval}, \eqref{eq:dilithium_sign}, \eqref{eq:kyber_enc},
\eqref{eq:dkg_share}.

\begin{equation}
\label{eq:m6_comm}
\Omega_{\text{comm}}(N) =
\frac{1}{N} \sum_{i=1}^{N} \bar{B}_i(N)
\end{equation}

As given in Equation~\eqref{eq:m6_comm}, $\bar{B}_i(N)$ is the mean
additional bytes per second per vehicle $v_i$ attributable to SHIELD-GH
protocol messages — ZKP proof submissions (Eq.~\eqref{eq:zkp_proof}),
VRF evaluation broadcasts (Eq.~\eqref{eq:vrf_eval}), threshold partial
signature messages (Eq.~\eqref{eq:threshold_partial}), and PQC-LKH
re-keying ciphertexts (Eq.~\eqref{eq:lkh_rekey}) — measured as a
function of increasing $N$.

\begin{equation}
\label{eq:m6_store}
\Omega_{\text{store}}(N, t) =
N \cdot
\left(
b_{\text{fwd}} + b_{\text{ZKP}} + b_{\text{grad}} + b_{\text{VRF}}
\right)
\cdot \frac{t}{W}
\end{equation}

As given in Equation~\eqref{eq:m6_store}, $b_{\text{fwd}}$,
$b_{\text{ZKP}}$, $b_{\text{grad}}$, and $b_{\text{VRF}}$ are the
per-vehicle per-window byte sizes of the forwarding record, ZKP
commitment (Eq.~\eqref{eq:pedersen}), gradient hash commitment
(Eq.~\eqref{eq:gradient_commit}), and VRF seed record
(Eq.~\eqref{eq:vrf_seed}) respectively, $t$ is the elapsed simulation
time, and $W$ is the observation window length. The total gives the
blockchain ledger growth rate as a function of $N$ and operating time.

\vspace{0.6em}
\noindent\textbf{Ablation-Only Metrics (M7--M12)}

These metrics evaluate design components absent from all three external
baselines B1--B3 and therefore cannot be fairly compared against existing
works. Each is evaluated through the ablation variants A1--A5 or
purpose-built controlled experiments within the NS-3 simulation.

\mypar
\textbf{M7 --- ZKP Proof Evasion Detection Rate (PEDR).}
PEDR measures the fraction of proof-evasion events — covering both
cryptographic failure (FAIL state) and deliberate proof withholding
(ABSENT state) — correctly flagged by the three-state ZKP model,
including fabricated commitments caught by RSU cross-reference and
single-strike ZKP violations in probationary mode. This metric
specifically evaluates the ABSENT-state extension and probationary
mode single-strike rule, both absent from any baseline. The metric
is given in Equation~\eqref{eq:m7_pedr}.

\begin{equation}
\label{eq:m7_pedr}
\text{PEDR} =
\frac{\left|\left\{
v_i \in \mathcal{V}_{\text{att}} :
\Pi_{\text{ZKP}}(v_i, t)
\in \{\texttt{FAIL},\,\texttt{ABSENT}\}
\right\}\right|}
{|\mathcal{V}_{\text{att}}|}
\end{equation}

As given in Equation~\eqref{eq:m7_pedr}, $\mathcal{V}_{\text{att}}$ is
the set of attacking vehicles that either submitted a false proof or
withheld a proof within $T_{\text{zkp}}$, and $\Pi_{\text{ZKP}}(v_i,t)$
is the three-state ZKP result from Eq.~\eqref{eq:zkp_state}. Evaluates:
three-state ZKP model (Eq.~\eqref{eq:zkp_state}), RSU cross-reference
(Eq.~\eqref{eq:rsu_crossref}), and probationary mode
(Eq.~\eqref{eq:vehicle_mode}). Ablation: three-state vs two-state ZKP
(A3 variant).

\mypar
\textbf{M8 --- Aggregate Controller Anomaly Detection Rate (CADR).}
CADR measures the rate at which low-and-slow controller attacks — where
no individual rule triggers S4--S6 but $\Psi_c(t)$ accumulates sufficient
sub-threshold evidence — are correctly detected and result in a controller
trust decrement. This metric directly evaluates the aggregate anomaly
score, which has no equivalent in any baseline. The metric is given in
Equation~\eqref{eq:m8_cadr}.

\begin{equation}
\label{eq:m8_cadr}
\text{CADR} =
\frac{\left|\left\{
c \in \mathcal{C}_{\text{mal}} :
\Psi_c(t) > \psi_{\text{thresh}}
\right\}\right|}
{|\mathcal{C}_{\text{mal}}|}
\end{equation}

As given in Equation~\eqref{eq:m8_cadr}, $\mathcal{C}_{\text{mal}}$ is
the set of compromised controllers conducting low-and-slow sub-threshold
attacks and $\Psi_c(t)$ is the aggregate behavioural anomaly score from
Eq.~\eqref{eq:controller_agg}. Evaluates:
Eqs.~\eqref{eq:controller_agg} and~\eqref{eq:controller_trust}. Ablation:
aggregate anomaly scoring enabled vs disabled.

\mypar
\textbf{M9 --- PQC-LKH Re-Keying Operation Efficiency Ratio.}
This metric empirically validates the theoretical $O(\log N)$ re-keying
complexity claim of the PQC-LKH scheme against the $O(N)$ unicast Kyber
re-keying baseline, measuring the ratio of KEM operations as $N$ scales.
The metric is given in Equation~\eqref{eq:m9_rekey}.

\begin{equation}
\label{eq:m9_rekey}
\eta_{\text{rekey}}(N) =
\frac{\lceil \log_2 N \rceil}{N - 1}
\end{equation}

As given in Equation~\eqref{eq:m9_rekey}, $\lceil \log_2 N \rceil$ is
the number of Kyber operations required by PQC-LKH
(Eq.~\eqref{eq:lkh_rekey}) and $N-1$ is the number required by unicast
re-keying for the same exclusion guarantee. As $N$ grows,
$\eta_{\text{rekey}}(N) \to 0$, confirming the logarithmic efficiency
gain. Evaluates: Eqs.~\eqref{eq:lkh_rekey}, \eqref{eq:lkh_grp}, and
the lazy membership policy (Eq.~\eqref{eq:lkh_keys}). Ablation:
PQC-LKH vs unicast Kyber re-keying across increasing $N$.

\mypar
\textbf{M10 --- FL Global Model Robustness under Gradient Poisoning.}
This metric evaluates the MCC of the global FL detection model as the
fraction $\rho_p$ of vehicles submitting poisoned gradients increases
from zero to the BFT bound $f_{\max}/n$, measuring the contribution of
the blockchain gradient integrity mechanism in maintaining model quality
under poisoning pressure. The metric is given in
Equation~\eqref{eq:m10_fl}.

\begin{equation}
\label{eq:m10_fl}
\text{MCC}_{\text{FL}}(\rho_p) =
\frac{
TP(\rho_p) \cdot TN(\rho_p) -
FP(\rho_p) \cdot FN(\rho_p)
}{\sqrt{
(TP(\rho_p)+FP(\rho_p))
(TP(\rho_p)+FN(\rho_p))
(TN(\rho_p)+FP(\rho_p))
(TN(\rho_p)+FN(\rho_p))
}}
\end{equation}

As given in Equation~\eqref{eq:m10_fl}, $\rho_p \in [0,\, f_{\max}/n]$
is the fraction of vehicles submitting poisoned gradient updates, and
$TP(\rho_p)$, $TN(\rho_p)$, $FP(\rho_p)$, $FN(\rho_p)$ are the global
model confusion matrix entries at poisoning fraction $\rho_p$. Evaluates:
gradient hash commitment (Eq.~\eqref{eq:gradient_commit}), acceptance
filter (Eq.~\eqref{eq:gradient_verify}), and weighted FedAvg
(Eq.~\eqref{eq:fedavg}). Ablation A4: blockchain-verified FedAvg vs
centralised static model.

\mypar
\textbf{M11 --- Controller Failover Recovery Time (CFRT).}
CFRT measures the elapsed time from the DEBSC detecting
$\mathcal{T}_c(t) < \theta_c$ to full segment re-routing operational
under the replacement controller, including key revocation, DKG share
refresh, and whitelist-governed flow rule reinstallation. The aggregate
metric is given in Equation~\eqref{eq:m11_cfrt} and the stage
decomposition in Equation~\eqref{eq:m11_cfrt_decomp}.

\begin{equation}
\label{eq:m11_cfrt}
\text{CFRT} = t_{\text{restored}} - t_{\text{trigger}}
\end{equation}

As given in Equation~\eqref{eq:m11_cfrt}, $t_{\text{trigger}}$ is
the timestamp at which $\mathcal{T}_c(t)$ falls below $\theta_c$ and
the DEBSC initiates failover, and $t_{\text{restored}}$ is the timestamp
at which the replacement controller's first whitelist-verified flow rule
takes effect in segment $z$. The stage decomposition is given in
Equation~\eqref{eq:m11_cfrt_decomp}.

\begin{equation}
\label{eq:m11_cfrt_decomp}
\text{CFRT} =
\underbrace{\Delta t_{\text{revoke}}}_{\text{Key revocation}}
\;+\;
\underbrace{\Delta t_{\text{notify}}}_{\text{Switch notification}}
\;+\;
\underbrace{\Delta t_{\text{DKG}}}_{\text{DKG share refresh}}
\;+\;
\underbrace{\Delta t_{\text{reinstall}}}_{\text{Flow rule reinstall}}
\end{equation}

As given in Equation~\eqref{eq:m11_cfrt_decomp}, the four stages
correspond to: blockchain key revocation via RSU consensus
(Algorithm~\ref{alg:controller_failover}, line~16); switch notification
of the revoked public key (line~17); DKG share refresh for the
replacement controller (Eq.~\eqref{eq:dkg_share}); and
whitelist-governed flow rule reinstallation (line~18). Evaluates:
Algorithm~\ref{alg:controller_failover},
Eqs.~\eqref{eq:failover_select}, \eqref{eq:controller_trust},
and~\eqref{eq:dkg_share}.

\mypar
\textbf{M12 --- VRF Endorser Pool Availability Rate (EPAR).}
EPAR measures the probability of each EndorserMode being activated as
the fraction of compromised RSUs $f_c$ increases from zero to beyond the
BFT bound, validating that the bypass mechanism activates at the correct
thresholds and that \texttt{NORMAL} mode is maintained with high
probability within the T4 adversary bound. The metric is given in
Equation~\eqref{eq:m12_epar}.

\begin{equation}
\label{eq:m12_epar}
P_{\text{mode}}(f_c) =
\Pr\!\left[
\mathrm{EndorserMode}(t) = \text{mode}
\;\middle|\;
|\mathcal{R}_{\text{comp}}| =
\left\lfloor f_c \cdot n \right\rfloor
\right]
\end{equation}

As given in Equation~\eqref{eq:m12_epar}, $\mathcal{R}_{\text{comp}}$
is the set of compromised RSUs, $f_c \in [0,1]$ is the RSU compromise
fraction, $n$ is the total RSU count, and the probability is computed
empirically over repeated simulation trials at each $f_c$ level.
Evaluated for each mode $\in \{\texttt{NORMAL},\,\texttt{RELAXED},\,
\texttt{DEFERRED},\,\texttt{EMERGENCY}\}$ as $f_c$ increases.
Evaluates: Eqs.~\eqref{eq:endorser_pool}, \eqref{eq:kend},
and~\eqref{eq:endorser_mode}. Ablation: VRF dynamic endorser selection
vs static first-$k$-to-respond scheme.

\vspace{0.6em}
\noindent\textbf{Metric-to-Component Coverage Summary}

\begin{table}[H]
\renewcommand{\arraystretch}{1.25}
\caption{Performance Evaluation Metric Coverage of Formulated Design Components}
\label{tab:metric_coverage}
\centering
\scriptsize
\begin{tabular}{|p{3.2cm}|p{0.55cm}|p{0.55cm}|p{0.55cm}|p{0.55cm}|p{0.55cm}|p{0.55cm}|p{0.55cm}|p{0.55cm}|p{0.55cm}|p{0.55cm}|p{0.55cm}|p{0.55cm}|}
\hline
\textbf{Design Component} &
\textbf{M1} & \textbf{M2} & \textbf{M3} & \textbf{M4} &
\textbf{M5} & \textbf{M6} & \textbf{M7} & \textbf{M8} &
\textbf{M9} & \textbf{M10} & \textbf{M11} & \textbf{M12} \\
\hline
Signatures S1--S6 & \checkmark & & \checkmark & & & & & & & & & \\
\hline
MATD + PDR correction & & & & \checkmark & & & & & & & & \\
\hline
DEBSC dual-evidence gate & & \checkmark & & \checkmark & \checkmark & & & & & & & \\
\hline
Three-state ZKP + cross-ref & & & & & & & \checkmark & & & & & \\
\hline
Probationary mode & & & & & & & \checkmark & & & & & \\
\hline
LLM-FL fusion & \checkmark & & & & & & & & & & & \\
\hline
Gradient integrity + FedAvg & & & & & & & & & & \checkmark & & \\
\hline
$\Psi_c$ aggregate anomaly & \checkmark & & & & & & & \checkmark & & & & \\
\hline
Controller trust + failover & & & & & & & & & & & \checkmark & \\
\hline
DKG key generation & & & & & & \checkmark & & & & & \checkmark & \\
\hline
PQC-LKH re-keying & & \checkmark & & & & \checkmark & & & \checkmark & & & \\
\hline
VRF endorser selection & & & & & & & & & & & & \checkmark \\
\hline
End-to-end response pipeline & & & & & \checkmark & & & & & & & \\
\hline
Protocol overhead + scalability & & & & & & \checkmark & & & & & & \\
\hline
\textbf{Type} &
\multicolumn{6}{c|}{\textit{State-of-art comparable}} &
\multicolumn{6}{c|}{\textit{Ablation-only}} \\
\hline
\end{tabular}
\end{table}
% --- END SUPERVISOR REVISION BLOCK 16 ---

% --- SUPERVISOR REVISION BLOCK 17: Upgraded Ablation Study ---
\subsection{Internal Ablation Study}
\label{sec:ablation}

The original five-ablation design is retired and replaced with seventeen
component-specific ablations covering every individually formulated design
element of SHIELD-GH. Each ablation isolates exactly one component by
removing or replacing it with a degraded alternative while holding all
other components constant. The independent variable ($x$) is selected per
ablation to maximally expose the contribution of the specific component
under test; a single $x$ per ablation is used to conserve experimental
space. All applicable metrics from the finalised set M1--M12 are
evaluated per ablation. Ablations A1--A3 are state-of-the-art comparable:
baselines B1--B3 can be run in the same experimental setup for direct
comparison. Ablations A4--A17 are ablation-only, as no reviewed baseline
implements these components. Table~\ref{tab:ablation_summary} provides an
overview; detailed experimental designs follow.

\begin{table}[!htbp]
\renewcommand{\arraystretch}{1.2}
\caption{Ablation Study Overview}
\label{tab:ablation_summary}
\centering
\scriptsize
\begin{tabular}{|p{0.45cm}|p{2.4cm}|p{3.0cm}|p{3.3cm}|p{1.5cm}|}
\hline
\textbf{ID} &
\textbf{Component} &
\textbf{Replaced With} &
\textbf{x Variable} &
\textbf{Metrics} \\
\hline
A1 &
MATD handoff correction &
Raw uncorrected PDR &
Speed $v$ (km/h): \{30, 60, 90, 120, 150\} &
M1, M4 \\
\hline
A2 &
DP signatures S1--S3 &
Progressive subset: S1 $\to$ S1+S2 $\to$ S1+S2+S3 &
DP attack variant $\in$ \{FR, IT, TS\} &
M1, M3 \\
\hline
A3 &
CP signatures S4--S6 &
DP detection only; CP attacks undetected &
CP attack variant $\in$ \{FR, IT, TS\} &
M1, M3 \\
\hline
A4 &
LLM semantic scorer &
Signature $+$ reputation only ($\mu_2 = 0$) &
Drop rate $\rho \in \{10, 20, 30, 40\}\%$ on S2/S5 &
M1, M3 \\
\hline
A5 &
Federated learning &
Centralised static model; no round updates &
Non-IID $\alpha \in \{0.1, 0.3, 0.5, 1.0, \text{IID}\}$ &
M1, M10 \\
\hline
A6 &
Gradient hash integrity &
FL without blockchain commitment &
Poisoning fraction $\rho_p \in \{0, 10, 20, 33\}\%$ &
M10 \\
\hline
A7 &
DEBSC cryptographic gate &
Statistical-only gate; ZKP not required &
Drop rate $\rho \in \{10, 30, 50, 70\}\%$ &
M2, M4 \\
\hline
A8 &
ZKP ABSENT state &
Two-state ZKP (ABSENT $=$ undefined) &
Proof-withholding fraction $\in \{0, 25, 50, 75, 100\}\%$ &
M7 \\
\hline
A9 &
RSU cross-reference &
Cryptographic ZKP only; no independent count &
Fabricated commitment fraction $\in \{0, 25, 50, 75, 100\}\%$ &
M7 \\
\hline
A10 &
Probationary mode &
New vehicles enter ACTIVE immediately &
Pseudonym rotation rate (rot./min): \{0, 1, 2, 5, 10\} &
M4, M7 \\
\hline
A11 &
Aggregate anomaly $\Psi_c$ &
Hard-threshold S4--S6 triggers only &
Sub-threshold rules/window: \{1, 3, 5, 8, 10\} &
M8 \\
\hline
A12 &
Multi-controller architecture &
Single-controller; no trust, no failover &
Compromised controllers: \{0, 1, 2, 3\} &
M2, M11 \\
\hline
A13 &
Ordered-list failover &
Highest-trust selection across all segments &
Concurrently failing segments: \{1, 2, 3, $M{-}1$\} &
M11 \\
\hline
A14 &
DKG joint key generation &
Single-controller key generation &
Compromised entities: \{0, 1, 2, $t{-}1$, $t$\} &
M2, M11 \\
\hline
A15 &
PQC-LKH tree re-keying &
Unicast Kyber to all $N{-}1$ vehicles &
Population $N \in \{16, 32, 64, 128, 200\}$ &
M6, M9 \\
\hline
A16 &
VRF dynamic endorser selection &
Static first-$k$-to-respond endorsement &
Compromised RSU fraction $f_c \in \{0, 5, 10, 20, 33\}\%$ &
M2, M12 \\
\hline
A17 &
Flow rule whitelist governance &
Direct rule install; no RSU co-signature &
Malicious rules/min: \{1, 5, 10, 20\} &
M1, M3 \\
\hline
\end{tabular}
\end{table}

\mypar
\textbf{A1 --- MATD Handoff Correction.}
MATD (Eq.~\eqref{eq:matd}) is disabled and raw uncorrected PDR is fed
directly to all six detection signatures. The x variable is vehicle speed
$v$, swept from 30 to 150~km/h, because MATD's mobility correction
benefit scales with speed: at higher speeds the handoff-induced PDR drop
grows and is more easily mistaken for deliberate dropping. M4 (FIR) is
expected to increase monotonically with speed in the ablated configuration,
while M1 (MCC) degrades as false positives accumulate; the proposed
framework maintains low FIR across all speed bands through handoff loss
subtraction (Eq.~\eqref{eq:mobility_pdr}).

\mypar
\textbf{A2 --- Data-Plane Signatures S1--S3.}
Three progressively augmented sub-configurations are evaluated: S1-only
(fixed-rate detection), S1+S2 (adding intermittent detection), and the
full S1+S2+S3 set. The x variable is the DP attack variant injected
(FR, IT, TS), as each additional signature closes exactly one variant
coverage gap. M3 (AVCR) quantifies how many of the three DP variants
are covered at each sub-configuration level, while M1 (MCC) shows the
aggregate detection quality improvement as signatures are added;
S3 is expected to produce the largest AVCR step for the TS variant.

\mypar
\textbf{A3 --- Controller-Plane Signatures S4--S6.}
CP signatures are removed entirely, leaving the framework with data-plane
detection only; CP attacks are therefore undetected at the flow-rule level.
The x variable is the CP attack variant injected (FR, IT, TS), which
directly exposes which controller-plane grey hole behaviours remain
invisible without S4--S6. M1 (MCC) collapses to near-chance performance
for all CP variants, and M3 (AVCR) drops from 1.0 to at most 0.5,
quantifying the controller-plane blind spot that motivates the full
six-signature taxonomy.

\mypar
\textbf{A4 --- LLM Semantic Scorer.}
The LLM weight in the fusion equation is set to $\mu_2 = 0$
(Eq.~\eqref{eq:fusion}), reducing detection to signature plus reputation
only. The x variable is attack drop rate $\rho$ at low values
(10--40\%) applied specifically to the S2 intermittent and S5 CP-IT
variants, where rule-based autocorrelation signatures are most unreliable
and the LLM's temporal pattern reasoning provides the greatest incremental
gain. M1 (MCC) and M3 (AVCR) are expected to degrade most sharply at
$\rho \leq 20\%$ where the periodic drop signal falls below the
autocorrelation significance threshold $\gamma_{it}$ but remains
detectable as a sequential log anomaly by the LLM.

\mypar
\textbf{A5 --- Federated Learning.}
The FL pipeline is replaced by a single centralised model trained on a
fixed representative dataset at deployment and never updated during
simulation. The x variable is the Dirichlet non-IID heterogeneity
parameter $\alpha$ — lower $\alpha$ produces more skewed local data
distributions across vehicles — because FL's advantage over a static
model is maximised precisely when heterogeneity is high, which is the
SDVN-specific motivation for distributed training. M1 (MCC) and M10
(FL poisoning robustness MCC at $\rho_p = 0$) are evaluated; the
performance gap between FL and the static model is expected to widen as
$\alpha$ decreases toward 0.1.

\mypar
\textbf{A6 --- Gradient Hash Integrity.}
The blockchain hash commitment step is removed and gradient updates are
accepted by the aggregator without pre-round commitment verification
(Eq.~\eqref{eq:gradient_commit} disabled). The x variable is the
poisoning fraction $\rho_p$, swept from 0\% (no poisoning) to 33\%
(BFT bound). M10 ($\text{MCC}_{\text{FL}}(\rho_p)$) is evaluated
exclusively: at $\rho_p = 0$ the two configurations should be
equivalent, confirming no degradation from the integrity mechanism
itself; as $\rho_p$ increases the ablated configuration's MCC collapses
while the proposed framework maintains model quality through poisoned
gradient rejection (Eq.~\eqref{eq:gradient_verify}).

\mypar
\textbf{A7 --- DEBSC Cryptographic Gate.}
The DEBSC reverts to a statistical-only gate: isolation fires when the
reputation deficit alone exceeds $\theta_R$, without requiring
$\Pi_{\text{ZKP}}(v_i,t) \in \{\texttt{FAIL}, \texttt{ABSENT}\}$
(Eq.~\eqref{eq:debsc} reduced). The x variable is attack drop rate
$\rho$, swept across a wide range, because the cryptographic gate's
primary benefit appears at low drop rates where legitimate
handoff-induced PDR drops cross the reputation threshold. M4 (FIR)
is expected to surge at low $\rho$ for the ablated configuration as
legitimate vehicles with handoff loss are incorrectly isolated;
M2 (GHSR) degrades in turn because false isolations disrupt routing
paths that legitimate vehicles occupy.

\mypar
\textbf{A8 --- ZKP ABSENT State.}
The three-state ZKP model reverts to two-state: a missing proof is
treated as undefined rather than \texttt{ABSENT}, so the cryptographic
gate never fires solely from proof withholding. The x variable is the
fraction of attackers withholding proofs (rather than submitting false
ones), swept from 0\% (all attackers submit false proofs, detectable
in both models) to 100\% (all attackers simply do not submit). M7
(PEDR) is evaluated exclusively: at 0\% withholding both models perform
identically; PEDR for the two-state ablation collapses to zero as the
withholding fraction increases, while the proposed three-state model
maintains high PEDR through the \texttt{ABSENT}-state trigger.

\mypar
\textbf{A9 --- RSU Cross-Reference.}
RSU-independent forwarding count observation is removed; ZKP acceptance
is based on cryptographic proof validity alone without checking
consistency against the independently observed count
(Eq.~\eqref{eq:rsu_crossref} disabled). The x variable is the fraction
of attackers submitting fabricated Pedersen commitments that pass
cryptographic verification but misreport forwarding counts — an evasion
vector that only the cross-reference closes. M7 (PEDR) is evaluated
exclusively; PEDR for the ablated configuration is expected to degrade
linearly with fabrication fraction, while the proposed framework
maintains high PEDR by classifying cross-reference failures as
\texttt{FAIL} regardless of cryptographic validity.

\mypar
\textbf{A10 --- Probationary Mode.}
New vehicles enter \texttt{ACTIVE} mode immediately upon registration
with no single-strike ZKP rule and full statistical gate threshold
$\theta_R$ applied. The x variable is the pseudonym rotation rate
(fresh identity registrations per minute), simulating an attacker
continuously cycling IEEE~1609.2 pseudonyms to present as a new vehicle
with no adverse blockchain history. M4 (FIR) and M7 (PEDR) are
evaluated: without probationary mode, a pseudonym-cycling attacker
permanently avoids the statistical gate and PEDR collapses; FIR also
degrades as the unconstrained system attempts compensatory isolation of
ambiguous nodes.

\mypar
\textbf{A11 --- Aggregate Controller Anomaly Score $\Psi_c$.}
$\Psi_c(t)$ is removed; controller trust decrements only when a
confirmed S4--S6 hard-threshold trigger fires
(Eq.~\eqref{eq:controller_agg} disabled). The x variable is the number
of sub-threshold drop rules installed per observation window — the
intensity of the low-and-slow controller attack — swept from 1 to 10
rules per window. M8 (CADR) is evaluated exclusively: without aggregate
scoring the detection rate remains at zero for all sub-threshold rule
counts since no individual rule exceeds $\tau_c$; the proposed
framework accumulates CADR as sub-threshold evidence exceeds
$\psi_{\text{thresh}}$, revealing the aggregate evasion attack surface.

\mypar
\textbf{A12 --- Multi-Controller Flat Architecture.}
The multi-controller topology reverts to a single-controller
architecture with no per-controller trust score, no ordered list, and
no failover capability. The x variable is the number of simultaneously
compromised controllers, swept from 0 (no compromise, establishing
baseline parity) to 3. M2 (GHSR) and M11 (CFRT) are evaluated: in the
single-controller ablation, GHSR collapses immediately when the sole
controller is compromised ($x = 1$) as no failover exists, and CFRT
becomes undefined (infinite); the proposed multi-controller framework
maintains positive GHSR and bounded CFRT through ordered-list failover.

\mypar
\textbf{A13 --- Ordered-List Failover Selection.}
Ordered-list traversal (Eq.~\eqref{eq:failover_select}) is replaced by
highest-trust selection applied globally: all failing segments
simultaneously select the single highest-trust registered controller.
The x variable is the number of concurrently failing segments, swept
from 1 to $M-1$. M11 (CFRT) is evaluated: under highest-trust selection
the single chosen controller accumulates all fallback sessions and
progressively approaches a new single point of failure under load;
CFRT degrades as the overloaded replacement controller's response time
increases, while the ordered-list mechanism distributes failover load
across the fleet.

\mypar
\textbf{A14 --- DKG Joint Root Key Generation.}
DKG is replaced by single-controller Dilithium key generation: the
active controller alone holds the complete root Kyber key material and
signs all re-keying broadcasts unilaterally. The x variable is the
number of simultaneously compromised entities (controller plus RSUs),
swept from 0 to $t$ (the DKG reconstruction threshold). M2 (GHSR) and
M11 (CFRT) are evaluated: in the ablated configuration, compromising
the controller alone is sufficient to broadcast fraudulent re-keying
messages and cause GHSR collapse at $x = 1$; the DKG-protected
configuration requires $t$ simultaneous compromises before the root
key is recoverable, as formalised in Eq.~\eqref{eq:dkg_share}.

\mypar
\textbf{A15 --- PQC-LKH Tree Re-Keying.}
PQC-LKH is replaced by unicast Kyber re-keying: on each isolation event
the RSUs send $N-1$ individual Kyber encapsulations directly to each
remaining vehicle. The x variable is the vehicle population $N$, swept
from 16 to 200. M9 ($\eta_{\text{rekey}}(N)$) and the computation and
communication sub-metrics of M6 (MDPOS) are evaluated: the unicast
baseline is expected to grow linearly in both KEM operation count and
bandwidth, while the proposed PQC-LKH configuration produces
$\lceil\log_2 N\rceil$ operations and messages, empirically validating
the $O(\log N)$ complexity claim formalised in
Eq.~\eqref{eq:lkh_rekey}.

\mypar
\textbf{A16 --- VRF Dynamic Endorser Selection.}
VRF selection (Eq.~\eqref{eq:endorser_select}) is replaced by static
first-$k$-to-respond endorsement: the $k$ RSUs that reply fastest
within $T_q$ are accepted regardless of trust score. The x variable is
the compromised RSU fraction $f_c$, swept from 0\% to 33\% (the BFT
bound). M2 (GHSR) and M12 (EPAR) are evaluated: the static scheme
allows compromised RSUs to participate in endorsement at low $f_c$
values that the VRF trust filter would exclude, producing manipulated
isolation decisions and GHSR degradation; M12 shows the EPAR mode
distribution in each configuration as $f_c$ increases, revealing the
bypass mechanism activating at the correct thresholds under VRF selection.

\mypar
\textbf{A17 --- Flow Rule Whitelist Governance.}
The whitelist smart contract (Assumption~A5) is removed: the active
controller installs flow rules directly on OpenFlow switches without
requiring $k_{wl}$ RSU co-signatures. The x variable is the malicious
rule injection rate (rules per minute), representing how aggressively a
compromised controller injects S6 target-specific CP-TS rules. M1 (MCC)
and M3 (AVCR) are evaluated: without whitelist governance the controller
can pre-whitelist its own malicious rules to bypass S6, causing AVCR to
drop for the CP-TS variant as injection rate increases; the proposed
governance mechanism blocks this at the smart contract level, maintaining
stable AVCR regardless of injection rate.
% --- END SUPERVISOR REVISION BLOCK 17 ---

% --- SUPERVISOR REVISION BLOCK 19: Benchmark Experiments ---
\subsection{Benchmark Experiments against State-of-the-Art}
\label{sec:benchmark_experiments}

Five benchmark experiments are designed for system-level evaluation of
SHIELD-GH against baselines B1--B3. Experiments E1--E4 produce graphs
along a single or two-dimensional independent variable; no experiment
evaluates attack variants in isolation since SHIELD-GH is proposed as
a unified system covering all six variants simultaneously. Experiment E5
produces a single comparison table under default settings, broken down
per variant, to answer RQ1. Default simulation settings follow
Table~\ref{tab:sim_settings} unless overridden per experiment.

\mypar
\subsubsection{Experiment E1: Attack Penetration $\times$ Attack Intensity}
\label{sec:exp1}

\mypar
\textbf{Objective.}
To evaluate how SHIELD-GH and baselines B1--B3 respond as both the
prevalence and severity of the attack jointly increase across the full
two-dimensional attack space, addressing RQ2.

\mypar
\textbf{Independent variables.}
The experiment sweeps a two-dimensional grid defined by attacker
penetration $p$ and attack drop rate intensity $\rho_a$. Attacker
penetration is the fraction of all network nodes acting as grey hole
attackers across both data-plane and controller-plane roles. The grid
is given in Equation~\eqref{eq:e1_grid}.

\begin{equation}
\label{eq:e1_grid}
(p,\, \rho_a) \in
\{0, 20, 40, 60, 80, 100\}\%
\;\times\;
\{0, 20, 40, 60, 80, 100\}\%
\end{equation}

As given in Equation~\eqref{eq:e1_grid}, $p \in [0,1]$ is the
fraction of total network nodes assigned as attackers and $\rho_a
\in [0,1]$ is the packet drop rate applied by each attacker across
all six attack variants simultaneously. The grid yields $6 \times
6 = 36$ evaluation points equidistant in both dimensions. At
$p = 0\%$ no attack is active regardless of $\rho_a$; this
establishes the no-attack performance baseline. Controller compromise
is not independently varied: as $p$ increases, a proportionally
higher fraction of CP attack events accumulate in the simulation,
causing controller trust scores $\mathcal{T}_c(t)$
(Eq.~\eqref{eq:controller_trust}) to decrement and controller-plane
signatures S4--S6 to fire at higher penetration levels — controller
compromise is therefore implicit in $p$.

\mypar
\textbf{Fixed parameters.}
$N = 200$ vehicles, $v = 80$~km/h, all six attack variants active
simultaneously, simulation settings per Table~\ref{tab:sim_settings}.

\mypar
\textbf{Metrics evaluated.} M1, M2, M3, M4, M5.

\mypar
\textbf{Baselines compared.} B1, B2, B3.

\mypar
\textbf{Visualization.}
Six $6\times6$ heatmaps (one per metric per framework: SHIELD-GH and
each baseline), enabling direct visual comparison of performance
contours across the 2D attack space. Colour encodes metric value;
iso-performance contour lines are overlaid to identify degradation
boundaries.

\mypar
\subsubsection{Experiment E2: Vehicle Mobility}
\label{sec:exp2}

\mypar
\textbf{Objective.}
To evaluate how increasing vehicle speed affects detection accuracy,
false isolation, and response latency, and to quantify the contribution
of MATD relative to non-mobility-corrected baselines, addressing RQ3.

\mypar
\textbf{Independent variable.}
Target vehicle speed $v_{\text{target}}$ swept across four levels
spanning urban, arterial, rural, and motorway conditions, given in
Equation~\eqref{eq:e2_speed}.

\begin{equation}
\label{eq:e2_speed}
v_{\text{target}} \in \{10,\; 60,\; 100,\; 140\}\;\text{km/h}
\;\equiv\;
\{2.78,\; 16.67,\; 27.78,\; 38.89\}\;\text{m/s}
\end{equation}

As given in Equation~\eqref{eq:e2_speed}, each target speed in km/h
is converted to SUMO's native unit of m/s by dividing by 3.6. The
four values are selected to span the full range from low-mobility
urban scenarios (10~km/h, where handoff loss is negligible) to
high-mobility motorway scenarios (140~km/h, where handoff-induced
PDR drops are maximally confusable with grey hole dropping). The
MATD handoff loss model (Eq.~\eqref{eq:handoff_loss}) predicts
monotonically increasing $\rho_{\text{ho}}(v_i,t)$ with speed; this
experiment empirically validates that prediction.

\mypar
\textbf{SUMO mobility configuration.}
To ensure all vehicles travel at exactly $v_{\text{target}}$ without
stochastic deviation, the following SUMO vType parameters must be set
for every vehicle class in the simulation:

\begin{itemize}
  \item \texttt{maxSpeed} = target speed in m/s
        (e.g.\ \texttt{maxSpeed="38.89"} for 140~km/h)
  \item \texttt{speedFactor="1"} — disables per-vehicle speed
        variation around the road limit
  \item \texttt{speedDev="0"} — eliminates all stochastic speed
        deviation so every vehicle travels at exactly
        $v_{\text{target}}$
  \item All road and edge \texttt{speed} attributes in the OSM
        network set to the same target m/s value, overriding
        OpenStreetMap road class defaults
\end{itemize}

\noindent This configuration ensures the mobility trace produced by
SUMO and consumed by the NS-3 TraCI bridge
(\texttt{sumo\_ns3\_bridge.py}) delivers uniform per-vehicle speed
logs to the MATD module at each simulation step
(Eq.~\eqref{eq:matd}).

\mypar
\textbf{Fixed parameters.}
$N = 200$ vehicles, $p = 40\%$ attacker penetration,
$\rho_a = 40\%$ drop rate, all six attack variants active,
simulation settings per Table~\ref{tab:sim_settings}.

\mypar
\textbf{Metrics evaluated.} M1, M2, M4, M5,
$\Omega_{\text{comm}}(N)$ sub-metric of M6
(communication overhead increases with handoff frequency at
higher speeds).

\mypar
\textbf{Baselines compared.}
B1 and B3 (neither implements mobility correction; their M4 FIR is
expected to increase monotonically with speed, directly isolating
the MATD contribution). B2 is not included as it is a blockchain-only
scheme with no mobility model.

\mypar
\textbf{Visualization.}
Line plots per metric vs $v_{\text{target}}$ for SHIELD-GH, B1, B3.

\mypar
\subsubsection{Experiment E3: Network Scalability}
\label{sec:exp3}

\mypar
\textbf{Objective.}
To evaluate how detection performance, response latency, and protocol
overhead scale as the vehicle population grows, and to empirically
validate the $O(\log N)$ PQC-LKH re-keying complexity claim, addressing
RQ4.

\mypar
\textbf{Independent variable.}
Vehicle population $N$ swept across four levels, given in
Equation~\eqref{eq:e3_scale}.

\begin{equation}
\label{eq:e3_scale}
N \in \{100,\; 200,\; 300,\; 400\}\;\text{vehicles}
\end{equation}

As given in Equation~\eqref{eq:e3_scale}, the four levels are
equidistant with step size 100, spanning from a small-scale
deployment to a dense urban scenario. The RSU count scales
proportionally at one RSU per 25 vehicles, giving
$|\mathcal{R}| \in \{4, 8, 12, 16\}$ RSUs respectively, consistent
with the target deployment density in Table~\ref{tab:sim_settings}.
The VRF endorser pool $\mathcal{E}(t)$ (Eq.~\eqref{eq:endorser_pool})
and the minimum endorser count $k_{\text{end}}$
(Eq.~\eqref{eq:kend}) scale with $|\mathcal{R}|$ at each level.

\mypar
\textbf{Fixed parameters.}
$v = 80$~km/h, $p = 40\%$ attacker penetration, $\rho_a = 40\%$
drop rate, all six attack variants active, simulation settings per
Table~\ref{tab:sim_settings}.

\mypar
\textbf{Metrics evaluated.}
M1, M2, M4, M5, M6 (all three sub-metrics $\Omega_{\text{comp}}(N)$,
$\Omega_{\text{comm}}(N)$, $\Omega_{\text{store}}(N,t)$ as functions
of $N$), M9 ($\eta_{\text{rekey}}(N)$ — PQC-LKH vs unicast Kyber
operation count ratio).

\mypar
\textbf{Baselines compared.}
B1, B2, B3 for M1/M2/M4/M5; SHIELD-GH vs unicast Kyber re-keying
baseline for M6/M9 (no reviewed baseline implements PQC-LKH).

\mypar
\textbf{Visualization.}
Line plots per metric vs $N$. M6 produces three subplots
($\Omega_{\text{comp}}$, $\Omega_{\text{comm}}$,
$\Omega_{\text{store}}$). M9 shows measured KEM operation count for
PQC-LKH vs unicast as $N$ increases, with the theoretical
$\lceil\log_2 N\rceil$ curve overlaid for validation.

\mypar
\subsubsection{Experiment E4: Attack Evasiveness Index}
\label{sec:exp4}

\mypar
\textbf{Objective.}
To evaluate how detection coverage and mitigation effectiveness degrade
as all grey hole attack variants are simultaneously tuned to maximise
evasion of the framework's detection mechanisms, and to determine at
what evasiveness level the lightweight rule-based mode fails and the
full LLM mode becomes the primary detection path, addressing RQ5.

\mypar
\textbf{Construction of the Attack Evasiveness Index.}
The three variant-specific evasion parameters are each normalised to
a unit evasiveness sub-score. The sub-scores for intermittent attack
period, malicious epoch duty cycle, and targeting selectivity are
given in Equation~\eqref{eq:aei_components}.

\begin{equation}
\label{eq:aei_components}
\xi_T = \frac{T^*}{T_{\max}},
\quad
\xi_d = 1 - d,
\quad
\xi_k = 1 - \frac{k_t}{K}
\end{equation}

As given in Equation~\eqref{eq:aei_components}, $T^*$ is the
intermittent attack cycle period, $T_{\max}$ is the upper bound of
the autocorrelation detection window in Eq.~\eqref{eq:sig_dpit},
$d \in (0,1]$ is the duty cycle (fraction of the attack cycle spent
in the malicious dropping state), $k_t$ is the number of traffic
flows or source vehicles targeted by the S3/S6 attacker, and $K$ is
the total number of active flows. Each sub-score is zero when the
corresponding variant is at its most detectable configuration and
approaches one when it is maximally evasive. The aggregate Attack
Evasiveness Index is given in Equation~\eqref{eq:aei}.

\begin{equation}
\label{eq:aei}
\xi =
\frac{1}{3}\!\left(
\xi_T + \xi_d + \xi_k
\right)
=
\frac{1}{3}\!\left(
\frac{T^*}{T_{\max}}
+ (1 - d)
+ \left(1 - \frac{k_t}{K}\right)
\right)
\end{equation}

As given in Equation~\eqref{eq:aei}, $\xi \in [0,1]$ is a unified
scalar that increases as all three attack variants simultaneously
become harder to detect. At $\xi = 0$ all three variants are
maximally detectable: the intermittent period is well inside
$[T_{\min}, T_{\max}]$, duty cycle is high, and many flows are
targeted. At $\xi = 1$ all three variants are maximally evasive:
the intermittent period reaches $T_{\max}$ (outside the
autocorrelation window), duty cycle is minimal (sparse malicious
epochs), and only a single flow is targeted (minimum KL divergence
signal in Eq.~\eqref{eq:sig_dpts}).

\mypar
\textbf{Independent variable.}
$\xi$ is swept across six equidistant levels, given in
Equation~\eqref{eq:e4_levels}.

\begin{equation}
\label{eq:e4_levels}
\xi \in \{0.0,\; 0.2,\; 0.4,\; 0.6,\; 0.8,\; 1.0\}
\end{equation}

As given in Equation~\eqref{eq:e4_levels}, each level $\xi_\ell$
maps to a unique triplet $(T^*, d, k_t)$ via the parameter functions
in Equations~\eqref{eq:e4_param_T}, \eqref{eq:e4_param_d},
and~\eqref{eq:e4_param_kt}.

\begin{equation}
\label{eq:e4_param_T}
T^*(\xi) = \xi \cdot T_{\max}
\end{equation}

\begin{equation}
\label{eq:e4_param_d}
d(\xi) = d_{\max} - \xi\,(d_{\max} - d_{\min})
\end{equation}

\begin{equation}
\label{eq:e4_param_kt}
k_t(\xi) =
\max\!\left(1,\;
\left\lfloor (1 - \xi) \cdot K \right\rfloor
\right)
\end{equation}

As given in Equation~\eqref{eq:e4_param_T}, the intermittent
period $T^*$ increases linearly from $0$ to $T_{\max}$ as $\xi$
increases, moving from deep inside the autocorrelation detection
window toward its outer boundary, where S2/S5 signature sensitivity
degrades. As given in Equation~\eqref{eq:e4_param_d}, the duty
cycle $d$ decreases linearly from $d_{\max}$ (dense malicious epochs,
easily detected) to $d_{\min}$ (sparse epochs, minimum evidence per
window). As given in Equation~\eqref{eq:e4_param_kt}, the targeted
flow count $k_t$ decreases from $K$ (all flows targeted,
maximum KL divergence, detectable as fixed-rate) to 1 (single flow
targeted, minimum divergence signal). Fixed-rate attacks S1/S4 remain
active at $\rho_a = 40\%$ throughout all $\xi$ levels — they have no
evasiveness dimension beyond intensity, which is already swept in E1.
The default simulation parameter values used are:
$T_{\max} = 30$~s, $d_{\max} = 0.8$, $d_{\min} = 0.05$,
$K = 20$ active flows.

\mypar
\textbf{Fixed parameters.}
$N = 200$ vehicles, $v = 80$~km/h, $p = 40\%$ attacker
penetration, $\rho_a = 40\%$ drop rate for S1/S4, simulation
settings per Table~\ref{tab:sim_settings}.

\mypar
\textbf{Metrics evaluated.} M1, M2, M3, M4, M5.

\mypar
\textbf{Baselines compared.}
B1 and B3. B1 relies purely on threshold-based detection and is
expected to collapse at moderate $\xi$ when lightweight signatures
fail; B3 (FL-BERT) has partial temporal reasoning but no S3/S6
controller-plane evasion model. SHIELD-GH is expected to maintain
coverage beyond the $\xi$ level at which B1/B3 fail, through the
LLM full-mode pipeline.

\mypar
\textbf{Visualization.}
Line plots per metric vs $\xi$ for SHIELD-GH and baselines, with a
vertical dashed line marking the $\xi$ level at which the lightweight
mode MCC falls below 0.5 (SHIELD-GH lightweight mode handoff point)
and the full LLM mode becomes the primary detection path.

\mypar
\subsubsection{Experiment E5: Per-Variant Comparison Table under Default Settings}
\label{sec:exp5}

\mypar
\textbf{Objective.}
To evaluate SHIELD-GH against baselines B1--B3 at the individual
attack variant level under a single default configuration, producing
a quantitative breakdown that answers RQ1 directly.

\mypar
\textbf{Setup.}
Single fixed configuration: $N = 200$ vehicles, $v = 80$~km/h,
$p = 40\%$ attacker penetration, $\rho_a = 40\%$ drop rate,
$\xi = 0.4$ (moderate evasiveness). Each of the six attack variants
S1--S6 is injected in isolation for a separate simulation run, with
all other variants inactive. This is the only experiment that
evaluates variants individually; it produces one data point per
cell rather than a curve.

\mypar
\textbf{Result format.}
A single comparison table with rows S1 (DP-FR), S2 (DP-IT),
S3 (DP-TS), S4 (CP-FR), S5 (CP-IT), S6 (CP-TS) and columns for
SHIELD-GH, B1, B2, B3. Each cell reports M1 (MCC), M2 (GHSR),
M4 (FIR), and M5 (ESRL). Cells for which a baseline has no
detection capability for the given variant are marked N/A:
B1 covers S1 and a single S2 variant only; B2 covers S1 and S2
data-plane only (no CP detection); B3 covers S1, S2, S3 data-plane
only. All four S4--S6 rows for B1, B2, B3 are therefore N/A,
revealing the controller-plane detection gap that SHIELD-GH closes.

\mypar
\textbf{Metrics evaluated.} M1, M2, M4, M5.

\mypar
\textbf{Baselines compared.} B1, B2, B3.
% --- END SUPERVISOR REVISION BLOCK 19 ---

%% ============================================================
\section{Current Simulation Results}
\label{sec:results_progress}
%% ============================================================

The following results were obtained from an ns-3.35 simulation of a five-node
Software-Defined Vehicular Network (SDVN) under a \textbf{fixed-rate grey hole attack}
(S1 variant). Node~3 acts as the malicious attacker, selectively dropping forwarded
packets at a fixed rate. Two simulation runs are presented, each capturing the
network state at a different simulation time under the same attack configuration.
No mitigation mechanism is active in either run; these results characterise the
raw attack impact for subsequent comparison against the full SHIELD-GH framework.

\subsection{Run 1 -- Fixed-Rate Attack (Higher Forwarding Rate)}
\label{subsec:run1}

Figure~\ref{fig:run1_attack} shows the ns-3.35 simulation console output for Run~1 of the fixed-rate grey hole attack, captured at simulation time 11:20.

\begin{figure}[H]
  \centering
  \includegraphics[width=\textwidth]{Attacks/Screenshot from 2026-05-20 11-20-32.png}
  \caption{ns-3.35 simulation log: fixed-rate grey hole attack (S1 variant), Run~1.
           Node~3 (attacker) achieves a per-node PDR of 95.24\%, indicating a
           lower drop rate in this run. Network-wide average PDR is 86.54\% and
           detection accuracy is 92.39\%.}
  \label{fig:run1_attack}
\end{figure}

Table~\ref{tab:run1_metrics} summarises the key metrics extracted from Run~1.

\begin{table}[H]
\renewcommand{\arraystretch}{1.3}
\caption{Run 1 -- Observed Metrics (Fixed-Rate Attack, No Mitigation)}
\label{tab:run1_metrics}
\centering
\small
\begin{tabular}{|l|c|}
\hline
\textbf{Metric} & \textbf{Value} \\
\hline
\multicolumn{2}{|l|}{\textit{Network measurements (no mitigation active)}} \\
\hline
Network-wide average PDR
  ($\overline{\text{PDR}}_{\text{attack}}$, S1 baseline)  & 86.54\% \\
Average end-to-end latency                                  & 7.77 ms \\
Node 3 PDR (attacker)                                       & 95.24\% \\
\hline
\multicolumn{2}{|l|}{\textit{Detection quality (M1) ---
  S1 variant only, no isolation active}} \\
\hline
True Positives (TP)   & 85 \\
False Positives (FP)  & 7 \\
Detection accuracy (raw) & 92.39\% \\
M1 --- MCC            & 0$^{\dagger}$ \\
\hline
\multicolumn{2}{|l|}{\textit{Metrics requiring active mitigation}} \\
\hline
M2 --- GHSR           & N/A$^{\ddagger}$ \\
M4 --- FIR            & N/A$^{\ddagger}$ \\
M5 --- ESRL           & N/A$^{\ddagger}$ \\
\hline
\end{tabular}
\vspace{1mm}
\footnotesize{%
$^{\dagger}$~MCC\,=\,0 because the early prototype records no true-negative
or false-negative events (all nodes flagged); this is a known prototype
limitation, not a framework result.
$^{\ddagger}$~M2~(GHSR), M4~(FIR), and M5~(ESRL) require the full
SHIELD-GH mitigation pipeline to be active; these runs characterise the raw
attack impact for Milestone~14 comparison.}
\end{table}

\subsection{Run 2 -- Fixed-Rate Attack (Higher Drop Rate)}
\label{subsec:run2}

Figure~\ref{fig:run2_attack} shows the ns-3.35 simulation console output for Run~2, where Node~3 exhibits a higher packet drop rate, captured at simulation time 11:22.

\begin{figure}[H]
  \centering
  \includegraphics[width=\textwidth]{Attacks/Screenshot from 2026-05-20 11-22-27.png}
  \caption{ns-3.35 simulation log: fixed-rate grey hole attack (S1 variant), Run~2.
           Node~3 [ATTACKER] drops a larger proportion of forwarded packets,
           resulting in a per-node PDR of only 68.25\%. Network-wide average PDR
           degrades to 68.71\% and detection accuracy falls to 71.74\%.}
  \label{fig:run2_attack}
\end{figure}

Table~\ref{tab:run2_metrics} summarises the key metrics extracted from Run~2.

\begin{table}[H]
\renewcommand{\arraystretch}{1.3}
\caption{Run 2 -- Observed Metrics (Fixed-Rate Attack, No Mitigation)}
\label{tab:run2_metrics}
\centering
\small
\begin{tabular}{|l|c|}
\hline
\textbf{Metric} & \textbf{Value} \\
\hline
\multicolumn{2}{|l|}{\textit{Network measurements (no mitigation active)}} \\
\hline
Network-wide average PDR
  ($\overline{\text{PDR}}_{\text{attack}}$, S1 baseline)  & 68.71\% \\
Average end-to-end latency                                  & 7.82 ms \\
Node 3 PDR (attacker)                                       & 68.25\% \\
\hline
\multicolumn{2}{|l|}{\textit{Detection quality (M1) ---
  S1 variant only, no isolation active}} \\
\hline
True Positives (TP)   & 66 \\
False Positives (FP)  & 26 \\
Detection accuracy (raw) & 71.74\% \\
M1 --- MCC            & 0$^{\dagger}$ \\
\hline
\multicolumn{2}{|l|}{\textit{Metrics requiring active mitigation}} \\
\hline
M2 --- GHSR           & N/A$^{\ddagger}$ \\
M4 --- FIR            & N/A$^{\ddagger}$ \\
M5 --- ESRL           & N/A$^{\ddagger}$ \\
\hline
\end{tabular}
\vspace{1mm}
\footnotesize{%
$^{\dagger}$~MCC\,=\,0; see Run~1 note.
$^{\ddagger}$~M2, M4, M5 not available without active mitigation pipeline.}
\end{table}

Comparing Run~1 and Run~2 shows the sensitivity of network performance to the
attacker's drop rate under the fixed-rate variant (S1). As the drop rate increases
from Run~1 to Run~2, network-wide PDR falls by 17.83 percentage points (86.54\%
to 68.71\%) and detection accuracy degrades from 92.39\% to 71.74\%. These preliminary results confirm that the ns-3.35
prototype correctly reproduces the expected grey hole attack degradation and
establish a baseline against which the full SHIELD-GH detection and mitigation
framework will be evaluated in Milestone~14.

%% ============================================================
\section{Expected Results}
\label{sec:expected_results}
%% ============================================================

The following sections present the expected performance results of the SHIELD-GH
framework based on analytical modelling of the six attack variants and comparison
with the three external baselines (B1--B3) and ablations A1--A17 defined in
Section~\ref{sec:ablation}. Since full simulation experiments are scheduled for
Milestone~14 (see Chapter~5), the values reported here are
derived from the formal attack models in Sections~\ref{sec:signatures}
and~\ref{sec:detection_methodology}, calibrated against the parameter ranges
reported in the closest existing works~\cite{malik2023greyhole,
alabdulatif2024mitigating, ahsan2024flbert}. Each table specifies the
metrics evaluated, the baseline of comparison, and the expected direction
of improvement for each scenario.

\subsection{How Detection Metrics Vary with Attack Variant and Intensity}

The six attack variants differ in their observability: fixed-rate variants
produce a stable low-PDR signal that rule-based detectors capture quickly,
while intermittent and target-specific variants require longer observation
windows and higher-capacity detectors (LLM, FL). The expected M1 (MCC) and
M3 (AVCR) for each variant are evaluated under both the lightweight and
full detection modes at a representative drop rate of 40\%.

Controller-plane signatures (S4--S6) are expected to achieve slightly higher
accuracy than their data-plane counterparts because they operate on the
deterministic flow-rule set rather than probabilistic PDR statistics, reducing
the influence of channel noise. Intermittent variants (S2, S5) show the largest
gap between lightweight and full mode: lightweight detection relies on
autocorrelation periodicity which fails when the attack cycle length falls
outside $[T_{\min}, T_{\max}]$, whereas the LLM captures irregular intermittent
bursts as sequential anomalies without requiring a fixed period assumption.

\subsection{How Packet Delivery Ratio Varies with Drop Rate}

The expected network-layer PDR is evaluated at varying attacker drop rates
$\rho \in \{20\%, 40\%, 60\%, 80\%\}$ for the fixed-rate data-plane variant
(S1), representing the most straightforward impact of a grey hole attack on
throughput. Three conditions are compared: no framework active, the lightweight
SHIELD-GH mode active, and the full SHIELD-GH mode active. The PDR is measured
after mitigation is complete (i.e., the attacker has been isolated).

The PDR recovery gap between no-framework and SHIELD-GH widens with drop rate,
because the severity of the attack increases the urgency and benefit of isolation.
At low drop rates (20\%), the attack itself causes less harm and B2's blockchain
contract provides partial benefit; at high drop rates (80\%), SHIELD-GH's faster
detection and more reliable isolation provide substantially greater PDR recovery.
The lightweight mode reaches a slightly lower post-isolation PDR than the full
mode because intermittent or borderline fixed-rate attackers may not be caught
until the next observation window.

\subsection{How Detection Accuracy Varies with Attack Intensity}

Detection accuracy is expected to vary as the attacker drop rate increases
from 10\% to 90\%, for both the lightweight and full detection modes. At very
low drop rates ($\leq$15\%), the attack signal approaches the noise level of
legitimate channel loss, causing both modes to approach random-chance
performance; below this threshold the MATD correction cannot isolate the
attack signal from handoff-induced loss. Above 30\%, the lightweight mode
saturates quickly; the full mode provides gains primarily in the 15--35\%
region where rule-based signatures are unreliable but temporal LLM patterns
remain distinguishable.

The proposed framework is expected to consistently outperform B1 (threshold-only,
no mobility correction) and B3 (FL-BERT, no controller-plane model) across all
drop rates, with the largest gap at low drop rates (10--30\%) where the
MATD mobility correction and DEBSC dual-evidence gate prevent false
positives that inflate B1's FPR and degrade B3's MCC.

\subsection{How False Positive Rate Varies with Vehicle Speed (Mobility Effect)}

A key concern in SDVN is that high-speed vehicles undergoing RSU handoffs
exhibit transient PDR drops that resemble attack behaviour. Without mobility
correction, a threshold-based detector will falsely isolate legitimate
high-speed vehicles. The expected M4 (FIR) for legitimate vehicles at increasing
speeds is evaluated by comparing three conditions: no mobility correction
(equivalent to B1 and ablation A1), MATD correction only (MATD enabled but
DEBSC cryptographic gate disabled, i.e.\ ablation A7 with MATD restored), and
the full DEBSC dual-evidence gate (proposed).

Without mobility correction (ablation A1), M4 (FIR) rises sharply above
60~km/h as handoff-induced PDR drops trigger the fixed-rate signature. MATD
correction reduces FIR by roughly 60\% at all speeds. The DEBSC cryptographic
gate (ablation A7 reversed) provides the final reduction: a legitimate fast
vehicle with a temporarily low PDR will still hold valid ZKP forwarding
proofs for the packets it did forward, so the cryptographic gate blocks false
isolation even when the statistical gate fires.

\subsection{Detection Latency and Mitigation Response Time}

The expected detection latency (time from attack onset to confirmed detection
verdict) and mitigation response time (time from verdict to completed node
isolation) are evaluated for both modes, broken down by pipeline stage.
PQC-LKH operational overhead relative to unicast Kyber re-keying is also
quantified via ablation A15 and the M6 (MDPOS) computation sub-metric.

The lightweight mode satisfies the 1-second end-to-end latency budget required
for safety-critical vehicular applications. The full mode incurs additional
latency from LLM inference, but the two-tier edge-cloud design (Section~3.7.3)
ensures that the majority of decisions are resolved at the edge tier
($\leq$80~ms) without escalation to the cloud. The PQC overhead of 12--18~ms
per Dilithium signing operation is acceptable relative to the total mitigation
pipeline, and the $O(\log N)$ PQC-LKH re-keying is expected to be 3--4$\times$
faster than naive unicast Kyber re-keying for networks of $N \geq 16$ vehicles.

\subsection{Ablation Study: Component Contribution to Detection and
Mitigation Performance}

The expected impact of removing each framework component (ablations
A1--A17, Table~\ref{tab:ablation_summary}) is evaluated per component at
the independent variable and metrics specified in each ablation design
(Section~\ref{sec:ablation}). The following qualitative expectations
correspond to key ablations under the S2 intermittent variant at the
default operating point ($\rho_a = 40\%$, $p = 40\%$, $N = 200$,
$v = 80$~km/h).

Disabling the LLM semantic scorer (A4, $\mu_2 = 0$) produces the largest
M1 (MCC) drop for the intermittent variant, confirming that temporal
pattern reasoning is essential when autocorrelation signatures fail at low
drop rates. Replacing FL with a centralised static model (A5) also degrades
M1 (MCC), reflecting the non-IID distribution problem quantified in
Section~\ref{sec:mobility_analysis}. Disabling MATD (A1) causes the largest
M4 (FIR) increase at high vehicle speeds, confirming that mobility
correction is the primary mechanism preventing false isolation of high-speed
vehicles. Removing the DEBSC cryptographic gate (A7) substantially increases
M4 (FIR), showing that the ZKP evidence gate independently blocks false
isolations that the statistical gate alone would trigger. Replacing PQC-LKH
with unicast Kyber re-keying (A15) does not affect M1 or M4 since
cryptography contributes to mitigation rather than detection, but it degrades
the M6 (MDPOS) computation and communication sub-metrics and eliminates the
$O(\log N)$ scaling benefit validated by M9 ($\eta_{\text{rekey}}$).

\chapter{Timeline and Resource Required}

\section{Timeline}

The timeline provides an overview of project tasks scheduled for December 2025 to August 2026, showing key project milestones along with their expected durations in weeks.

\begin{itemize}
    \item \textbf{MILESTONE 01: Project Topic Selection} \\
    Tasks: Select topic, read proposal description, get supervisor approval.
    \item \textbf{MILESTONE 02: Supervisor Discussion} \\
    Tasks: Discuss project feasibility and approach.
    \item \textbf{MILESTONE 03: Literature Review} \\
    Tasks: Study SDVN, Grey Hole attacks, blockchain, and related work.
    \item \textbf{MILESTONE 04: Research Gap Analysis} \\
    Tasks: Identify gaps and limitations in existing work.
    \item \textbf{MILESTONE 05: Proposal Preparation} \\
    Tasks: Draft and finalize project proposal.
    \item \textbf{MILESTONE 06: Technology Study} \\
    Tasks: Study NS3, SUMO, Blockchain, ML tools.
    \item \textbf{MILESTONE 07: System Architecture Design} \\
    Tasks: Design overall system structure.
    \item \textbf{MILESTONE 08: Threat Model Design} \\
    Tasks: Identify possible Grey Hole attack scenarios.
    \item \textbf{MILESTONE 09: NS3-SUMO Setup} \\
    Tasks: Install and configure simulation environment.
    \item \textbf{MILESTONE 10: Basic SDVN Simulation} \\
    Tasks: Simulate SDVN without attacks.
    \item \textbf{MILESTONE 11: Blockchain Implementation} \\
    Tasks: Deploy Hyperledger Fabric for node logging.
    \item \textbf{MILESTONE 12: LLM Integration} \\
    Tasks: Integrate LLM for monitoring network logs.
    \item \textbf{MILESTONE 13: Federated Learning Setup} \\
    Tasks: Configure FL across simulated nodes.
    \item \textbf{MILESTONE 14: Attack Scenarios Modeling and Simulation} \\
    Tasks: Implement Grey Hole attacks and test mitigation.
    \item \textbf{MILESTONE 15: System Integration and Testing} \\
    Tasks: Integrate all components and run tests.
    \item \textbf{MILESTONE 16: Performance Evaluation} \\
    Tasks: Analyze detection accuracy, latency, and network performance.
    \item \textbf{MILESTONE 17: Final Documentation} \\
    Tasks: Write final report and prepare diagrams.
    \item \textbf{MILESTONE 18: Final Presentation and Viva Preparation} \\
    Tasks: Prepare slides and rehearse presentation.
\end{itemize}

The full schedule is shown in Figure~\ref{fig:timeline}.


\clearpage % Force new page
\thispagestyle{empty} % Optional: remove page number for timeline page

% \begin{figure}[h!]
%     \centering
%     \includegraphics[width=\textwidth, height=0.45\textheight, keepaspectratio]{timeline.png}
%     \caption{Project timeline from December 2025 to August 2026.}
%     \label{fig:timeline}
% \end{figure}

\begin{figure*}[p]
    \centering
    \includegraphics[
        width=\textwidth,
        height=0.95\textheight,
        keepaspectratio
    ]{timeline.png}
    \caption{Project timeline from December 2025 to August 2026.}
    \label{fig:timeline}
\end{figure*}

\clearpage % Force next section to start on a new page



\section{Resources Required}

This section provides a succinct description of the hardware and software tools that are necessary to complete the proposed project successfully on detecting and mitigating gray hole attacks in software defined vehicular networks

\subsection{Hardware Resources}
\begin{itemize}
    \item \textbf{Computer / Workstation:}
    A system with at least an intel i7 or Ryzen 7 processor and a minimum of 16GB RAM also sufficient storage to execute network simulations, blockchain components, and machine learning tasks.
\end{itemize}

\subsection{Network Simulation Tools}
\begin{itemize}
    \item \textbf{NS-3 (Network Simulator 3):}
    Used to simulate the Software-Defined Vehicular Network environment and Grey Hole attack scenarios.
    \item \textbf{SUMO (Simulation of Urban Mobility):}
    Used to generate realistic vehicular mobility patterns for the simulations.
    \item \textbf{C++ and Python:}
    Used for NS-3 scripting, simulation control, and preprocessing simulation data.
\end{itemize}

\subsection{Machine Learning and Security Tools}
\begin{itemize}
    \item \textbf{Python:}
    Used for implementing machine learning models and processing simulation data.
    \item \textbf{TensorFlow / PyTorch:}
    Used for developing and training learning models.
    \item \textbf{Hugging Face Transformers:}
    Used for integrating Large Language Models (LLMs) to analyze network behavior and assist in Grey Hole attack detection and response.
    \item \textbf{Flower Framework:}
    Used to implement Federated Learning for distributed, privacy-preserving model training across vehicular nodes.
    \item \textbf{PyCryptodome:}
    Used for cryptographic operations and security-related functions.
\end{itemize}

\subsection{Blockchain Platform}
\begin{itemize}
    \item \textbf{Hyperledger Fabric:}
    Used to implement a secure and tamper-proof distributed ledger for logging node behaviors.
    \item \textbf{Docker:}
    Used for containerized deployment of blockchain network components.
\end{itemize}

\subsection{Analysis and Visualization Tools}
\begin{itemize}
    \item \textbf{Python Libraries (NumPy, Pandas, Matplotlib):}
    Used to analyze simulation results and visualize performance metrics.
    \item \textbf{Wireshark (Optional):}
    Used for packet-level analysis and validation.
\end{itemize}

\subsection{Version Control and Documentation}
\begin{itemize}
    \item \textbf{GitHub:}
    Used for version control and collaboration among team members.
    \item \textbf{Visual Studio Code / PyCharm:}
    Used as integrated development environments for coding and debugging.
    \item \textbf{LaTeX / Overleaf:}
    Used for preparing the project proposal, reports, and final documentation.
    \item \textbf{Draw.io / Lucidchart:}
    Used for designing system architecture and workflow diagrams.
\end{itemize}


\chapter{Conclusion}
% Summarize the key points of your proposal and reiterate the importance of the project.
The project suggests an innovative and feasible solution to minimize the grey hole attacks in the software defined vehicular networks. Grey hole attacks are also hard to detect since malicious vehicles selectively drop packets in various ways, including dropping packets only at specific times or specific messages, and most of ways. Thus, such types of attacks can have a severe impact on road safety and network performance of vehicular network. This project also considers scenarios where the network controller may be malicious and place malicious flow rules between the vehicle network.

\mypar
To address this issue in vehicular network, the project integrates three current technologies which are Blockchain, Large Language Models, and Federated Learning. The packet forwarding activities are safely recorded using blockchain to ensure that no behavior of the nodes is modified or concealed. Large Language Models are useful in network data analysis and detecting abnormal patterns that point to Grey Hole attacks. Federated Learning enables vehicles to jointly train detection models without exchanging their own data, and enhances privacy and scalability of a distributed network.

\mypar
Simulations will be used to assess how well the system is able to detect various types of attack by Gray Hole and ensure that the network communication is not disrupted. The combination of security, learning, and analysis techniques is likely to enhance the trust between vehicles. The project will also assist in learning how the modern technologies can be used to address the real world security issues in the vehicular networks.


% References
\renewcommand{\bibname}{References}
\bibliographystyle{ieeetr}
\addcontentsline{toc}{chapter}{References} % Add to table of contents
\bibliography{bibliography} 
%\printbibliography

\end{document}
